\documentclass[UTF8,a4paper,12pt,fontset=fandol,scheme=chinese]{ctexart}
% ============================================================
%  SY02-36 约化群的 Kloosterman 层 —— 导言区
%  原文：J. Heinloth, B.-C. Ngô, Z. Zhiwei Yun,
%        "Kloosterman sheaves for reductive groups"
%        arXiv:1005.2765v1 (16 May 2010), 58 页
% ============================================================

\usepackage[a4paper,margin=2.6cm]{geometry}

\PassOptionsToPackage{leqno}{amsmath}% 原文公式号在左
\usepackage{amsmath,amssymb,amsfonts,amsthm}
\usepackage{mathtools}
\usepackage{mathrsfs}
\usepackage{bm}
\usepackage{graphicx}
\usepackage{booktabs}
\usepackage{array}
\usepackage{tabularx}
\usepackage{enumitem}
\usepackage{xcolor}
\usepackage[most]{tcolorbox}
\usepackage{xspace}
\usepackage{tikz}
\usetikzlibrary{cd}
\usepackage{tikz-cd}
\usepackage{microtype}

\usepackage[bookmarks=true,bookmarksnumbered=true,
  colorlinks=true,linkcolor=blue!45!black,citecolor=blue!45!black,
  urlcolor=teal!70!black,filecolor=teal!70!black,
  pdftitle={约化群的 Kloosterman 层},
  pdfauthor={J. Heinloth, B.-C. Ngô, Z. Yun（译者：数译）}]{hyperref}

\linespread{1.22}
\setlength{\parskip}{.25em}
\setlength{\emergencystretch}{3em}
\sloppy

% ---------- 中文排版 ----------
\xeCJKsetup{CJKecglue={\hskip 0.24em plus 0.05em minus 0.05em}}
\ctexset{
  section    = {name={,}, number=\thesection,    aftername={\enspace}},
  subsection = {name={,}, number=\thesubsection, aftername={\enspace}},
  subsubsection = {name={,}, number=\thesubsubsection, aftername={\enspace}},
}
% 引言为「第 0 节」：section 计数器从 -1 起，显示为 0
\renewcommand{\thesection}{%
  \ifnum\value{section}<0 0\else\arabic{section}\fi}
\renewcommand{\thesubsection}{\thesection.\arabic{subsection}}
\renewcommand{\thesubsubsection}{\thesubsection.\arabic{subsubsection}}

% ---------- 编号 ----------
% 原文：主定理 Theorem 1–4 全局编号（用 mainthm 手写号）；
% 其余定理类与编号公式共用一个按节递增的计数器。
\numberwithin{equation}{section}
\theoremstyle{plain}
\newtheorem{theorem}[equation]{定理}
\newtheorem{lemma}[equation]{引理}
\newtheorem{proposition}[equation]{命题}
\newtheorem{corollary}[equation]{推论}
\newtheorem{definition}[equation]{定义}
\newtheorem{example}[equation]{例}
\newtheorem{remark}[equation]{注}
\newtheorem{claim}[equation]{断言}
\newtheorem{notation}[equation]{记号}
\newtheorem{conjecture}[equation]{猜想}
\newtheorem{fact}[equation]{事实}
\theoremstyle{definition}

% 主定理：全局号 1–4，不占节内计数器
\newenvironment{mainthm}[1][]{%
  \par\addvspace{\medskipamount}%
  \noindent\textbf{定理 #1.}\ \itshape\ignorespaces
}{%
  \par\addvspace{\medskipamount}%
}
\newenvironment{mainthmstar}[1][]{%
  \par\addvspace{\medskipamount}%
  \noindent\textbf{定理 #1}\ \itshape\ignorespaces
}{%
  \par\addvspace{\medskipamount}%
}

\setcounter{tocdepth}{2}
\setcounter{secnumdepth}{3}

% ---------- 译注 ----------
\newcounter{transnote}
\newtcolorbox{yizhu}{
  colback=orange!4, colframe=orange!55!black, boxrule=.4pt,
  left=6pt, right=6pt, top=4pt, bottom=4pt,
  fontupper=\small,
  before skip=6pt, after skip=6pt,
}
\makeatletter
\renewenvironment{yizhu}{%
  \refstepcounter{transnote}%
  \begin{tcolorbox}[colback=orange!4, colframe=orange!55!black, boxrule=.4pt,
    left=6pt, right=6pt, top=4pt, bottom=4pt, fontupper=\small,
    before skip=6pt, after skip=6pt]%
  \noindent\textbf{译注 \thetransnote：}\ignorespaces
}{%
  \end{tcolorbox}%
}
\makeatother

% 行内译注：橙色小字脚注
\newcommand{\zhnote}[1]{%
  \footnote{\textcolor{orange!70!black}{【译注】}#1}}

% 路线注：每块开头标原文页码
\newcommand{\routenote}[1]{%
  \par\noindent{\footnotesize\color{gray}［对应原文 #1］}\par\addvspace{.2em}}

% ---------- 常用记号 ----------
\newcommand{\Bun}{\mathrm{Bun}}
\newcommand{\Gr}{\mathrm{Gr}}
\newcommand{\Loc}{\mathrm{Loc}}
\newcommand{\Rep}{\mathrm{Rep}}
\renewcommand{\Vec}{\mathrm{Vec}}
\newcommand{\Perv}{\mathrm{Perv}}
\newcommand{\IC}{\mathrm{IC}}
\newcommand{\AS}{\mathrm{AS}}
\newcommand{\Kum}{\mathrm{Kum}}
\newcommand{\Swan}{\mathrm{Swan}}
\newcommand{\Kl}{\mathrm{Kl}}
\newcommand{\Hk}{\mathrm{Hk}}
\newcommand{\Hecke}{\mathcal{H}\mathrm{ecke}}
\newcommand{\Fl}{\mathrm{Fl}}
\newcommand{\Aut}{\mathrm{Aut}}
\newcommand{\Out}{\mathrm{Out}}
\newcommand{\Stab}{\mathrm{Stab}}
\newcommand{\Sat}{\mathcal{S}\mathrm{at}}
\newcommand{\semiSat}{\mathcal{S}}
\newcommand{\cO}{\mathcal{O}}
\newcommand{\cG}{\mathcal{G}}
\newcommand{\cS}{\mathcal{S}}
\newcommand{\cC}{\mathcal{C}}
\newcommand{\cT}{\mathcal{T}}
\newcommand{\cF}{\mathcal{F}}
\newcommand{\cL}{\mathcal{L}}
\newcommand{\cE}{\mathcal{E}}
\newcommand{\cH}{\mathcal{H}}
\newcommand{\ck}{\check}
\newcommand{\ckG}{\check G}
\newcommand{\ckT}{\check T}
\newcommand{\ckg}{\check{\mathfrak g}}
\newcommand{\ckb}{\check{\mathfrak b}}
\newcommand{\ckn}{\check{\mathfrak n}}
\newcommand{\LG}{{}^L G}
\newcommand{\frakg}{\mathfrak{g}}
\newcommand{\frakt}{\mathfrak{t}}
\newcommand{\frakb}{\mathfrak{b}}
\newcommand{\frakn}{\mathfrak{n}}
\newcommand{\frakh}{\mathfrak{h}}
\newcommand{\fraku}{\mathfrak{u}}
\newcommand{\frakm}{\mathfrak{m}}
\newcommand{\frakz}{\mathfrak{z}}
\newcommand{\ev}{\mathrm{ev}}
\newcommand{\pr}{\mathrm{pr}}
\newcommand{\ad}{\mathrm{ad}}
\newcommand{\Ad}{\mathrm{Ad}}
\newcommand{\mult}{\mathrm{mult}}
\newcommand{\inv}{\mathrm{inv}}
\newcommand{\St}{\mathrm{Std}}
\newcommand{\dd}{\mathrm{d}}
\newcommand{\sgn}{\mathrm{sgn}}
\newcommand{\CC}{\mathbb{C}}
\newcommand{\PP}{\mathbb{P}}
\newcommand{\GG}{\mathbb{G}}
\newcommand{\Aff}{\mathbb{A}}
\newcommand{\ZZ}{\mathbb{Z}}
\newcommand{\QQ}{\mathbb{Q}}
\newcommand{\FF}{\mathbb{F}}
\newcommand{\Ql}{\bar{\mathbb{Q}}_\ell}
\newcommand{\Wt}{\widetilde{W}}
\newcommand{\It}{I^t}
\newcommand{\Ip}{I^+}
\newcommand{\Satgeo}{\Sat^{\mathrm{geo}}}
\newcommand{\Grtriv}{\Gr^{\mathrm{triv}}}
\newcommand{\KlG}{\Kl_{\ckG}}
\newcommand{\Klg}{\Kl_{\mathfrak g}}
\newcommand{\red}{\mathrm{red}}
\newcommand{\der}{\mathrm{der}}
\newcommand{\loc}{\mathrm{loc}}
\newcommand{\rs}{r_s}
\newcommand{\rl}{r_\ell}

\pagestyle{plain}


\begin{document}

\begin{titlepage}
\centering
\vspace*{2.2cm}
{\Huge\bfseries 约化群的\\[0.35em] Kloosterman 层\par}
\vspace{1.0cm}
{\Large Kloosterman sheaves for reductive groups\par}
\vspace{1.6cm}
{\large Jochen Heinloth\quad Bao-Ch\^au Ng\^o\quad Zhiwei Yun\par}
\vspace{0.6cm}
{\normalsize \texttt{arXiv:1005.2765v1 [math.AG] 16 May 2010}\par}
\vspace{0.4cm}
{\normalsize 原文 58 页\par}
\vfill
{\large 译者：数译（DSH agent）\par}
\vspace{0.3cm}
{\large 版本 v1.0\quad 2026 年 11 月\par}
\vspace{1.2cm}
\end{titlepage}

\section*{译者说明}
\addcontentsline{toc}{section}{译者说明}

\textbf{原文性质。}这是 J. Heinloth、B.-C. Ng\^o 与 Z. Yun 的研究论文《Kloosterman sheaves for reductive groups》（arXiv:1005.2765v1，2010 年 5 月 16 日版，共 58 页；已发表于 \emph{Annals of Mathematics} \textbf{177} (2013), 241--310）。原文不是讲义也不是综述：它宣布并证明四条主定理（定理 1--4），另附四个技术附录，几乎每页都有对既有文献的引用跳步。译本以 arXiv v1 为底本，因为 v1 的编号（定理 1--4、按节混排的定理类与公式）与已发表版一致，便于读者按原文页码回查。

\textbf{编号约定。}译本完全沿用原文编号：四条主定理沿用原文的全局编号（定理 1、定理 2、定理 3、定理 4）；其余定理类条目（引理、命题、推论、定义、例、注、断言、记号、猜想）与编号公式\emph{共用}一个按节递增的计数器。因此 §2 中出现的是 2.1, 2.2, \dots, 2.14 这样定理与公式混排的号（2.13 是推论、2.14 是猜想）；§5 的 5.1--5.43 同样混排且编号连续。凡译本标注「(n.m)」处，其编号与原文逐条一致，可据以回查原文。原文以方括号缩写引用文献（如 [19]、[30]），译本一并保留，出处见文末参考文献。

\textbf{译注标记。}凡在译文中出现橙色小字的「【译注】」，都是译者所加，不属于原文；其用途限于：说明原文的排印错误或记号歧义、补充一个不解释便会误解的背景事实、标明某处推导在原文中确实被跳过的步骤。每条译注都尽量给出可核查的出处（原文献、nLab、Stacks Project 等）；若纯系译者的推断而无外部依据，则明确标为「译者按（无外部依据）」。译者不对原文作数学上的改写或补证。

\textbf{术语约定。}术语以本库《数语规范》已冻结条目为准（如 Hecke 特征层、局部系统、反常层、模叠、约化群、Langlands 对偶群、仿射 Grassmannian）；本库未收的术语（如 pinning、parahoric 子群、affine generic character、quasi-minuscule 表示、Swan 导子）在首次出现处给出定译与依据，并登记于工程文档的术语表。英文术语与汉字之间留半角空格。记号 $G$ 指原文的约化群，$\GG$（花体）指 $\PP^1$ 上的群概形，$\ckG$ 指 Langlands 对偶群。

\newpage
\tableofcontents
\newpage

% ============================================================
%  ch01 —— §0 引言（原文 p.1–p.6）
% ============================================================
\routenote{p.1--p.6}

\begin{center}
{\large\bfseries 摘\quad 要}
\end{center}

Deligne 在 $\PP^1-\{0,\infty\}$ 上构造了一个引人注目的局部系统，它附着于一族 Kloosterman 和。Katz 计算了它的单值性，并问：是否存在一般约化群的 Kloosterman 层？以及在 Langlands 对应下，这些局部系统应当附着于哪些自守形式？受 Gross 与 Frenkel--Gross 工作的启发，我们在几何 Langlands 纲领的框架内找到了一族显式的这类自守形式，甚至一族简单的自守层。我们用这些自守层以统一的方式对任意约化群构造 $\ell$-进 Kloosterman 层，并描述这些 Kloosterman 层的局部与整体单值性。特别地，它们给出具有例外单值群 $G_2$、$F_4$、$E_7$ 与 $E_8$ 的母题 Galois 表示。这也给出任意约化群带野分歧的几何 Langlands 对应的一个例子。

\setcounter{section}{-1}
\section{引言}

\subsection{Kloosterman 和与 Kloosterman 层的回顾}

设 $n$ 是正整数，$p$ 是素数。对每个 $\FF_q\supset\FF_p$ 与 $a\in\FF_q^\times$，$n$ 变量的 \emph{Kloosterman 和}定义为指数和
\[
\Kl_n(a;q):=(-1)^{n-1}\sum_{x_1x_2\cdots x_n=a,\ x_i\in\FF_q}
\exp\Big(\frac{2\pi i}{p}\mathrm{Tr}_{\FF_q/\FF_p}(x_1+x_2+\cdots+x_n)\Big).
\]
Kloosterman 和出现在模形式的 Fourier 系数中。它们满足 Weil 界
\begin{equation}
|\Kl_n(a;q)|\le n q^{(n-1)/2}.
\end{equation}
当 $n=2$ 时，对任意 $a\in\FF_p^\times$，定义角 $\theta(a)\in[0,\pi]$ 使
\[
2p^{\deg(a)/2}\cos\theta(a)=\Kl_2(a;p^{\deg(a)}),
\]
其中 $\deg(a)$ 是 $a$ 在 $\FF_p$ 上的次数。于是当 $\deg(a)$ 趋于 $\infty$ 时，角 $\{\theta(a)\mid a\in\FF_p^\times\}$ 按 Sato--Tate 测度 $\frac{2}{\pi}\sin^2\theta\,d\theta$ 在 $[0,\pi]$ 上等分布。一般地，关于 $\Kl_n(a;q)$ 的类似等分布定理由 Katz [25] 借助 Deligne 在 [11] 中的结果证明。

上述 Kloosterman 和的性质是借助一个层论的化身证明的。让我们回顾这个构造。在 [10] 中，Deligne 考虑图表
\[
\begin{tikzcd}
& \GG_m^n \arrow[dl, "\sigma"'] \arrow[dr, "\pi"] & \\
\GG_a & & \GG_m
\end{tikzcd}
\]
这里 $\GG_m$ 是乘法群，$\GG_a\cong\Aff^1$ 是加法群，而 $\sigma$（相应地 $\pi$）是取 $\GG_m^n$ 的 $n$ 个坐标之和（相应地之积）的映射。设 $\psi:\FF_p\to\mathbb{Q}_\ell(\mu_p)^\times$ 是非平凡特征（这里 $\mu_p$ 是 $p$ 次单位根之集），$\mathrm{AS}_\psi$ 是 $\FF_p$ 上加法群 $\GG_a$ 上相应的 Artin--Schreier 局部系统。Deligne 于是把 Kloosterman 层定义为 $\FF_p$ 上 $\GG_m$ 上如下带 $\mathbb{Q}_\ell(\mu_p)$-系数的层复形：
\[
\Kl_n:=R\pi_!\sigma^*\mathrm{AS}_\psi[n-1].
\]
固定一个嵌入 $\iota:\mathbb{Q}_\ell(\mu_p)\hookrightarrow\CC$，使 $\iota\psi(x)=\exp(2\pi i x/p)$ 对 $x\in\FF_p$ 成立。对任意 $a\in\GG_m(\FF_q)=\FF_q^\times$，记 $\mathrm{Frob}_a$ 为 $a$ 处的几何 Frobenius，它作用在几何茎 $(\Kl_n)_a$ 上。则由 Grothendieck--Lefschetz 迹公式有
\[
\Kl_n(a;q)=\iota\,\mathrm{Tr}\big(\mathrm{Frob}_a,(\Kl_n)_a\big).
\]
在这个意义上，$\Kl_n$ 是 Kloosterman 和族 $\{\Kl_n(a;q)\}_{a\in\FF_q^\times}$ 的层论化身。在 [10, Théorème 7.4, 7.8] 中 Deligne 证明了：

\begin{enumerate}[label=(\arabic*), leftmargin=1.6em, itemsep=.15em]
\item $\Kl_n$ 集中在次数 $0$，且是秩 $n$ 的局部系统。
\item $\Kl_n$ 在 $\{0\}$ 附近驯顺分歧，且单值是只有一个 Jordan 块的幂幺变换。
\item $\Kl_n$ 在 $\{\infty\}$ 附近完全野分歧（即 $\infty$ 处的野惯性群在 $\Kl_n$ 的茎上没有非零不动向量），且 Swan 导子 $\mathrm{Swan}_\infty(\Kl_n)=1$。
\item $\Kl_n$ 是纯权 $n-1$ 的（这蕴涵上文给出的 Weil 界估计）。
\end{enumerate}

在 [25, §8.7] 中，Katz 证明了 Kloosterman 层的唯一性：若 $\GG_m$ 上秩 $n$ 的局部系统 $L$ 满足上面列出的性质 (2) 与 (3)，则 $L$ 在 $\GG_m\otimes_{\FF_p}\bar{\FF}_p$ 上同构于 $\Kl_n$ 经 $\GG_m$ 上平移作用后的结果。此外在 [25, §11] 中，Katz 进一步研究了 $\Kl_n$ 的整体几何单值性。固定 $\GG_m$ 的一个几何点 $\eta$，记
\[
\phi:\pi_1(\GG_m,\eta)\to GL_n(\mathbb{Q}_\ell(\mu_p))
\]
为与 Kloosterman 层 $\Kl_n$ 相伴的单值表示。设 $\phi^{\mathrm{geo}}$ 是 $\phi$ 到 $\pi_1(\GG_m\otimes_{\FF_q}\bar{\FF}_q,\eta)$ 的限制。Katz 定出了 $\phi^{\mathrm{geo}}$ 的像的 Zariski 闭包 $\ckG^{\mathrm{geo}}$：
\begin{equation}
\ckG^{\mathrm{geo}}=
\begin{cases}
\mathrm{Sp}_n, & n \text{ 偶},\\
\mathrm{SL}_n, & n \text{ 奇},\ p \text{ 奇},\\
\mathrm{SO}_n, & n \text{ 奇},\ n\ne 3,\ p=2,\\
G_2, & n=7,\ p=2.
\end{cases}
\end{equation}
借助 Deligne 在 [11] 中的定理，他由此推出 Frobenius 元 $\mathrm{Frob}_a$ 的共轭类按一个「Sato--Tate」测度等分布。

\subsection{本文的动机与目标}

鉴于例外群 $G_2$ 作为整体几何单值群的神秘出现，Katz 问道（[26] p. I-5）：是否所有半单群都会作为 $\GG_m$ 上局部系统的几何单值群出现？换言之，是否存在指数和，其等分布律由任意的单群、特别是例外群支配？

在本文中我们找到这类局部系统的一个统一构造。对任意分裂约化群 $\ckG$，我们将在 $\GG_m=\PP^1\setminus\{0,\infty\}$ 上构造一个 $\ckG$-局部系统 $\Kl_{\ckG}$，其局部分歧与 $\Kl_n$ 类似；而当 $\ckG=GL_n$ 时，我们恢复 Deligne 的 Kloosterman 层 $\Kl_n$。我们将定出其整体几何单值性的 Zariski 闭包（它结果是「大的」），证明该层的纯度并推出等分布律。最后，我们给出关于这类 Galois 表示（或局部系统）唯一性的一个猜想。

为引言的目的，让我们限制到如下情形。设 $G$ 是 $k=\FF_q$ 上分裂、几乎单且单连通的群，或 $G$ 是 $k$ 上的 $GL_n$。设 $\ckG$ 是其在 $\mathbb{Q}_\ell(\mu_p)$ 上的 Langlands 对偶群，即一个分裂、单且伴随的群，或 $GL_n$。

我们构造的动机来自有理函数域 $K=k(t)$ 的 Langlands 对应。Katz [25] 在其研究 Kloosterman 层时已经提出：

\begin{quote}
「……把这些结果与 Langlands 哲学所给出的、用自守形式对这些层所作的猜想性描述相比较，将会是很有意思的。」
\end{quote}

在一系列工作 [20]、[21]、[22]、[15] 中，Gross（部分与 Reeder、与 Frenkel 合作）提出了 $G(\mathbb{A}_K)$ 的一个候选自守表示 $\pi$；在 $G=GL_n$ 的情形，它应当给出 Kloosterman 层 $\Kl_n$。让我们简要回顾他们的工作。

在 [20] 中，Gross 与 Reeder 给出了 $G(k((s)))$ 的表示 $V_\varphi$ 的如下构造。固定 $G$ 的一个 Borel 子群 $B$，记 $U$ 为 $B$ 的幂幺根。记 $I(0):=\{g\in G(k[[s]])\mid g\bmod s\in B\}$ 为 $G(k[[s]])$ 的 Iwahori 子群，$I(1):=\{g\in G(k[[s]])\mid g\bmod s\in U\}$ 为 $I(0)$ 的幂幺根。对任意仿射一般特征 $\varphi:I(1)\to\mathbb{Q}_\ell(\mu_p)$（见 §1.3），令 $V_\varphi:=\mathrm{c\text{-}Ind}_{I(1)\times Z(G)}^{G(k((s)))}(\varphi\otimes 1)$ 为 $G(k((s)))$ 的紧诱导表示。Gross 与 Reeder 证明这个表示不可约且超尖。

对任意整体域 $F$，Gross [21] 设法用迹公式得到了 $G(\mathbb{A}_F)$ 的自守表示 $\pi$ 的重数表达式，这些 $\pi$ 的分歧局部成分或是 Steinberg 表示，或是表示 $V_\varphi$。特别地，当整体域是 $K=k(t)$ 时，这个公式蕴涵：对任意半单单连通群 $G$，存在唯一的尖自守表示 $\pi=\pi(\varphi)$，满足

\begin{enumerate}[label=(\arabic*), leftmargin=1.6em, itemsep=.15em]
\item $\pi$ 在 $\{0,\infty\}$ 之外非分歧；
\item $\pi_0$ 是 $G(k((t)))$ 的 Steinberg 表示；
\item $\pi_\infty$ 是 $G(k((s)))$ 的简单超尖表示 $V_\varphi$，其中 $s=t^{-1}$。
\end{enumerate}

受 Langlands 哲学启发，Gross [22] 提出了如下猜想：自守表示 $\pi=\pi(\varphi)$ 应当对应于 $\GG_m=\PP^1\setminus\{0,\infty\}$ 上的一个 $\ckG$-局部系统 $\Kl_{\ckG}(\varphi)$——我们称之为与 $\ckG$ 及 $\varphi$ 相伴的 Kloosterman 层——它具有与 $\pi$ 平行的如下性质：

\begin{enumerate}[label=(\arabic*), leftmargin=1.6em, itemsep=.15em]
\item $\Kl_{\ckG}(\varphi)$ 是 $\GG_m=\PP^1\setminus\{0,\infty\}$ 上的 $\ckG$-局部系统。
\item $\Kl_{\ckG}(\varphi)$ 在 $\{0\}$ 附近驯顺分歧，且单值是 $\ckG$ 中的正则幂幺元。
\item 与 $\ckG$ 的伴随表示相伴的局部系统 $\Kl^{\mathrm{Ad}}_{\ckG}(\varphi)$（它是 $\PP^1\setminus\{0,\infty\}$ 上秩为 $\dim\ckG$ 的局部系统）在 $\{\infty\}$ 附近完全分歧（即 $\infty$ 处的惯性群没有非零不动向量），且 $\mathrm{Swan}_\infty(\Kl^{\mathrm{Ad}}_{\ckG}(\varphi))=r(\ckG)$，即 $\ckG$ 的秩。
\item 对 $\ckG$ 的任意不可约表示 $V$，相应的局部系统 $\Kl^V_{\ckG}(\varphi)$ 是纯的。
\end{enumerate}

对 $\ckG=GL_n$，层 $\Kl_n$ 具有上述性质。对 $A$、$B$、$C$ 与 $G_2$ 型的 $\ckG$，具有上述性质的 $\ckG$-局部系统此前已由 Katz [26] 用逐情形构造、并以 $\Kl_n$ 为基本构件造出。

在 [15] 中，Gross 与 Frenkel 在复数上构造了这类局部系统的类比物。确切地说，他们定义了 $\PP^1_\CC\setminus\{0,\infty\}$ 上的 $\ckG_\CC$-联络 $\nabla^{\ckG}$（依赖于在 $\check{\mathfrak g}$ 的仿射单根空间中选取非零向量的选择）。这个联络在 $\{0\}$ 处有正则奇点、在 $\infty$ 处有不正则奇点，与上面的性质 (2) 和 (3) 平行。他们进一步证明 $\nabla^{\ckG}(\check X)$ 在上同调意义下刚性 [15, Theorem 1]，并计算了这些联络的微分 Galois 群。对 $A$、$B$、$C$ 与 $G_2$ 型的 $\ckG$，他们验证 $\nabla^{\ckG}$ 与 Katz [26] 构造的联络一致——后者是上述 Katz 的 $\ell$-进 $\ckG$-局部系统的类比物。所有这些都给出强有力的证据，表明 $\nabla^{\ckG}(\check X)$ 应当是猜想中的局部系统 $\Kl_{\ckG}(\varphi)$ 的正确 de Rham 类比物。

[15] 中关于猜想性局部系统 $\Kl_{\ckG}(\varphi)$ 的预测，是我们工作的指导方针。

\subsection{构造方法}

我们对 Kloosterman 层的构造可以概括如下。

我们从上面提到的自守表示 $\pi$ 出发。关键的观察是：$\pi$ 含有一个 Hecke 特征函数 $f_\varphi$，它可以利用双 coset
\[
G(K)\backslash G(\mathbb{A}_K)\Big/I^0\times I^\infty(2)\times\prod_{x\ne0,\infty}G(\mathcal O_x)
\]
的组合学显式写出（§2.1）。这个双 coset 的点是 $\PP^1$ 上带特定水平结构的 $G$-丛的模叠 $\Bun_G(0,2)$ 的有理点，于是我们可以把函数 $f_\varphi$ 升级为此叠上的层 $A_\varphi$。然后我们证明这个层确实是 Hecke 特征层。特别地，我们证明：几何 Hecke 算子作用在 $A_\varphi$ 上所得的特征值定义了 $\PP^1\setminus\{0,\infty\}$ 上的一个 $\ckG$-局部系统——即我们的 Kloosterman 层 $\Kl_{\ckG}(\varphi)$。在技术上，这个结果的证明依赖于如下事实：层 $A_\varphi$ 是某个具有仿射支撑的局部系统的 clean 扩张。

在 $G=GL_n$ 的情形，我们显式算出局部系统 $\Kl_{\ckG}(\varphi)$。对特定的 $\varphi$，与 $\Kl_{\ckG}(\varphi)$ 及 $\ckG=GL_n$ 的标准表示相伴的秩 $n$ 局部系统，结果正是 Deligne 的 Kloosterman 层 $\Kl_n$。这里 Deligne 的图表将作为几何 Hecke 变换的一部分重新出现。

当把 $G$ 换成一类特定的、$\PP^1\setminus\{0,\infty\}$ 上的拟分裂群概形 $\mathbb{G}$ 时，我们的构造同样有效；此外还有一个依赖于额外 Kummer 特征的变体。这些推广看起来正是比较不同群的 Kloosterman 层时所需的自然框架。由于这些推广并不需要额外的论证，我们将在更一般的框架下给出构造。

\subsection{Kloosterman 层的性质}

本文最长的一部分（§4–§6）随后用于研究单伴随群 $\ckG$ 的 Kloosterman 层 $\Kl_{\ckG}(\varphi)$ 的局部与整体性质。

首先，借助 [11]，所期望的纯度性质 (4) 是我们构造的推论。当 $\ckG$ 是伴随群时，我们还可以把 $\Kl^V_{\ckG}(\varphi)$ 归一化为纯权 $0$。于是，像前面一样固定嵌入 $\iota:\mathbb{Q}_\ell(\mu_p)\hookrightarrow\CC$，指数和
\[
\Kl^V_{\ckG}(\varphi;a;q)=\iota\,\mathrm{Tr}\big(\mathrm{Frob}_a,(\Kl^V_{\ckG}(\varphi))_a\big),\qquad a\in\FF_q^\times,\ V\in\Rep(\ckG)
\]
满足 Weil 界
\[
|\Kl^V_{\ckG}(\varphi;a;q)|\le\dim V.
\]
关于 $0$ 处单值性的性质 (2)（见定理 1(2)）将在 §4.3 中证明。关于 $\infty$ 处单值性的性质 (3)（见定理 2）将在 §5 中证明。Swan 导子的计算涉及对仿射 Grassmannian 中某些 Schubert 簇的几何作细致研究，这占据了 §5 的很大一部分。利用 Gross--Reeder 的一个结果 [20]，可以给出 $\infty$ 处单值性的显式描述（推论 2.13）。这些结果合起来表明，局部系统 $\Kl_{\ckG}(\varphi)$ 满足 Gross 所期望的性质 (1)--(4)。

再设
\[
\phi:\pi_1(\PP^1\setminus\{0,\infty\},\eta)\to\ckG(\mathbb{Q}_\ell(\mu_p))
\]
为与 Kloosterman 层 $\Kl_{\ckG}(\varphi)$ 相伴的单值表示，$\phi^{\mathrm{geo}}$ 为其到 $\pi_1(\PP^1\setminus\{0,\infty\}\otimes_k\bar k,\eta)$ 的限制。我们发现（定理 3）：当 $p>2$（且 $\ckG=B_3$ 时 $p>3$）时，$\phi^{\mathrm{geo}}$ 的像的 Zariski 闭包 $\ckG^{\mathrm{geo}}$ 与 Frenkel--Gross 在 [15, Cor. 9, 10] 中算出的 $\nabla^{\ckG}(\check X)$ 的微分 Galois 群 $\ckG^{\nabla}$ 一致；我们在表 1 中列出它。

\begin{table}[htbp]
\centering
\small
\caption{$\Kl_{\ckG}(\varphi)$ 的整体几何单值群}\label{tab:1}
\begin{tabular}{@{}ll@{}}
\toprule
$\ckG$ & $\ckG^{\mathrm{geo}}$\\
\midrule
$A_{2n}$ & $A_{2n}$\\
$A_{2n-1}$, $C_n$ & $C_n$\\
$B_n$, $D_{n+1}$ ($n\ge4$) & $B_n$\\
$E_7$ & $E_7$\\
$E_8$ & $E_8$\\
$E_6$, $F_4$ & $F_4$\\
$B_3$, $D_4$, $G_2$ & $G_2$.\\
\bottomrule
\end{tabular}
\end{table}

利用 Katz 的论证（[25, §3, §13]），这蕴涵关于共轭类 $\phi(\mathrm{Frob}_a)$ 的如下等分布律。

\medskip
\noindent\textbf{推论.} 设 $\ckG$ 单且伴随（此时 $\phi$ 的整个像都落在 $\ckG^{\mathrm{geo}}$ 中）。设 $\ckG^{\mathrm{geo},c}\subset\ckG^{\mathrm{geo}}(\CC)$ 是一个紧实形式（利用嵌入 $\iota:\mathbb{Q}_\ell(\mu_p)\hookrightarrow\CC$ 使 $\ckG^{\mathrm{geo}}(\CC)$ 有意义）。记 $\ckG^{\natural}_{\mathrm{geo},c}$ 为 $\ckG^{\mathrm{geo},c}$ 的共轭类之集。

对 $a\in k^\times$，$\phi(\mathrm{Frob}_a)$ 在 $\ckG^{\mathrm{geo}}(\CC)$ 中的共轭类事实上属于 $\ckG^{\natural}_{\mathrm{geo},c}$。当 $\deg(a)$（即域扩张 $k(a)/k$ 的次数）趋于 $\infty$ 时，共轭类
\[
\{\phi(\mathrm{Frob}_a)\mid a\in k^\times\}\subset\ckG^{\natural}_{\mathrm{geo},c}
\]
按 Haar 测度从 $\ckG^{\mathrm{geo},c}$ 到 $\ckG^{\natural}_{\mathrm{geo},c}$ 的推前而等分布。
\medskip

\subsection{未决问题}

与关于 $\Kl_n$ 的已知结果相比，对 $\Kl_{\ckG}(\varphi)$ 唯一缺少的一环是唯一性。我们证明 $\Kl_{\ckG}(\varphi)$ 在上同调意义下刚性，即
\[
H^1\big(\PP^1,j_{!*}\Kl^{\mathrm{Ad}}_{\ckG}(\varphi)\big)=0.
\]
这给出 $\Kl_{\ckG}(\varphi)$ 的物理刚性的证据：任何其他满足性质 (2) 与 (3) 的 $\PP^1\setminus\{0,\infty\}$ 上的局部系统 $L$，都应当在 $\PP^1\setminus\{0,\infty\}\otimes_k\bar k$ 上同构于 $\Kl_{\ckG}(\varphi)$。我们将在猜想 7.1 与 7.2 中陈述确切的唯一性猜想。

表 1 还提示：表中同一行出现的不同群 $\ckG$ 的 Kloosterman 层应当本质上是同一个（见猜想 7.3）。这可以看作 Kloosterman 层的函子性陈述。

此外，Kloosterman 层的单值性只是在分裂群的情形被详细研究过。在 §7.2 中，我们陈述关于 $\Kl_{^LG}(\varphi)$ 的局部与整体单值性的预测，其中 $\mathbb{G}$ 是 $\PP^1$ 上拟分裂、单且单连通的群概形，在 $\{0,\infty\}$ 之外的所有位处有良好约化，并由一个驯 Kummer 覆盖 $[N]:\PP^1\setminus\{0,\infty\}\to\PP^1\setminus\{0,\infty\}$（$z\mapsto z^N$）分裂。

\subsection{全文结构}

在 §1 中，我们定义 $\PP^1$ 上各种群概形 $\mathbb{G}(m_0,m_\infty)$，它们编码了 $G$-丛模叠的水平结构——自守层就定义在这些模叠上。我们还给出这些模叠的几何描述。在 §2 中，我们构造自守层 $A_\varphi$ 并陈述本文的主要结果。证明在 §4–§6 中给出。§3 专门处理 $G=GL_n$ 的情形，我们在此恢复 Deligne 的经典 Kloosterman 层。在 §7 中，我们陈述关于 Kloosterman 层刚性的猜想，以及关于拟分裂群的 Kloosterman 层的预期。

本文包含四个附录。在 §8 中，我们在拟分裂群的一般框架下证明模空间 $\Bun_G$ 的结构定理。在 §9 中，我们证明外自同构与几何 Satake 等价之间的相容性。这是描述 Kloosterman 层的整体单值像所需要的，而我们未能找到该结果的文献出处。在 §10 中，我们收集关于 $\ckG$ 的拟极小表示的一些事实。在 §11 中，我们分析 $G_2$ 的伴随 Schubert 簇的几何——它不适合其他类型群的统一处理。

\medskip
\noindent\textbf{致谢.} 作者们 2009--2010 年在高等研究院（IAS）逗留期间相遇并开始了本项目。他们感谢 IAS 提供的一流工作条件。

% ============================================================
%  ch02 —— §1 结构群（原文 p.7–p.11）
% ============================================================
\routenote{p.7--p.11}

\section{结构群}

我们将在特征为 $p$ 的固定有限域 $k$ 上工作。固定底曲线 $\PP^1$ 的坐标 $t$，使得 $\Aff^1=\operatorname{Spec}(k[t])\subset\PP^1$。记 $s:=t^{-1}$ 为 $\infty\in\PP^1$ 附近的坐标。

对任意闭点 $x\in\PP^1$，记 $\mathcal O_x$ 为 $x$ 处的完全局部环，$K_x$ 为 $\mathcal O_x$ 的分式域。

由于我们关心的是（野）分歧局部系统的几何 Langlands 对应，我们需要考虑带各种水平结构的主丛。把这些主丛看作 $\PP^1$ 上群概形 $\mathbb{G}$ 的 torsor 会很方便。此外，为了陈述猜想中的函子性，允许拟分裂群概形是有益的。我们将分几步引入这些群概形。

\subsection{$\PP^1\setminus\{0,\infty\}$ 上的拟分裂群概形}

我们假设 $\mathbb{G}|_{\PP^1\setminus\{0,\infty\}}$ 是拟分裂约化群。此外假设存在有限扩张 $k'/k$ 与整数 $N$，满足 $(N,p)=1$，使得 $\mathbb{G}$ 在驯扩张 $[N]:\GG_{m,k'}\to\GG_{m,k'}$（由 $t\mapsto t^N$ 定义）上分裂。我们可以且将会假设 $k'$ 包含所有 $N$ 次单位根。记 $\mu_N$ 为 $N$ 次单位根之群。

固定子群 $S\subset T\subset B\subset G$，其中 $S$ 是极大分裂环面，$T$ 是极大环面，$B$ 是 Borel 子群。我们还固定 $G$ 的一个拟钉住（quasi-pinning）。

为描述我们的群概形，为简单起见假设 $k=k'$。一般情形下，我们给出的构造在 $\operatorname{Gal}(k'/k)$ 下不变，于是 Galois 下降即给出一般情形。

由 [13]（Exp.~XXIV, Thm~3.11），这些群可以如下显式描述。按假设存在 $k$ 上分裂约化群 $G$，使得 $[N]^*\mathbb{G}|_{\GG_m}=G\times\GG_m$。覆盖 $[N]$ 的自同构群为 $\mu_N$。固定 $T\subset B\subset G$ 为分裂极大环面与 Borel 子群，以及 $G$ 的一个钉住 $\dagger$。

这些数据定义一个态射 $\sigma:\mu_N\to\operatorname{Aut}^\dagger(G)$，其中 $\operatorname{Aut}^\dagger(G)$ 是保持钉住的 $G$ 的自同构群。记 $[N]^*(G\times\GG_m)$ 为 $G\times\GG_m$ 沿覆盖 $[N]$ 的 Weil 限制。则 $\mu_N$ 通过 $\sigma$ 在 $G$ 上的作用与在覆盖 $[N]$ 上的作用而作用于这个 Weil 限制。我们仍把对角作用记为 $\sigma$。于是下降给出：
\[
\mathbb{G}|_{\PP^1\setminus\{0,\infty\}}\cong\big([N]^*(G\times\PP^1\setminus\{0,\infty\})\big)^\sigma
=(G\times\GG_m)/(\mu_N).
\]

\subsection{$0$ 与 $\infty$ 处的水平结构}

接下来需要描述 $\mathbb{G}|_{\PP^1\setminus\{0,\infty\}}$ 到 $\PP^1$ 的若干扩充。仍记 $[N]:\PP^1\to\PP^1$ 为由 $t\mapsto t^N$ 给出的映射。

群 $\mathbb{G}|_{\PP^1\setminus\{0,\infty\}}$ 将被扩充为由凸函数（在所有根上取 $0$）定义的（特殊）Bruhat--Tits 群 $\mathbb{G}$。由于我们假设 $\mathbb{G}|_{\PP^1\setminus\{0,\infty\}}$ 在驯扩张上分裂，这个群可以描述为群 $\big([N]^*(G\times\PP^1)\big)^\sigma$ 的连通分支。子群 $S\subset T|_{\PP^1\setminus\{0,\infty\}}\subset B|_{\PP^1\setminus\{0,\infty\}}\subset\mathbb{G}|_{\PP^1\setminus\{0,\infty\}}$ 定义闭子群 $S\subset\mathbb{T}\subset\mathbb{B}\subset\mathbb{G}$。

我们将需要在 $0$ 与 $\infty$ 处引入水平结构。它们对应于 $\mathbb{G}$ 的 Moy--Prasad 滤链的最初几步。关于这些群概形构造的自足论述可见 Yu 的预印本 [36]。

先考虑环面。对分裂环面 $T=\GG_m^r\times\PP^1$，记 $T(m_0,m_\infty)$ 为 $\PP^1$ 上光滑群概形，使得
\begin{gather*}
T(m_0,m_\infty)|_{\PP^1\setminus\{0,\infty\}}=\GG_m^r\times\PP^1\setminus\{0,\infty\},\\
T(m_0,m_\infty)(\mathcal O_{\PP^1,0})=\{g\in T(\mathcal O_{\PP^1,0})\mid g\equiv1\bmod t^{m_0}\},\\
T(m_0,m_\infty)(\mathcal O_{\PP^1,\infty})=\{g\in T(\mathcal O_{\PP^1,\infty})\mid g\equiv1\bmod s^{m_\infty}\}.
\end{gather*}
这也为诱导环面定义了滤链。对任意环面 $T$，取嵌入到诱导环面 $T\hookrightarrow I$，并通过拉回 $I$ 上的滤链来定义滤链。由 [36] 第 4 节，这个定义与所取嵌入无关；且由于我们假设 $T$ 在驯扩张上分裂，这些子群定义了连通群。

对约化群 $\mathbb{G}|_{\PP^1\setminus\{0,\infty\}}$，我们要为 $m_0,m_\infty\in\{0,1,2\}$ 定义模型 $\mathbb{G}(m_0,m_\infty)$，它们对应于 $0$ 处（相应地 $\infty$ 处）Iwahori 子群的 Moy--Prasad 滤链的第 $m_0$ 步（相应地第 $m_\infty$ 步）。

先假设 $\mathbb{G}=G\times\PP^1$ 是分裂半单群。像前面一样，固定 $S=T\subset B\subset G$ 为分裂极大环面、Borel 子群以及 $G$ 的钉住。记 $U\subset B$ 为幂幺子群。记 $\Phi=\Phi(G,S)$ 为根系，$\Phi^\pm\subset\Phi$ 为相对于所选 Borel 子群的正、负根之集。记 $\alpha_1,\dots,\alpha_r$ 为正单根，$\{\theta_j\}$ 为 $G$ 的最高根。最后，对每个根 $\alpha$，记 $U_\alpha$ 为 $G$ 中相应的根子群。

我们考虑 $G(k[[s]])$ 的如下有界子群：
\begin{itemize}[leftmargin=1.8em, itemsep=.2em]
\item $I(0):=\{g\in G(k[[s]])\mid g\bmod s\in B\}$ 是 Iwahori 子群。
\item $I(1):=\{g\in G(k[[s]])\mid g\bmod s\in U\}$ 是 $I(0)$ 的幂幺根。
\end{itemize}
为描述 $I(2)$，令 $f:\Phi\cup\{0\}\to\mathbb N$ 为由 $f(0)=1$ 以及
\[
f(\alpha)=
\begin{cases}
0 & \text{若 } \alpha\in\Phi^+\setminus\{\alpha_i\} \text{ 为正但非单根},\\
1 & \text{若 } \alpha=\alpha_i \text{ 是正单根},\\
1 & \text{若 } \alpha\in\Phi^-\setminus\{-\theta_j\},\\
2 & \text{若 } \alpha=-\theta_j \text{ 是某个最高根的负根}
\end{cases}
\]
定义的凹函数。$I(2)\subset G(k[[s]])$ 表示由凹函数 $f$ 定义的有界子群，即由 $\{u\in U_\alpha\mid u\equiv1\bmod s^{f(i)}\}$ 与 $\{g\in T(k[[s]])\mid g\equiv1\bmod s\}$ 生成的子群。

注意按定义 $I(1)/I(2)\cong\oplus_{\alpha\ \text{单仿射}}\GG_a$，同构于那些使 $\alpha$ 为单仿射根的根子群 $U_\alpha\subset G(k[[s]])$ 的直和。

类似地，我们用对偶 Borel 子群 $B^{\mathrm{opp}}$ 代替上述定义中的 Borel 子群，定义 $I(i)^{\mathrm{opp}}\subset G(k[[t]])$ 为相应的群。

\medskip
\noindent\textbf{例.} 对 $G=SL_n$，取 $S=T$ 为对角矩阵，$B$ 为上三角矩阵。则子群 $I(0)\subset SL_n(k[[s]])$ 由下三角项被 $s$ 整除的矩阵组成；$I(1)\subset I(0)$ 由对角项属于 $1+sk[[s]]$ 的矩阵组成。

$SL_n$ 关于单根的根子群由上对角项 $a_{i,i+1}$ 给出，而最长根之负根的 $U_{-\tilde\alpha}$ 由左下角的项给出。于是 $I(2)$ 由形如
\[
\begin{pmatrix}
1+sk[[s]] & sk[[s]] & k[[s]] & k[[s]]\\
sk[[s]] & 1+sk[[s]] & sk[[s]] & k[[s]]\\
sk[[s]] & sk[[s]] & 1+sk[[s]] & sk[[s]]\\
s^2k[[s]] & sk[[s]] & sk[[s]] & 1+k[[s]]
\end{pmatrix}
\]
的矩阵组成。把矩阵 $(a_{ij})$ 映到各项 $a_{i,i+1}$ 与 $a_{n,1}$ 的首项系数，得到同构 $I(1)/I(2)\cong\Aff^n$。

对一般分裂约化群 $\mathbb{G}=G\times\PP^1$，考虑导群 $G^{\mathrm{der}}$ 与中心的连通分支 $Z(G)^\circ$——后者是环面。暂时记 $I_{G^{\mathrm{der}}}(i)$、$I^{\mathrm{opp}}_{G^{\mathrm{der}}}(i)$ 为上面为半单群 $G^{\mathrm{der}}$ 定义的群，并定义
\[
I(i):=Z(G)^\circ(i)(k[[s]])\cdot I_{G^{\mathrm{der}}}(i)\subset G(k[[s]]),
\]
类似地 $I(i)^{\mathrm{opp}}=Z(G)^\circ(i,m_\infty)(k[[t]])\,I^{\mathrm{opp}}_{G^{\mathrm{der}}}(i)$。

对分裂约化群 $\mathbb{G}=G\times\PP^1$，群 $\mathbb{G}(m_0,m_\infty)$ 表示 Bruhat--Tits 群概形，使得
\begin{gather*}
\mathbb{G}(m_0,m_\infty)|_{\PP^1\setminus\{0,\infty\}}=G\times\GG_m,\\
\mathbb{G}(m_0,m_\infty)(\mathcal O_\infty)=I(m_\infty),\qquad
\mathbb{G}(m_0,m_\infty)(\mathcal O_0)=I(m_0)^{\mathrm{opp}}.
\end{gather*}
最后，对一般拟分裂群定义
\[
\mathbb{G}(m_0,m_\infty):=\Big(\big([N]^*\mathbb{G}(m_0,m_\infty)\big)^\sigma\Big)^\circ.
\]
注意对环面这并不给出新的定义。

我们将缩写：
\[
I(m_\infty):=\mathbb{G}(m_0,m_\infty)(\mathcal O_\infty),\qquad
I^-(m_0):=\mathbb{G}(m_0,m_\infty)(\PP^1-\{\infty\}).
\]
回忆 [32]：群 $I(m)$ 具有 $k$ 上（无限维）群概形的自然结构。

\subsection{仿射一般特征}

如前面所述，我们的构造依赖于 $I(1)$ 的一个特征的选择。称线性函数 $\varphi:I(1)/I(2)\to\Aff^1$ 为\emph{一般}的，若对任意单仿射根 $\alpha$，$\varphi$ 在 $U_\alpha$ 上的限制非平凡。全文固定这样一个一般 $\varphi$。

固定非平凡加性特征 $\psi:\FF_p\to\mathbb{Q}_\ell$，并把由 $\psi\circ\mathrm{Tr}_{k/\FF_p}$ 定义的 $k$ 上特征仍记为 $\psi$。用这个记号，特征 $\psi\circ\varphi$ 称为 $I(1)$ 的\emph{仿射一般特征}。

\subsection{主丛}

定义了群之后，我们需要收集关于 $G$-丛模叠 $\Bun_G$ 的几何的一些基本结果。所有这些对常群都是熟知的（例如可参看 [14] 的新近论述）。为把这些结果推广到我们的设定，我们依赖 [32] 与 [23]，其中解释了扭 loop 群的相应结果。需要指出的是，除 $\Bun_G$ 连通分支的计算外，所有结果都是关于 $\PP^1$ 上、在驯分歧覆盖 $[N]:\PP^1\to\PP^1$ 上分裂的群概形的特有问题。

首先回忆 [23] 中的一些结果与记号。记 $N$ 为 $T$ 的正规化子，$\widetilde W:=N(k((s)))/T(k[[s]])$ 为 Iwahori--Weyl 群。进一步记 $W_0:=N(k((s)))/T(k((s)))$ 为相对 Weyl 群。它同构于特殊纤维 $\mathbb{G}_{\infty,\mathrm{red}}$ 的约化商群的 Weyl 群（同处 Proposition 13），且
\[
\widetilde W\cong X_*(T)^{\pi_1(\PP^1\setminus\{0,\infty\})}\ltimes W_0.
\]
记 $W_a$ 为 $G(k((s)))$ 的根系（的）仿射 Weyl 群，它可以等同于 $G$ 的导群的单连通覆盖的 Iwahori--Weyl 群（[23] p.~196）。这是一个 Coxeter 群。于是 $\widetilde W\cong W_a\ltimes\Omega$，其中 $\Omega\cong X_*(Z(\ckG)^{\pi_1(\GG_m)})$ 是一个 alcove 的稳定子。

最后，对 $x\in\PP^1-\{\infty\}$，记 $\Gr_{G,x}$ 为仿射 Grassmannian（见 [32]），使得 $\Gr_{G,x}(k)=G(K_x)/G(\mathcal O_x)$。它也可以定义为参数化 $\PP^1$ 上 $G$-丛 $P$ 连同平凡化 $\phi:P|_{\PP^1\setminus\{x\}}\xrightarrow{\sim}\mathbb{G}_{\PP^1\setminus\{x\}}$ 的 ind-概形（例如 [24]）。

\begin{proposition}
假设基域 $k$ 或是有限域，或是代数闭域。设 $\mathbb{G}$ 为 $\PP^1$ 上 §1 中定义的拟分裂群概形。则下列结论成立：
\begin{enumerate}[label=(\arabic*), leftmargin=1.8em, itemsep=.2em]
\item 对任意 $x\in\PP^1$，典范映射 $\Gr_{G,x}\to\Bun_G$ 在 $\Bun_G$ 的光滑拓扑下局部有截面。此外，该映射在 $k$-点上本质满，即：$\PP^1$ 上任意 $G$-丛到 $\PP^1\setminus\{x\}$ 的限制都是平凡的。
\item $\pi_0(\Bun_G)\cong\pi_1\big([N]^*\mathbb{G}|_{\GG_m}\big)^{\pi_1(\GG_m)}\cong\Omega$。
\item $\PP^1$ 上每个 $G$-丛都可约化到 $T$。
\item （Birkhoff--Grothendieck 分解）
\[
\mathbb{G}(0,0)(k((s)))=\bigsqcup_{w\in\widetilde W}I^-(0)\cdot w\cdot I(0).
\]
\end{enumerate}
\end{proposition}

若 $\mathbb{G}=G\times\PP^1$ 是分裂群，本结果的证明可在 [14] 中找到。由于分裂群情形是我们的出发点，一般情形的证明将推迟到附录。对 $\gamma\in\Omega$，记 $\Bun^\gamma_{G(m,n)}$ 为 $\Bun_{G(m,n)}$ 的相应连通分支。

让我们收集该结果的一些推论。

\begin{corollary}
设 $n\in\{0,1,2\}$，$m\in\{0,1\}$。任意 $\gamma\in\Omega$ 定义同构 $\mathrm{Hk}_\gamma:\Bun_{G(m,n)}\to\Bun_{G(m,n)}$。它诱导连通分支之间的同构 $\Bun^0_{G(m,n)}\to\Bun^\gamma_{G(m,n)}$。特别地，$\Bun_{G(m,n)}$ 的所有连通分支互相同构。
\end{corollary}

\begin{proof}
由命题的第 (2) 部分，$\Bun_{G(m,n)}$ 的连通分支由 $\Omega\subset\widetilde W$ 编号。由于 $\Omega$ 是定义 $I(0)^{\mathrm{opp}}$ 的 alcove 的稳定子，任意 $\gamma\in\Omega$ 都正规化 $I(0)^{\mathrm{opp}}$ 及其幂幺根 $I(1)^{\mathrm{opp}}$。于是 $\Omega$ 通过右乘作用在 $\Gr_{G(m,n),0}=G(k((t)))/I(m)^{\mathrm{opp}}$ 上。这定义同构 $\mathrm{Hk}_\gamma:\Bun_{G(m,n)}\to\Bun_{G(m,n)}$。此外，[23] Lemma 14 蕴涵该映射在 $\Bun^0_{G(m,n)}$ 上的限制诱导同构 $\Bun^0_{G(m,n)}\to\Bun^\gamma_{G(m,n)}$。
\end{proof}

\begin{corollary}\label{cor:1.3}
\begin{enumerate}[label=(\arabic*), leftmargin=1.8em, itemsep=.2em]
\item 记 $\mathbb{G}(\PP^1)$ 为平凡 $G$-丛的自同构群。则包含 $\mathrm B\mathbb{G}(\PP^1)\to\Bun_G$ 是仿射开嵌入。
\item 记 $T(\PP^1)=H^0(\PP^1,T)=H^0(\PP^1,\mathbb{G}(0,0))$ 为平凡 $\mathbb{G}(0,0)$-丛的自同构群。平凡丛的包含定义仿射开嵌入 $\mathrm B(T(\PP^1))\hookrightarrow\Bun_{\mathbb{G}(0,0)}$。
\item 平凡 $\mathbb{G}(0,1)$-丛定义仿射开嵌入 $\operatorname{Spec}(k)\hookrightarrow\Bun^0_{\mathbb{G}(0,1)}$。
\item 把 $I(1)/I(2)$ 在 $\Bun_{\mathbb{G}(0,2)}$ 上的作用施于平凡 $\mathbb{G}(0,2)$-丛，得到典范映射 $j_0:I(1)/I(2)\hookrightarrow\Bun_{\mathbb{G}(0,2)}$。它是仿射开嵌入。
\item 对任意 $\gamma\in\Omega$，映射
\[
j_\gamma:=\mathrm{Hk}_\gamma\circ j_0:I(1)/I(2)\hookrightarrow\Bun_{\mathbb{G}(0,2)}
\]
是仿射开嵌入，称为\emph{大胞腔}。
\item 把 $T_0^{\mathrm{red}}\times I(1)/I(2)$ 在 $\Bun_{\mathbb{G}(1,2)}$ 上的作用施于平凡 $\mathbb{G}(1,2)$-丛，得到典范仿射嵌入：
\[
\tilde j_0:T_0^{\mathrm{red}}\times I(1)/I(2)\hookrightarrow\Bun_{\mathbb{G}(1,2)},\qquad
\tilde j_\gamma:=\mathrm{Hk}_\gamma\circ\tilde j_0:T_0^{\mathrm{red}}\times I(1)/I(2)\hookrightarrow\Bun_{\mathbb{G}(1,2)}.
\]
\end{enumerate}
\end{corollary}

\begin{proof}
先回忆为何本推论对 $GL_n$ 成立。对 (1)，注意 $\PP^1$ 上唯一上同调平凡的秩 $n$ 向量丛是 $\mathcal O(-1)^n$。因此 $\Bun_{GL_n}$ 上上同调行列式线丛之逆有一个零点恰在平凡丛上的截面。这证明了本情形的 (1)。其次回忆 $GL_n(0,0)$-丛就是向量丛连同 $0$ 与 $\infty$ 处的完整旗。记 $\mathrm{Mod}_{i,0}:\Bun^0_{GL_n(0,0)}\to\Bun^i_{GL_n(0,0)}$ 为沿 $0$ 处旗的第 $i$ 次上修改。其逆由第 $i$ 次下修改给出，我们记为 $\mathrm{Mod}_{-i,0}$。为证 (2)，记 $\mathrm{Mod}_{i,\infty}$ 为 $\infty$ 处的第 $i$ 次修改。平凡 $GL_n(0,0)$-丛是由 $\mathcal O^n$ 以及 $0$ 处纤维 $(\mathcal O^n)_0$ 的典范对偶旗 $(V_{i,0})_{i=0,\dots,n}$ 与 $\infty$ 处的 $(V_{i,\infty})_{i=0,\dots,n}$ 给出的丛。它是 $\PP^1$ 上唯一的、使复形 $E\to E_0/V_{i,0}\oplus E_\infty/V_{n-i,\infty}$ 对所有 $i$ 都上同调平凡的次数 $0$ 的 $GL_n(0,0)$-丛 $(E,V_{i,0},V_{i,\infty})$。因此平凡 $GL_n(0,0)$-丛在 $\Bun_{GL_n(0,0)}$ 中的包含又一次由线丛截面的非消没定义。

对一般 $G$，取忠实表示 $\rho:\mathbb{G}\to GL_n\times\PP^1$。它定义映射 $\Bun_G\to\Bun_{GL_n}$。Birkhoff--Grothendieck 分解蕴涵：$G$-丛平凡当且仅当相伴的 $GL_n$-丛平凡。此外，为验证 $0$、$\infty$ 处的约化到 $B$ 是互相相反的，也只需在诱导的 $GL_n$-丛上验证。这证明了 (1) 与 (2)。

(3) 由 (2) 以及命题 1.1 证明中的同构 $H^0(\PP^1,T)\cong T_0^{\mathrm{red}}$ 推出。(4) 由 (3) 推出，因为映射 $\Bun_{\mathbb{G}(0,2)}\to\Bun_{\mathbb{G}(0,1)}$ 是 $I(1)/I(2)$-torsor。由推论 1.2，(5) 由 (4) 推出。最后，(6) 由 (5) 推出，因为 $\Bun_{\mathbb{G}(1,2)}\to\Bun_{\mathbb{G}(0,2)}$ 是 $T_0^{\mathrm{red}}$-torsor。
\end{proof}

% ============================================================
%  ch03 —— §2.1–2.3（前半，原文 p.12–p.15）
% ============================================================
\routenote{p.12--p.15}

\section{特征层与特征值——主要结果的陈述}

本节将构造自守层 $A_\varphi$。我们还将陈述关于 Kloosterman 层 $\Kl^{L}_G(\varphi,\chi)$ 的局部与整体单值性的主要结果；这个层将被定义为 $A_\varphi$ 的特征值。

\subsection{特征函数}

本小节给出 Gross 的自守表示 $\pi$ 中特征函数的一个简单公式。令人惊讶的是，这个计算不依赖于 Gross 的结果。我们将使用 §1.3 中选定的仿射一般特征 $\psi\circ\varphi:I(1)/I(2)\to\mathbb{Q}_\ell$。

由命题 1.1 知
\begin{gather*}
\Bun_{G(0,2)}(k)=G(k(t))\backslash G(0,2)(\mathbb{A}_{k(t)})\Big/\prod_x G(0,2)(\mathcal O_x)
=I^-(0)\backslash G(k((s)))/I(2),\\
G(k((s)))=\bigsqcup_{w\in\widetilde W}I^-(0)\,w\,I(1).
\end{gather*}
假设如引言中那样给定 $G(\mathbb{A}_{k(t)})$ 的一个自守表示 $\pi=\otimes'\pi_\nu$，使得对 $\nu\notin\{0,\infty\}$ 局部表示 $\pi_\nu$ 非分歧，$\pi_0$ 是 Steinberg 表示，且 $\pi_\infty$ 出现在 $\mathrm{c\text{-}Ind}_{Z(G(0,1))\times I(1)}^{G(k((s)))}\psi\circ\varphi$ 中。则存在 $\Bun_{G(0,2)}(k)$ 上的函数 $f$，使得对所有 $k\in I(1)$ 有 $f(gk)=\psi(\varphi(k))f(g)$。下面的引理以初等方式刻画这类函数：

\begin{lemma}
设 $f:\Bun_{G(0,2)}(k)\to\mathbb{Q}_\ell$ 是函数，使得对所有 $k\in I(1)$、$g\in G(0,2)(k((s)))$ 有 $f(gk)=\psi(\varphi(k))f(g)$。则 $f$ 由它在 $\gamma\in\Omega$ 处的值 $f(\gamma)$ 唯一决定。此外，对所有 $w\in\widetilde W-\Omega$ 有 $f(w)=0$。
\end{lemma}

\begin{proof}
对任意 $w\in\widetilde W$ 且 $l(w)>0$，存在单仿射根 $\alpha_i$ 使 $w.\alpha_i$ 为负。这蕴涵 $wU_{\alpha_i}w^{-1}\subset I^-(0)$。于是对所有 $u\in U_{\alpha_i}$ 有
\[
\psi(\varphi(u))f(w)=f(wu)=f(w),
\]
故 $f(w)=0$。对 $l(w)=0$ 的 $w$，即 $w\in\Omega$，由推论 1.3 可知 $f(\gamma)$ 的值可任意。
\end{proof}

\begin{remark}
上述引理中的任意函数 $f$ 自动是尖的。事实上，对任意抛物子群 $Q\subset G$ 及其幂幺根 $N_Q$，常数项函数
\[
f_{N_Q}(g)=\int_{N_Q(K)\backslash N_Q(\mathbb{A}_K)}f(ng)\,dn
\]
在 $N_Q(K_\infty)$ 下左不变，且相对于特征 $\psi\varphi$ 在 $I_\infty(1)$ 下右等变。与引理 2.1 类似的论证表明这样的函数必为零。
\end{remark}

\subsection{特征层}

让我们把前面的观察几何化。记 $\mathrm{AS}_\psi$ 为由特征 $\psi$ 定义的 $\Aff^1$ 上的 Artin--Schreier 层。令 $\mathrm{AS}_\varphi:=\varphi^*(\mathrm{AS}_\psi)$，即 Artin--Schreier 层到 $I(1)/I(2)\cong\Aff^d$ 的拉回。

记 $\mathrm{Perv}(\Bun_{G(0,2)})^{I(1),\mathrm{AS}_\varphi}$ 为 $\Bun_{G(0,2)}$ 上 $(I(1),\mathrm{AS}_\varphi)$-等变的 perverse 层范畴。对带 $I(1)/I(2)$ 作用的任意叠，我们也将一般地使用这个记号。

对任意 $\gamma\in\Omega$，记 $j_\gamma:I(1)/I(2)\hookrightarrow\Bun^\gamma_{G(0,2)}$ 为大胞腔的嵌入，$i_\gamma:\operatorname{Spec}(k)\to\Bun^\gamma_{G(0,2)}$ 为由 $\gamma$ 定义的 $G(0,2)$-丛给出的映射。下面是引理 2.1 的几何类比。

\begin{lemma}
层 $\mathrm{AS}_\varphi$ 满足 $j_{\gamma,!}\mathrm{AS}_\varphi=j_{\gamma,*}\mathrm{AS}_\varphi$。此外，对任意 $\gamma\in\Omega$，函子
\[
\mathrm{Perv}(\GG_m)\to\mathrm{Perv}\big(\Bun^\gamma_{G(0,2)}\times\GG_m\big),\qquad
F\mapsto j_{\gamma,!}\mathrm{AS}_\varphi\boxtimes F[\dim(\Bun_{G(0,2)})]
\]
是范畴等价。其逆由 $K\mapsto(i_\gamma\times\mathrm{id}_{\PP^1\setminus\{0,\infty\}})^*K[-\dim(\Bun_{G(0,2)})]$ 给出。
\end{lemma}

\begin{proof}
对任意 $w\in\widetilde W-\Omega$，在 $N(T((t)))$ 中取代表元，仍记为 $w$。考虑由 $w$ 定义的 $G(0,2)$-丛 $P_w$。设 $U_\alpha\subset I(1)$ 是对应于某个单仿射根 $\alpha$ 的根子群，使得 $w.\alpha$ 为负，即 $U_{w.\alpha}\subset I^-(0)$。这定义包含 $U_\alpha\hookrightarrow\operatorname{Aut}(P_w)$。于是得到交换图
\[
\begin{tikzcd}
U_\alpha\times\operatorname{Spec}(k) \arrow[r, "{(\mathrm{id},P_w)}"] \arrow[d] & \operatorname{Spec}(k) \arrow[d, "{P_w}"]\\
U_\alpha\times\Bun_{G(0,2)} \arrow[r, "{\text{作用}}"] & \Bun_{G(0,2)}.
\end{tikzcd}
\]
这蕴涵 $(\mathrm{AS}_\varphi|_{U_\alpha})\boxtimes K|_{P_w}\cong\mathbb{Q}_\ell|_{U_\alpha}\boxtimes K|_{P_w}$。由于我们假设 $\mathrm{AS}_\varphi|_{U_\alpha}$ 由 $U_\alpha$ 的非平凡特征定义，可知 $K$ 在 $P_w$ 处的茎为零。对偶地，同样的结果对 $K$ 在 $P_w$ 处的余茎成立。

这证明了第一个断言。对第二个断言，它还蕴涵：$\Bun^\gamma_{G(0,2)}$ 上任意 $(I(1),\mathrm{AS}_\varphi)$-等变 perverse 层都是其从子叠 $j_\gamma(I(1)/I(2))$ 的 $!$-扩张。在这个子叠上，与局部系统 $\mathrm{AS}_\varphi$ 作张量给出 $(I(1)/I(2))$-等变层与 $(I(1)/I(2),\mathrm{AS}_\varphi)$-等变层之间的等价。这证明了我们的断言。
\end{proof}

利用这个引理，我们现在可以定义自守层：

\begin{definition}
定义 $A_\varphi\in\mathrm{Perv}(\Bun_{G(0,2)})^{I(1),\mathrm{AS}_\varphi}$ 为在分支 $\Bun^\gamma_{G(0,2)}$ 上由 $j_{\gamma,!}\mathrm{AS}_\varphi[\dim(\Bun_{G(0,2)})]$ 给出的 perverse 层。记 $A^\gamma_\varphi$ 为 $A_\varphi$ 到分支 $\Bun^\gamma_{G(0,2)}$ 的限制。
\end{definition}

\begin{remark}[带乘性特征的变体]
我们可以把 $A_\varphi$ 的上述构造稍作推广。回忆 $\Bun^\gamma_{G(1,2)}$ 中的开胞腔典范同构于 $T_0^{\mathrm{red}}\times I(1)/I(2)$\zhnote{原文此处写作 $I(1)/I(0)$；由推论 1.3(6) 与 $I(1)/I(2)\cong\Aff^n$ 的记号，应为 $I(1)/I(2)$，属原文笔误。以下同。[来源：原文 p.13 与推论 1.3(6)]}。任意特征 $\chi:T_0^{\mathrm{red}}(k)\to\mathbb{Q}_\ell$ 在 $T_0^{\mathrm{red}}$ 上定义特征层 $\mathrm{Kum}_\chi$。于是上述引理 2.3 对 $\Bun_{G(1,2)}$ 上 $(T_0^{\mathrm{red}}\times I(1)/I(2),\mathrm{Kum}_\chi\boxtimes\mathrm{AS}_\varphi)$-等变层也成立。记 $\Bun_{G(1,2)}$ 上相应的 perverse 层为 $A_{\varphi,\chi}$。
\end{remark}

\subsection{几何 Hecke 算子}

为陈述主要结果，我们需要回忆几何 Hecke 算子的定义。

Hecke 修改之叠是
\[
\begin{aligned}
\Hecke_{\mathbb{G}(m,n)}^{\Aff^1}(S):=\Big\{(E_1,E_2,x,\phi)\ \Big|\ &E_i\in\Bun_{\mathbb{G}(m,n)}(S),\ x:S\to\Aff^1,\\
&\phi:E_1|_{\PP^1-x\times S}\xrightarrow{\sim}E_2|_{\PP^1-x\times S}\Big\}.
\end{aligned}
\]
这个叠有自然的遗忘映射：
\begin{equation}
\begin{tikzcd}[column sep=large]
& \Hecke^{\Aff^1}_{\mathbb{G}(m,n)} \arrow[dl, "\mathrm{pr}_1"'] \arrow[dr, "\mathrm{pr}_2"] \arrow[d, "\mathrm{pr}_{\Aff^1}"] & \\
\Bun_{\mathbb{G}(m,n)} & \Bun_{\mathbb{G}(m,n)}\times\Aff^1 \arrow[l, "\mathrm{proj}_1"'] \arrow[r] & \Aff^1
\end{tikzcd}
\end{equation}
其中 $\mathrm{pr}_1(E_1,E_2,x,\phi)=E_1$，$\mathrm{pr}_2(E_1,E_2,x,\phi)=(E_2,x)$，$\mathrm{pr}_{\Aff^1}(E_1,E_2,x,\phi)=x$。记
\[
\Hecke_{\mathbb{G}(m,n)}^{\PP^1\setminus\{0,\infty\}}:=\mathrm{pr}_{\Aff^1}^{-1}\big(\PP^1\setminus\{0,\infty\}\big).
\]

\begin{remark}\label{rem:2.7}
\begin{enumerate}[label=(\arabic*), leftmargin=1.8em, itemsep=.2em]
\item $\mathrm{pr}_1$ 在平凡丛 $\mathbb{G}(m,n)\in\Bun_G$ 上的纤维称为 Beilinson--Drinfeld Grassmannian，记作 $\mathrm{GR}_{\mathbb{G}(m,n)}$。$\mathrm{GR}_{\mathbb{G}}\to\Aff^1$ 在点 $x\in\PP^1\setminus\{0,\infty\}$ 上的纤维同构于仿射 Grassmannian $\mathrm{Gr}_{\mathbb{G},x}$，即商 $\mathbb{G}(K_x)/\mathbb{G}(\mathcal O_x)$。
\item $\mathrm{pr}_2$ 在 $\Bun_{\mathbb{G}(m,n)}\times\PP^1\setminus\{0,\infty\}$ 上的几何纤维（非典范地）同构于仿射 Grassmannian $\mathrm{Gr}_{\mathbb{G}}$。在 $\Bun_{\mathbb{G}}\times\PP^1\setminus\{0,\infty\}$ 的光滑拓扑下局部，这个纤维化是平凡的（例如注 4.1）。
\item 图 (2.6) 有一大族对称性。群 $I(0)/I(2)=(I(1)/I(2))\rtimes T^{\mathrm{red}}_\infty$ 通过改变 $\infty$ 处的 $I(2)$-水平结构作用在 $\Bun_{\mathbb{G}(m,2)}$ 上，且这个作用扩充到图 (2.6)（即它也作用在 $\Hecke^{\Aff^1}_{\mathbb{G}(m,2)}$ 上，且映射 $\mathrm{pr}_i$ 在这些作用下等变）。一维环面 $\GG_m^{\mathrm{rot}}$ 作用在固定点 $\{0\}$ 与 $\{\infty\}$ 的曲线 $\PP^1$ 上，故也作用在 (2.6) 上。最后，钉住自同构 $\operatorname{Aut}^\dagger(\mathbb{G})$ 作用在 (2.6) 上。于是我们看到群
\[
I(0)/I(2)\rtimes\big(\GG_m^{\mathrm{rot}}\times\operatorname{Aut}^\dagger(\mathbb{G})\big)
\]
作用在图 (2.6) 上。
\end{enumerate}
\end{remark}

先回忆常群概形的 Hecke 算子。为此我们收集几何 Satake 等价的一些事实（见 [30]、[18]）。

设 $\mathrm{Gr}_G=G((\tau))/G[[\tau]]$ 为抽象的仿射 Grassmannian，不参照 $\PP^1\setminus\{0,\infty\}$ 上的任何点。令 $\mathcal O=k[[\tau]]$，$\operatorname{Aut}_{\mathcal O}$ 为 $\mathcal O$ 的连续（在 $\tau$-进拓扑下）自同构的 pro-代数 $k$-群。则 $G[[\tau]]\rtimes\operatorname{Aut}_{\mathcal O}$ 从左侧作用在 $\mathrm{Gr}_G$ 上。$G[[\tau]]$ 在 $\mathrm{Gr}_G$ 上的轨道由支配余特征 $\mu\in X_*(T)^+$ 编号。轨道记为 $\mathrm{Gr}_{G,\mu}$，其闭包（Schubert 簇）记为 $\mathrm{Gr}_{G,\le\mu}$。记 $\mathrm{Gr}_{G,\le\mu}$ 的相交上同调层为 $\mathrm{IC}_\mu$。我们将把 $\mathrm{IC}_\mu$ 归一化为权 $0$（见下面的注 2.9）。

Satake 范畴 $\mathrm{Sat}=\mathrm{Perv}_{\operatorname{Aut}_{\mathcal O}}(G[[\tau]]\backslash G((\tau))/G[[\tau]])$ 是 $\mathrm{Gr}_G$ 上 $G[[\tau]]\rtimes\operatorname{Aut}_{\mathcal O}$-等变 perverse 层（带有限型支撑）的范畴。类似地，把情形基变换到 $\bar k$ 定义 $\mathrm{Sat}^{\mathrm{geo}}$。最后把规范化的半单 Satake 范畴 $\mathcal S$ 定义为 $\mathrm{Sat}$ 中由 $\mathrm{IC}_\mu$ 的直和组成的满子范畴。

[30] 与 [18] 证明了：
\begin{itemize}[leftmargin=1.6em, itemsep=.2em]
\item $\mathrm{Sat}^{\mathrm{geo}}$ 带有自然的张量结构（它对 $\mathrm{Sat}$ 也有定义），使得整体上同调函子 $h=H^*(\mathrm{Gr}_G,-):\mathrm{Sat}^{\mathrm{geo}}\to\mathrm{Vec}$ 是纤维函子。
\item $\operatorname{Aut}^\otimes(h)$ 是 $\mathbb{Q}_\ell$ 上连通约化群，它 Langlands 对偶于 $G$。令 $\ckG:=\operatorname{Aut}^\otimes(h)$，则 Tannaka 形式化给出张量范畴的几何 Satake 等价
\[
\mathrm{Sat}^{\mathrm{geo}}\cong\Rep(\ckG).
\]
\item 按构造，$\ckG$ 配有极大环面 $\ckT$ 与自然同构 $X^*(\ckT)\cong X_*(T)$（事实上 $\ckG$ 配有典范钉住；见引理 9.3）。
\end{itemize}

几何 Satake 等价把 $\mathrm{IC}_\mu$ 送到极值权为 $\mu$ 的不可约表示 $V_\mu$。记这个等价之逆为 $V\mapsto\mathrm{IC}_V$。

\begin{remark}
在 [1, §3.5] 中论证了 $\mathcal S$ 在 $\mathrm{Sat}$ 的张量结构下封闭。因此 $\mathcal S$ 自然是张量范畴。

沿 $\mathrm{Gr}_G\otimes_k\bar k\to\mathrm{Gr}_G$ 的拉回给出张量函子 $\mathcal S\to\mathrm{Sat}^{\mathrm{geo}}$；由于两个范畴都是半单的且单对象显式已知，容易看出它是等价。因此 [30] 与 [18] 中的上述结果都适用于 $\mathcal S$。特别地，我们得到（$k$-版本的半单化）几何 Satake 等价
\[
\mathcal S\cong\mathrm{Sat}^{\mathrm{geo}}\cong\Rep(\ckG).
\]
\end{remark}

\begin{remark}[权的归一化]
我们采用使该复形纯权 $0$ 的归一化，即在 $\mathbb{Q}_\ell$ 中取 $q$ 的一个平方根，并把 $\mathrm{IC}_\mu$ 记为相交复形张量上 $\mathbb{Q}_\ell\big(\tfrac12\dim(\mathrm{Gr}_\mu)\big)$。

正如 [15] 所指出的，并不必要作这个相当不自然的选择——作它是为了从范畴 $\mathcal S$ 得到群 $\ckG$。替代地，我们可以把范畴 $\mathcal S$ 扩大，纳入相交上同调层 $\mathrm{IC}_\mu(n)$ 的所有 Tate 扭转。由前一注，这仍是一个中性张量范畴，定义群 $\ckG_1$，它是 $\ckG$ 被一个中心一维环面所作的扩张。
\end{remark}

% ============================================================
%  ch04 —— §2.3（后半）–2.5（原文 p.16–p.18）
% ============================================================
\routenote{p.16--p.18}

几何 Hecke 算子定义为函子（见 [17]）：
\begin{align*}
\Hk:\ \Rep(\ckG)\times D^b(\Bun_{\mathbb{G}(m,n)})&\longrightarrow D^b\big(\Bun_{\mathbb{G}(m,n)}\times\PP^1\setminus\{0,\infty\}\big),\\
(V,K)&\longmapsto \Hk_V(K):=\mathrm{pr}_{2,!}\mathrm{pr}_1^*(K\otimes\mathrm{IC}_V).
\end{align*}
为了复合这些算子，把 $\Hk_V$ 扩充为算子 $D^b(\Bun_{\mathbb{G}(m,n)}\times\PP^1\setminus\{0,\infty\})\to D^b(\Bun_{\mathbb{G}(m,n)}\times\PP^1\setminus\{0,\infty\})$：
\[
K\longmapsto \mathrm{pr}_{2,!}\Big(\big(\mathrm{pr}_1\times\mathrm{pr}_{\PP^1\setminus\{0,\infty\}}\big)^*K\otimes\mathrm{IC}_V\Big).
\]

设 $E$ 是 $\PP^1\setminus\{0,\infty\}$ 上的 $\ckG$-局部系统，视作张量函子
\[
E:\ \mathcal S\cong\Rep(\ckG)\to\Loc(\PP^1\setminus\{0,\infty\}),
\]
把 $V\in\Rep(\ckG)$ 送到 $E_V\in\Loc(\PP^1\setminus\{0,\infty\})$。

以 $E$ 为特征值的 Hecke 特征层是 perverse 层 $K\in\mathrm{Perv}(\Bun_{\mathbb{G}(m,n)})$ 连同同构
\[
\Hk_V(K)\xrightarrow{\ \sim\ }K\boxtimes E_V,
\]
这些同构与 $\Rep(\ckG)$ 上的对称张量结构以及 Hecke 对应的复合相容（关于相容性公理的详细阐述见 [17]）。

\begin{remark}
由于局部系统通常是以别的方式引入的，让我们简要回忆张量函子 $E$ 如何重构局部系统的单值表示。这对于固定单值表示的记号是有用的。取 $\PP^1\setminus\{0,\infty\}$ 中 $\operatorname{Spec}K_0$ 上方的一个几何点 $\eta$。限制到 $\eta$ 定义张量函子
\begin{equation}
\omega_E:\ \mathcal S\xrightarrow{\ E\ }\Loc(\PP^1\setminus\{0,\infty\})\xrightarrow{\ j_\eta^*\ }\mathrm{Vec},
\end{equation}
亦即（见 [12, Theorem 3.2]）一个带 $\pi_1(\PP^1\setminus\{0,\infty\},\eta)$ 作用的 $\ckG$-挠子。取该挠子的一个点，便定义了单值表示
\begin{equation}
\phi:\ \pi_1\big(\PP^1\setminus\{0,\infty\},\eta\big)\to\ckG(\mathbb{Q}_\ell).
\end{equation}
记 $\phi^{\mathrm{geo}}$ 为 $\phi$ 到 $\pi_1^{\mathrm{geo}}(\PP^1\setminus\{0,\infty\},\eta)=\pi_1(\PP^1\setminus\{0,\infty\}\otimes_k\bar k,\eta)$ 的限制，称之为 $E$ 的整体几何单值表示。
\end{remark}

上述构造对扭群 $\mathbb{G}$ 的类比将产生 $\LG$-局部系统。这里把 $\mathbb{G}$ 的 $L$-群定义为 $\LG=\ckG\rtimes\langle\sigma\rangle$，其中 $\sigma$ 是 $\mathbb{G}$ 的定义中所用的 $G$ 的 $N$ 阶自同构。这个自同构通过几何 Satake 同构诱导 $\ckG$ 的自同构。在引理 9.3 中我们将验证这个自同构确实保持 $\ckG$ 的钉住。

让我们给出对本文目的足够的 $\LG$-局部系统的定义。我们不能立即应用几何 Satake 同构的原因是 $\mathrm{pr}_2$ 的纤维沿 $\PP^1\setminus\{0,\infty\}$ 不是常值的，故只有部分 Hecke 算子 $\Hk_V$ 会定义 $\PP^1\setminus\{0,\infty\}$ 上的整体 Hecke 算子。不过，我们可以沿映射 $[N]:\PP^1\setminus\{0,\infty\}\to\PP^1\setminus\{0,\infty\}$ 拉回卷积图。把这个覆盖记作 $\widetilde{\PP^1\setminus\{0,\infty\}}\to\PP^1\setminus\{0,\infty\}$。

拉回之后，我们可以像以前一样在 $\Hecke^{\widetilde{\PP^1\setminus\{0,\infty\}}}_{\mathbb{G}(m,n)}$ 上定义 Hecke 算子
\[
\Hk:\ \Rep(\ckG)\times D^b(\Bun_{\mathbb{G}(m,n)})\cong\mathcal S\times D^b(\Bun_{\mathbb{G}(m,n)})\to D^b\big(\Bun_{\mathbb{G}(m,n)}\times\widetilde{\PP^1\setminus\{0,\infty\}}\big).
\]
此外，覆盖群 $\mu_N$ 作用在 $\widetilde{\PP^1\setminus\{0,\infty\}}$ 上的卷积图上。这在函子 $\Hk$ 上定义了 $\mu_N$-等变结构：在源上，$\mu_N$ 通过 $\sigma:\mu_N\to\operatorname{Aut}^\dagger(\mathbb{G})$ 作用在 $\mathcal S$ 上，而后者可等同于 $\operatorname{Aut}^\dagger(\ckG)$ 在 $\Rep(\ckG)$ 上的作用（引理 9.3）；在靶上，$\mu_N$ 作用在 $\widetilde{\PP^1\setminus\{0,\infty\}}$ 上。

设 $E$ 是 $\widetilde{\PP^1\setminus\{0,\infty\}}$ 上的 $\ckG$-局部系统，连同对 $\zeta\in\mu_N$ 的相容同构 $\zeta^*E\cong E\times^{\ckG,\sigma(\zeta)}\ckG$。我们把 $E$ 视作带 $\mu_N$-等变结构的张量函子
\[
E:\ \Rep(\ckG)\to\Loc(\widetilde{\PP^1\setminus\{0,\infty\}}).
\]
和前一样，可以定义以 $E$ 为特征值的 Hecke 特征层 $K\in\mathrm{Perv}(\Bun_{\mathbb{G}(m,n)})$，但现在须指定与 $\Rep(\ckG)$ 上张量结构相容的函子同构
\[
\epsilon(V):\ \Hk_V(K)\xrightarrow{\ \sim\ }K\boxtimes E_V,
\]
它与函子 $V\mapsto\Hk_V(K)$ 和函子 $E$ 的 $\mu_N$-等变结构交换。

注意，若 $\sigma:\mu_N\to\operatorname{Aut}(\mathbb{G})$ 平凡，则同构 $\zeta^*E\cong E\times^{\ckG,\sigma(\zeta)}\ckG$ 给出 $E$ 的下降数据。故此时定义与常群情形一致。

有了这些定义，我们可以陈述第一个主要结果：

\begin{mainthm}[1]
\begin{enumerate}[label=(\arabic*), leftmargin=1.9em, itemsep=.25em]
\item 层 $A=A_\varphi$ 与 $A_{\varphi,\chi}$ 是 Hecke 特征层。记 $A_\varphi$（相应地 $A_{\varphi,\chi}$）的特征值为 $\Kl^{\LG}(\varphi)$（相应地 $\Kl^{\LG}(\varphi,\chi)$）。
\item 若 $\mathbb{G}=\PP^1\times G$ 是常群概形，则局部系统 $\Kl^{\ckG}(\varphi)$ 在 $0$ 处驯顺分歧。$\Kl^{\ckG}(\varphi)$ 在 $0$ 处的单值作用由 $\ckG$ 中的一个主幂幺元给出。
\item 对任意不可约表示 $V\in\Rep(\ckG)$，层 $\Kl^{\LG}(\varphi,\chi)^V$ 是纯的。
\end{enumerate}
\end{mainthm}

由于我们用几何 Hecke 算子定义了 $\Kl^{\ckG}(\varphi)$，对任意点 $x\in\PP^1\setminus\{0,\infty\}$，由该局部系统定义的 $\mathrm{Frob}_x$ 的 $\ckG$-共轭类由 Gross 的自守形式 $\pi(\varphi)$ 在 $x$ 处的 Satake 参数给出。

\subsection{单值表示}

本节假设 $\mathbb{G}=G\times\PP^1$，其中 $G$ 是 $k$ 上几乎单的分裂群。如注 2.10 那样，记 $\phi$ 为 $\Kl^{\ckG}(\varphi)$ 的单值表示，$\phi^{\mathrm{geo}}$ 为其几何单值表示。

下一个结果是关于 $\Kl^{\ckG}(\varphi,\chi)$ 在 $\infty$ 处的几何单值的。记 $\Kl^{\mathrm{Ad}}_{\ckG}(\varphi,\chi)$ 为由伴随表示 $V=\ckg$ 诱导的 $\PP^1\setminus\{0,\infty\}$ 上的局部系统。固定局部 Galois 群 $\operatorname{Gal}(K^{\mathrm{sep}}_\infty/K_\infty)$ 到 $\pi_1(\PP^1\setminus\{0,\infty\},\eta)$ 的一个嵌入，便得到 $\operatorname{Gal}(K^{\mathrm{sep}}_\infty/K_\infty)$ 在 $\Kl^{\mathrm{Ad}}_{\ckG}(\varphi,\chi)$ 于形式去心圆盘 $\operatorname{Spec}K^{\mathrm{sep}}_\infty$ 处的几何茎上的作用。把这个茎与 $\ckg$ 等同，便得到 $\operatorname{Gal}(K^{\mathrm{sep}}_\infty/K_\infty)$ 在 $\ckg$ 上的作用（在 $\ckG$-共轭意义下良定义）。

设 $I_\infty\subset\operatorname{Gal}(K^{\mathrm{sep}}_\infty/K_\infty)$ 为惯性群，$I^+_\infty\subset I_\infty$ 为野惯性群，$I^t_\infty=I_\infty/I^+_\infty$ 为驯惯性群。

\begin{mainthm}[2]
设 $p=\operatorname{char}(k)$ 对 $G$ 良好（即当 $G$ 是 $B_n$、$C_n$ 型时 $p>2$，当 $G$ 是 $F_4$、$G_2$ 型时 $p>3$）。则
\begin{itemize}[leftmargin=1.6em, itemsep=.2em]
\item $\Swan_\infty\big(\Kl^{\mathrm{Ad}}_{\ckG}(\varphi,\chi)\big)=r(\ckG)$，即 $\ckG$ 的秩；
\item $\ckg^{I_\infty}=0$。
\end{itemize}
\end{mainthm}

Swan 等式将在推论 5.1 中证明；$I_\infty$-不变量的消没将在命题 5.3(2) 中证明。

\begin{corollary}
（Gross--Reeder [20, Proposition 5.6]）设 $p=\operatorname{char}(k)$ 不整除 $\#W$。则局部几何 Galois 表示 $\phi^{\mathrm{geo}}_\infty:I_\infty\to\ckG$ 是 [20, §6] 中定义的简单野参数。更确切地说，在 $\ckG$-共轭意义下，有正合列的交换图
\[
\begin{tikzcd}[column sep=large, row sep=small]
I^+_\infty \arrow[r] \arrow[d, "\phi^{\mathrm{geo},+}_\infty"'] & I_\infty \arrow[r] \arrow[d, "\phi^{\mathrm{geo}}_\infty"] & I^t_\infty \arrow[d] \\
\ckT \arrow[r] & N(\ckT) \arrow[r] & W
\end{tikzcd}
\]
其中
\begin{itemize}[leftmargin=1.6em, itemsep=.2em]
\item 驯惯性群 $I^t_\infty$ 的一个拓扑生成元映到 Coxeter 元 $\mathrm{Cox}\in W$（在 $W$-共轭意义下良定义，见 [6, Ch.~V, §6]）；
\item 野惯性群 $I^+_\infty$ 满射到一个子群 $\ckT(\zeta)\subset\ckT[p]$。这里 $\zeta\in\overline{\mathbb{F}}_p^\times$ 是一个本原 $h$ 次单位根（$h$ 为 Coxeter 数），而 $\ckT(\zeta)\subset\ckT[p]$ 是 $\ckT[p]$ 中唯一的 $\mathbb{F}_p[\mathrm{Cox}]$-子模，它同构于 $\mathbb{F}_p[\zeta]$（$\mathrm{Cox}$ 在其上以 $\zeta$ 乘法作用）。
\item $I_\infty$-表示 $\Kl^{\mathrm{Ad}}_{\ckG}(\varphi,\chi)$ 的非零跳跃等于 $1/h$。
\end{itemize}
\end{corollary}

最后我们可以陈述关于 $\Kl^{\ckG}(\varphi)$ 的整体几何单值的结果。设 $\ckG^{\mathrm{geo}}\subset\ckG$ 为整体几何单值表示 $\phi^{\mathrm{geo}}$ 之像的 Zariski 闭包。

\begin{mainthm}[3]
设 $\operatorname{char}(k)>2$。则 $\Kl^{\ckG}(\varphi)$ 的几何单值群 $\ckG^{\mathrm{geo}}$ 连通，且
\begin{itemize}[leftmargin=1.6em, itemsep=.2em]
\item 若 $\ckG$ 不是 $A_{2n}$（$n\ge 2$）型或 $B_3$ 型，则 $\ckG^{\mathrm{geo}}=\ckG^{\operatorname{Aut}^\dagger(\ckG),\circ}$；
\item 若 $\ckG$ 是 $A_{2n}$ 型，则 $\ckG^{\mathrm{geo}}=\ckG$；
\item 若 $\ckG$ 是 $B_3$ 型且 $\operatorname{char}(k)>3$，则 $\ckG^{\mathrm{geo}}=G_2$。
\end{itemize}
\end{mainthm}

证明将在 §6.2 给出。

\subsection{变体}

我们的构造有一个用 $D$-模代替 $\ell$-进层的变体。此时基域取为 $k=\mathbb{C}$。Artin--Schreier 局部系统 $\mathrm{AS}_\psi$ 被 $\Aff^1_\mathbb{C}=\operatorname{Spec}\mathbb{C}[x]$ 上的指数 $D$-模 $\mathbb{C}\langle x,\partial_x\rangle/(\partial_x-1)$ 替代，构造的其余部分照搬，于是得到张量函子
\[
\Kl^{\LG}(\varphi)_{dR}:\ \Rep({}^L G_\mathbb{C})\to\mathrm{Conn}(\PP^1\setminus\{0,\infty\},\mathbb{C}),
\]
其中 $\mathrm{Conn}(\PP^1\setminus\{0,\infty\},\mathbb{C})$ 是 $\PP^1\setminus\{0,\infty\},\mathbb{C}$ 上带联络的向量丛的张量范畴。

我们猜想本文的构造应给出与 Frenkel--Gross 构造相同的联络。为精确陈述，设 $G$ 几乎单，$\mathbb{G}$ 是 $G$ 在 $\PP^1\setminus\{0,\infty\}$ 上的拟分裂形式，由 $\sigma:\mu_N\to\operatorname{Aut}^\dagger(G)$ 给出。回忆 [15, §5] 中 $\PP^1\setminus\{0,\infty\},\mathbb{C}$ 上平凡 $\LG$-丛上的 Frenkel--Gross 联络：
\[
\nabla^{\LG}(\check X_0,\cdots,\check X_{r_\sigma})=d+\sum_{i=0}^{r_\sigma}\check X_i\,\frac{d}{dz},
\]
其中 $r_\sigma$ 是 $\mathbb{G}^\sigma$ 的秩，$\check X_i$ 是与 $\ckg$ 及 $\sigma$ 相伴的（扭）仿射 Kac--Moody 李代数的 $(-\alpha_i)$-根空间的一组基。

\begin{conjecture}
存在一般线性函数 $\varphi:I(1)/I(2)\to\mathbb{G}_{a,\mathbb{C}}$ 之集与基 $(\check X_0,\cdots,\check X_{r_\sigma})$ 之集之间的双射，使得只要 $\varphi$ 在这个双射下对应于 $(\check X_0,\cdots,\check X_{r_\sigma})$，就有 $\PP^1\setminus\{0,\infty\},\mathbb{C}$ 上 $\LG$-联络之间的自然同构
\[
\Kl^{\LG}(\varphi)_{dR}\cong({}^L G,\nabla^{\LG}(\check X_0,\cdots,\check X_{r_\sigma})).
\]
\end{conjecture}

% ============================================================
%  ch05 —— §3 GL_n 的例（原文 p.19–p.22）
% ============================================================
\routenote{p.19--p.22}

\section{例：$GL_n$ 的 Kloosterman 层}

本节对 $\PP^1$ 上的常群 $G=GL_n$ 计算 Kloosterman 层 $\Kl_{GL_n}(\varphi,\chi)$。它的 Langlands 对偶是 $GL_{n,\mathbb{Q}_\ell}$，记标准表示为 $\mathrm{Std}$。于是描述 $\PP^1\setminus\{0,\infty\}$ 上的这个 $GL_n$-局部系统，就等同于描述秩 $n$ 的局部系统 $\Kl^{\mathrm{Std}}_{GL_n}(\varphi,\chi)$。我们将看到这个秩 $n$ 局部系统与 Deligne 在 [10] 中定义的经典 Kloosterman 层一致。

\subsection{另一种模解释}

我们想用向量丛来解释 $GL_n(1,2)$-丛。先定义 $\Bun_{GL_n(1,2)}$ 的一个变体。设 $\mathrm{Bun}_{n,1,2}$ 为分类如下数据的叠 $(E,F^*E,\{v_i\},F_*E,\{v_i\})$：
\begin{enumerate}[label=(\arabic*), leftmargin=1.9em, itemsep=.15em]
\item $E$ 是 $\PP^1$ 上秩 $n$ 的向量丛；
\item 一个递减滤过 $F^*E$，给出 $E$ 在 $0$ 处纤维的完全旗：
\[
E=F^0E\supset F^1E\supset\cdots\supset F^nE=E(-\{0\});
\]
\item 对每个 $i=1,\cdots,n$，一个非零向量 $v_i\in F^{i-1}E/F^iE$；
\item 一个递增滤过 $F_*E$，给出 $E$ 在 $\infty$ 处纤维的完全旗：
\[
E(-\{\infty\})=F_0E\subset F_1E\subset\cdots\subset F_nE=E;
\]
\item 对 $i=1,\cdots,n$，一个向量 $v_i\in F_iE/F_{i-2}E$，它不落在 $F_{i-1}E/F_{i-2}E$ 中（这里约定 $F_{-1}E=(F_{n-1}E)(-\{\infty\})$）。
\end{enumerate}
注意 $\mathrm{Bun}_{n,1,2}$ 是 $\PP^1$ 上 $G$-挠子的模叠，其中 $\mathbb{G}$ 是 $\PP^1$ 上的 Bruhat--Tits 群概形，满足
\begin{itemize}[leftmargin=1.6em, itemsep=.15em]
\item $\mathbb{G}|_{\PP^1\setminus\{0,\infty\}}=GL_n\times(\PP^1\setminus\{0,\infty\})$；
\item $\mathbb{G}(\mathcal O_0)=I_{GL_n}(1)^{\mathrm{opp}}$；
\item $\mathbb{G}(\mathcal O_\infty)=Z_{GL_n}(1)(k[[s]])\cdot I_{SL_n}(2)\supset I_{GL_n}(2)$。
\end{itemize}
$\mathbb{G}$ 与 $GL_n(1,2)$ 的唯一差别在于：它们在 $\infty$ 处对中心 $\mathbb{G}_m=Z_{GL_n}$ 取不同的水平结构。因此有自然态射 $GL_n(1,2)\to\mathbb{G}$，从而有自然态射 $\Bun_{GL_n(1,2)}\to\mathrm{Bun}_{n,1,2}$，它是 $\mathbb{G}_a$-挠子。

选取 $E$ 在 $\PP^1\setminus\{0,\infty\}$ 上的平凡化，可以把 $\mathrm{Bun}_{n,1,2}$ 的模问题改写如下。设 $\Lambda$ 是以 $\{e_1,\cdots,e_n\}$ 为基的自由 $k[t,t^{-1}]$-模。令 $e_{i+jn}=t^je_i$（$1\le i\le n$，$j\in\mathbb{Z}$），则 $\Lambda$ 是以 $\{e_i\}_{i\in\mathbb{Z}}$ 为基的 $k$-向量空间。对任意 $k$-代数 $R$，$R\otimes\Lambda$ 中的 $R[t]$-格是指 $R[t]$-子模 $\Lambda'\subset R\otimes\Lambda$，使得存在 $M\in\mathbb{Z}_{>0}$ 满足
\[
\operatorname{Span}_R\{e_i\mid i>M\}\subset\Lambda'\subset\operatorname{Span}_R\{e_i\mid i\ge -M\},
\]
且 $\Lambda'/\operatorname{Span}_R\{e_i\mid i>M\}$ 与 $\operatorname{Span}_R\{e_i\mid i\ge-M\}/\Lambda'$ 都是投射 $R$-模。类似地可定义 $R\otimes\Lambda$ 中的 $R[t^{-1}]$-格。

设 $\widehat{\mathrm{Bun}}_{n,1,2}(R)$ 分类数据 $(\Lambda_*,\{\mathbf v_i\}_{0\le i\le n-1},\Lambda^*,\{\mathbf v_i\}_{1\le i\le n})$：
\begin{enumerate}[label=(\arabic*), leftmargin=1.9em, itemsep=.15em]
\item $R\otimes\Lambda\supset\Lambda_0\supset\Lambda_1\supset\cdots\supset\Lambda_n=t\Lambda_0$ 是 $R[t]$-格链，使得 $\Lambda_i/\Lambda_{i+1}$ 是秩 $1$ 的投射 $R$-模。对任意 $1\le i\le n$、$j\in\mathbb{Z}$，令 $\Lambda_{i+jn}=t^j\Lambda_i$。
\item $\mathbf v_i\in\Lambda_i/\Lambda_{i+1}$ 是 $R$-基，$i=0,\cdots,n-1$；
\item $R\otimes\Lambda\supset\Lambda^n\supset\Lambda^{n-1}\supset\cdots\supset\Lambda^0=t^{-1}\Lambda^n$ 是 $R[t^{-1}]$-格链，使得 $\Lambda^i/\Lambda^{i+1}$ 是秩 $1$ 的投射 $R$-模。对任意 $1\le i\le n$、$j\in\mathbb{Z}$，令 $\Lambda^{i+jn}=t^j\Lambda^i$。
\item $\mathbf v_i\in\Lambda^i/\Lambda^{i-2}$，它在 $\Lambda^i/\Lambda^{i-1}$ 中的像是 $R$-基，$i=1,\cdots,n$。
\end{enumerate}
群 $GL_n(k[t,t^{-1}])$（视作 $k$ 上的 ind-群概形）作用在 $\Lambda$ 上，从而作用在叠 $\widehat{\mathrm{Bun}}_{n,1,2}$ 上。此外，$\mathrm{Bun}_{n,1,2}$ 自然同构于商叠 $\widehat{\mathrm{Bun}}_{n,1,2}/GL_n(k[t,t^{-1}])$。

与格链 $(\Lambda_*,\Lambda^*)$ 相伴的是 $\operatorname{Spec}R$ 上局部常值的整值函数
\[
\deg(\Lambda_*,\Lambda^*):=\chi_R\Big(\Lambda_0\oplus\Lambda^0\xrightarrow{\ (\iota_0,\iota^0)\ }R\otimes\Lambda\Big),
\]
其中 $\iota_0,\iota^0$ 是包含映射，右端的 Euler 特征定义为 $\operatorname{rk}_R\ker(\iota_0,\iota^0)-\operatorname{rk}_R\operatorname{coker}_R(\iota_0,\iota^0)$。

对 $d\in\mathbb{Z}$，设 $\mathrm{Bun}^d_{n,1,2}\subset\mathrm{Bun}_{n,1,2}$ 为分类满足 $\deg(\Lambda_*,\Lambda^*)=d$ 的 $\{\Lambda_*,\mathbf v_i,\Lambda^*,\mathbf v_i\}$ 的子叠。

大胞腔的嵌入 $j_d:T\times\mathbb{G}_a^n\hookrightarrow\mathrm{Bun}^d_{n,1,2}$ 可如下固定。对 $a=(a_1,\cdots,a_n)\in T(R)=(R^\times)^n$ 与 $b=(b_1,\cdots,b_n)\in R^n$，点 $j_d(a,b)\in\mathrm{Bun}^d_{n,1,2}(R)$ 由如下数据的 $GL_n(R[t,t^{-1}])$-轨道给出：
\begin{enumerate}[label=(\arabic*), leftmargin=1.9em, itemsep=.15em]
\item $\Lambda_i=\operatorname{Span}_R\{e_j\mid j>i\}\subset R\otimes\Lambda$；
\item $\mathbf v_i=a_{i+1}e_{i+1}\in\Lambda_i/\Lambda_{i+1}$；
\item $\Lambda^i=\operatorname{Span}_R\{e_j\mid j\le i+d\}\subset R\otimes\Lambda$；
\item $\mathbf v_i=e_{i+d}+b_ie_{i+d-1}\in\Lambda^i/\Lambda^{i-2}$。
\end{enumerate}
特别地，有基点 $\star_d=j_d(1,0)\in\mathrm{Bun}^d_{n,1,2}$。记 $\star_d$ 的 $k[t]$-链（相应地 $k[t^{-1}]$-链）为 $\Lambda_*(\star_d)$（相应地 $\Lambda^*(\star_d)$）。注意任意点 $j_d(a,b)$ 的底向量丛是丛
\[
E_d=\mathcal O(m+1)^{\oplus r}\oplus\mathcal O(m)^{\oplus n-r},
\]
其中 $m\in\mathbb{Z}$、$0\le r\le n-1$ 由 $d=mn+r$ 唯一确定。

\subsection{与 $GL_n$ 的标准表示相伴的 Kloosterman 层}

设 $\chi:T(k)\to\mathbb{Q}_\ell^\times$ 是特征，它在 $T$ 上定义 Kummer 局部系统 $\mathrm{Kum}_\chi$。由定理 1，$\mathrm{Bun}^d_{n,1,2}$ 上的 perverse 层
\[
A_{\varphi,\chi}^d=j_{d,!}\big(\mathrm{Kum}_\chi[n]\boxtimes\varphi^*\mathrm{AS}_\psi[n]\big)
\]
构成 $\mathrm{Bun}_{n,1,2}$ 上的 Hecke 特征层。

考虑由余权 $\omega_1=(1,0,\cdots,0)$ 给出的 Hecke 对应 $\Hecke_{\omega_1}$（$\omega_1$ 定义表示 $\mathrm{Std}$）。把它限制到 $\mathrm{Bun}^0$ 与 $\mathrm{Bun}^1$ 的分支以及曲线 $\PP^1\setminus\{0,\infty\}$ 上：
\begin{equation}
\begin{tikzcd}[column sep=huge, row sep=large]
& \Hecke_{\omega_1} \arrow[dl, "\mathrm{pr}_1"'] \arrow[dr, "\mathrm{pr}_2"] & \\
\mathrm{Bun}^0_{n,1,2} & & \mathrm{Bun}^1_{n,1,2}\times(\PP^1\setminus\{0,\infty\})
\end{tikzcd}
\end{equation}
我们想计算 $\PP^1\setminus\{0,\infty\}$ 上由下式刻画的特征值局部系统 $\Kl^{\mathrm{Std}}_{GL_n}(\varphi,\chi)$：
\[
\mathrm{pr}_{2,!}\mathrm{pr}_1^*A^0_{\varphi,\chi}[n-1]\cong A^1_{\varphi,\chi}\boxtimes\Kl^{\mathrm{Std}}_{GL_n}(\varphi,\chi).
\]
$\Hecke_{\omega_1}$ 上的相交上同调层 $\mathrm{IC}_{\omega_1}$ 就是 $\mathbb{Q}_\ell[n-1]$，因为 $\mathrm{pr}_1:\Hecke_{\omega_1}\to\mathrm{Bun}_{n,1,2}$ 是 $\mathbb{P}^{n-1}$-丛\footnote{在这个计算中，我们不把 $\mathrm{IC}_{\omega_1}$ 归一化为权 $0$，如注 2.9 那样。故严格说来，层 $\Kl^{\mathrm{Std}}_{GL_n}(\varphi,\chi)$ 与 §2.3 中定义的那个相差一个 Tate 扭转 $\big(\tfrac{n-1}{2}\big)$。}。把图 (3.1) 限制到 $\star_1\times(\PP^1\setminus\{0,\infty\})\subset\mathrm{Bun}^1_{n,1,2}\times(\PP^1\setminus\{0,\infty\})$ 在 $\mathrm{pr}_2$ 下的纤维，以及大胞腔在 $\mathrm{pr}_1$ 下的纤维 $\mathrm{GR}^\circ_{\omega_1}$，得到
\begin{equation}
\begin{tikzcd}[column sep=large, row sep=large]
\mathrm{GR}^\circ_{\omega_1} \arrow[r, hook] \arrow[d, "\mathrm{pr}_1"'] \arrow[dr, "\pi"] & \mathrm{GR}_{\omega_1} \arrow[d, "\mathrm{pr}_1"] & \\
T\times\mathbb{G}_a^n \arrow[r, "j_0"'] & \mathrm{Bun}^0_{n,1,2} & \PP^1\setminus\{0,\infty\}
\end{tikzcd}
\end{equation}
开子概形 $\mathrm{GR}^\circ_{\omega_1}\subset\mathrm{GR}_{\omega_1}$ 的 $R$-点分类 $R[t,t^{-1}]$-同态 $M:R\otimes\Lambda\to R\otimes\Lambda$，它把链 $(\Lambda_*(\star_0),\Lambda^*(\star_0))$ 送到 $(\Lambda_*(\star_1),\Lambda^*(\star_1))$，模去左复合以 $(R\otimes\Lambda,\Lambda_*(\star_0),\Lambda^*(\star_0))$ 的自同构。关于 $R\otimes\Lambda$ 的 $R[t,t^{-1}]$-基 $\{e_1,\cdots,e_n\}$，任何这样的 $M$ 形如
\[
\begin{pmatrix}
a_1 & tb_n & & \\
b_1 & a_2 & & \\
& b_2 & a_3 & \\
& & \ddots & \ddots\\
& & & b_{n-1} & a_n
\end{pmatrix},
\]
模去右乘对角矩阵（这里 $a_i,b_i\in R^\times$）。因此可以把 $M$ 的矩阵归一化为
\begin{equation}
M=\begin{pmatrix}
a_1 & t & & \\
1 & a_2 & & \\
& 1 & \ddots & \\
& & \ddots & a_{n-1} & t\\
& & & 1 & a_n
\end{pmatrix}.
\end{equation}
换句话说，把 $M$ 送到 $(-a_1,\cdots,-a_n)$ 给出等同 $U\cong\mathbb{G}_m^n$。

现在用这个等同描述投影 $\mathrm{pr}_1$ 与 $\pi$。利用同构 $M$，可以把 $\star_1$ 的数据中的向量 $\{\mathbf v_i=e_{i+1}\}$、$\{\mathbf v_i=e_{i+1}\}$ 拉回，得到点 $\mathrm{pr}_1(M)=(\Lambda,\Lambda_*(\star_0),\mathbf v_i,\Lambda^*(\star_0),\mathbf v_i)\in\mathrm{Bun}^0_{n,1,2}$ 的相应向量：
\[
\mathbf v_i=a_{i+1}^{-1}e_{i+1},\quad i=0,\cdots,n-1;\qquad
\mathbf v_i=e_i-a_ie_{i-1},\quad i=1,\cdots,n.
\]
这意味着 (3.3) 中的 $M$ 对应的 $\mathrm{pr}_1(M)$ 的坐标为
\[
\mathrm{pr}_1(M)=j_0(a_1^{-1},\cdots,a_n^{-1},-a_1,\cdots,-a_n).
\]
另一方面，点 $\pi(M)\in\PP^1\setminus\{0,\infty\}$ 是使 $\det(M)=0$ 的 $t$ 的值，即
\[
\pi(M)=(-1)^na_1\cdots a_n.
\]
总之，$\mathrm{GR}^\circ_{\omega_1}$ 上的映射 $\mathrm{pr}_1$ 与 $\pi$ 可等同于如下映射
\[
\begin{tikzcd}[column sep=huge, row sep=large]
& \mathbb{G}_m^n \arrow[dl, "{(-\operatorname{inv},\iota)}"'] \arrow[dr, "\operatorname{mult}"] & \\
\mathbb{G}_m^n\times\mathbb{G}_a^n & & \mathbb{G}_m=\PP^1\setminus\{0,\infty\}
\end{tikzcd}
\]
其中 $\mathrm{inv}:\mathbb{G}_m^n\to\mathbb{G}_m^n$ 是逐坐标求逆，$\iota:\mathbb{G}_m^n\hookrightarrow\mathbb{G}_a^n$ 是自然包含。

于是得到

\begin{proposition}
与 Hecke 特征层 $A_{\varphi,\chi}$ 及对偶群 $GL_{n,\mathbb{Q}_\ell}=\widehat{GL}_n$ 的标准表示相伴的 Kloosterman 层 $\Kl^{\mathrm{Std}}_{GL_n}(\varphi,\chi)$ 形如
\[
\Kl^{\mathrm{Std}}_{GL_n}(\varphi,\chi)\cong\mathrm{mult}_!\big((-\mathrm{inv})^*\mathrm{Kum}_\chi\otimes\varphi^*\mathrm{AS}_\psi\big)[n-1].
\]
这里把 $\varphi:\mathbb{G}_a^n\to\mathbb{G}_a$ 到 $\mathbb{G}_m^n$ 的限制仍记为 $\varphi$。
\end{proposition}

\begin{remark}
当 $\varphi=\mathrm{add}:\mathbb{G}_a^n\to\mathbb{G}_a$ 是坐标求和，且把 $\chi$ 写成 $n$ 个乘性特征 $(\chi_1,\cdots,\chi_n)$ 时，Kloosterman 层 $\Kl^{\mathrm{Std}}_{GL_n}(\varphi,\chi)$ 与 Katz 在 [25, §4.1] 中定义的 Kloosterman 层 $\mathrm{Kl}(\psi;\chi_1,\cdots,\chi_n;1,\cdots,1)$ 相同，后者是 Deligne 的 Kloosterman 层 [10, §7] 的推广。
\end{remark}

% ============================================================
%  ch06 —— §4 定理 1 的证明（原文 p.22–p.26）
% ============================================================
\routenote{p.22--p.26}

\section{定理 1 的证明}

本节证明定理 1。

\subsection{第一步：$\Hk_V(A)$ 是 perverse 的}

我们要证明对每个 $V\in\Rep(\ckG)$，复形 $\Hk_V(A_{\varphi,\chi})[-1]$ 是 perverse 层。为简化记号，只考虑层 $A=A_\varphi$；对 $A_{\varphi,\chi}$ 的证明完全相同，只需在论证中把 $G(0,2)$ 处处换为 $G(1,2)$。

回忆 §2.3 的卷积图：
\[
\begin{tikzcd}[column sep=large, row sep=large]
& \Hecke^{\Aff^1}_{G(0,2)} \arrow[dl, "\mathrm{pr}_1"'] \arrow[dr, "\mathrm{pr}_2"] \arrow[d, "\mathrm{pr}_{\Aff^1}"] & \\
\Bun_{G(0,2)} & \Bun_{G(0,2)}\times\Aff^1 \arrow[l, "\mathrm{proj}_1"'] & \Aff^1
\end{tikzcd}
\]
我们的证明基于几个简单的几何观察。首先回忆上图中映射 $\mathrm{pr}_i$ 是局部平凡纤维化：

\begin{remark}
映射 $\mathrm{pr}_1$ 是局部平凡纤维化，即存在光滑图册 $U\to\Bun_{G(0,2)}$，使得
\[
U\times_{\Bun_{G(0,2)}}\Hecke^{\Aff^1}_{G(0,2)}\cong U\times\mathrm{GR}_{G(0,2)}.
\]
此外，$\mathrm{pr}_2$ 在这个图册上也局部平凡，即
\[
(U\times\Aff^1)\times_{\Bun_{G(0,2)}\times\Aff^1}\Hecke^{\Aff^1}_{G(0,2)}\cong U\times\mathrm{GR}_{G(0,2)}.
\]
\end{remark}

\begin{proof}
为找到满足这些条件的图册 $p:U\to\Bun_{G(0,2)}$，我们只需要 $U$ 上对应于 $p$ 的 $G(0,2)$-丛族在 $U\times\Aff^1$ 上平凡。由命题 1.1，对任意 $U$，这个条件在 $U$ 的 étale 拓扑下局部地成立。
\end{proof}

对 $\gamma\in\Omega$，记 $j_\gamma:I(1)/I(2)\hookrightarrow\Bun_{G(0,2)}$ 为典范嵌入。回忆 $A^\gamma=j_{\gamma,!}\mathrm{AS}_\varphi[\dim(\Bun_{G(0,2)})]$。记 $j'_\gamma:\mathrm{pr}_1^{-1}(I(1)/I(2))\hookrightarrow\Hecke^{\Aff^1}_{G(0,2)}$ 为这个嵌入到 Hecke 叠中的原像。

\begin{remark}
$\mathrm{pr}_2$ 到 $\mathrm{pr}_1^{-1}(I(1)/I(2))$ 的限制
\[
\mathrm{pr}_2:\ \mathrm{pr}_1^{-1}(I(1)/I(2))\to\Bun_{G(0,2)}\times\Aff^1
\]
是仿射的。
\end{remark}

\begin{proof}
由推论 1.3，开子集 $I(1)/I(2)\hookrightarrow\Bun_{G(0,2)}$ 由线丛 $L_i$ 的截面的非零性定义，而这些线丛可定义为忠实表示 $G\to GL_n$ 下 $GL_n$ 的相应线丛的拉回。若 $G=GL_n$，则对任意 $x\in\Aff^1$，这些丛的拉回生成 Picard 群 $\mathrm{Gr}_{GL_n(0,2),x}$，后者是 ind-投射概形（例如 [14, Theorem 7 and 8]）。因此 $I(1)/I(2)$ 在 $\mathrm{GR}_{GL_n(0,2)}$ 中的原像在 $\Aff^1$ 上是仿射的。

对一般的 $G$，ind-概形 $\mathrm{Gr}_{G(0,2),x}$ 通常构造为 $\mathrm{Gr}_{GL_n}$ 对某个合适忠实表示的闭子概形（[32]）。故断言又由 $G=GL_n$ 的情形推出。
\end{proof}

现在可以证明第一步。固定任意 $V\in\Rep(\ckG)$ 及 $\Hecke^{\Aff^1}_{G(0,2)}$ 上相应的 perverse 层 $\mathrm{IC}_V$。我们断言注 4.1 蕴涵
\[
j'_!(\mathrm{pr}_1^*\mathrm{AS}_\varphi\otimes\mathrm{IC}_V)\cong j'_*(\mathrm{pr}_1^*\mathrm{AS}_\varphi\otimes\mathrm{IC}_V).
\]
为看出这点，回忆 $j_!$ 与 $j_*$ 的构造与光滑基变换交换。故为证明所断言同构，可取光滑图册 $p:U\to\Bun_{G(0,2)}$ 使得 $\mathrm{pr}_1$ 的拉回同构于投影 $U\times\mathrm{GR}_G\to U$。但此时断言由同构 $j_!\mathrm{AS}_\varphi=j_*\mathrm{AS}_\varphi$（引理 2.3）推出。

特别地，我们发现 $j'_!(\mathrm{pr}_1^*A\otimes\mathrm{IC}_V)[1]$ 是 perverse 层，且
\[
\Hk_V(A)=(\mathrm{pr}_2\circ j')_!(\mathrm{pr}_1^*(A)\otimes\mathrm{IC}_V)
=(\mathrm{pr}_2\circ j')_*(\mathrm{pr}_1^*(A)\otimes\mathrm{IC}_V).
\]
由注 4.2，映射 $(\mathrm{pr}_2\circ j')$ 是仿射的。因此 $(\mathrm{pr}_2\circ j')_*$ 关于 perverse $t$-结构右正合（[3, §4.1.1]）。于是 $\Hk_V(A)[1]$ 必是 perverse 的。

\subsection{第二步：$A$ 是特征层}

我们已在注 2.7 中指出，$I(1)/I(2)$ 在 $\Bun_{G(0,2)}$ 上的作用扩充为卷积图上的作用。故 $\Hk_V(A)$ 是 $(I(1)/I(2),\varphi)$-等变的。由引理 2.3，这意味着对任意 $\gamma\in\pi_0(\Bun_{G(0,2)})$ 有
\[
\Hk_V(A)|_{\Bun^\gamma_{G(0,2)}\times X}=A^\gamma\boxtimes E^\gamma_V,
\]
其中 $E^\gamma_V[1]$ 是 $X$ 上的 perverse 层。

我们断言层 $E^\gamma_V$ 不依赖于 $\gamma$。为此可假设 $\mathrm{IC}_V$ 只支撑在一个连通分支 $\mathrm{GR}^\gamma_G\subset\mathrm{GR}_G$ 上，因为函子 $\Hk_V$ 同构于由 $\mathrm{IC}_V$ 到各连通分支的限制所定义的函子之直和。

注意对任意 $\gamma'\in\Omega$，Hecke 算子 $\Hk_{\gamma'}$ 与 $\Hk_V$ 交换，因为支撑在 $\PP^1$ 上不同点的 Hecke 算子交换。按定义 $A$ 是算子 $\Hk_{\gamma'}$ 的特征层。于是
\[
A^0\boxtimes E^0_V=\Hk_V(A^{-\gamma})=\Hk_V(\Hk_{-\gamma}(A^0))
=\Hk_{-\gamma}(\Hk_V(A^0))=\Hk_{-\gamma}(A^\gamma\boxtimes E^\gamma_V)
=A^0\boxtimes E^\gamma_V.
\]
这说明层 $E^\gamma_V$ 典范同构于 $E^0_V$，故可去掉指标 $\gamma$。

诸 $\Hk_V$ 与表示的张量积相容，故特别有
\[
\Hk_{V\otimes V'}(A)=\Hk_V\circ\Hk_{V'}(A)=A\boxtimes(E_V\otimes E_{V'}).
\]
这意味着 $E_V\otimes E_{V'}[1]$ 仍是 perverse 层。因此 $E_V$ 必是集中在上同调次数 $0$ 的 perverse 层。故对任意 $V$，复形 $E_V$ 是一个层，且 $E_V=j_{\eta,*}E_{V,\eta}$，其中 $j_\eta:\eta\hookrightarrow\PP^1$ 是一般点的包含。

此外，我们得到取值于 $\eta$ 上局部系统范畴的张量函子 $V\mapsto(E_V)_\eta$。这是刚性张量范畴，故由 [12, Theorem 3.2] 它定义 $\eta$ 上的 $\ckG$-局部系统。我们需要证明 $E_{V,\eta}$ 扩充为 $\PP^1\setminus\{0,\infty\}$ 上的局部系统。

对平凡 Hecke 算子 $\Hk_1$ 有典范同构 $E_1=\mathbb{Q}_\ell$，连同映射 $E_1\to E_V\otimes E_{V^*}\to E_1$，其复合等于乘 $\dim(V)$。我们已知 $E_V=j_{\eta,*}E_{V,\eta}$。故对 $\mathbb{G}_m$ 的任意几何点 $x$，纤维 $E_{V,x}$ 是几何一般纤维 $E_{V,\eta}$ 的子空间，且典范映射 $\mathrm{id}_\eta:\mathbb{Q}_{\ell,\eta}\to(E_V\otimes E_{V^*})_\eta$ 经过 $E_{V,x}\otimes E_{V^*,x}$ 分解。这意味着层 $E_V$ 局部常值，因为 $\mathrm{id}_\eta$ 对应于 $E_{V,\eta}$ 的恒等。

于是诸 $E_V$ 定义从 $\Rep(\ckG)$ 到 $\PP^1\setminus\{0,\infty\}$ 上局部系统范畴 $\Loc(\PP^1\setminus\{0,\infty\})$ 的张量函子。由于这又是刚性张量范畴，它定义 $\PP^1\setminus\{0,\infty\}$ 上的 $\ckG$-局部系统。

\subsection{第三步：驯点处的单值}

如定理 1(2) 的陈述那样，现在假设 $\mathbb{G}=G\times\PP^1$ 是常分裂群。为计算 $\Kl^{\ckG}(\varphi)$ 在 $\{0\}$ 处的单值，我们改述 Bezrukavnikov [5] 的一个论证。他的论证依赖于 Gaitsgory [16] 关于中心层的结果以及 Gabber 的结果——邻近闭链上的单值滤过与权滤过一致 [2]。

为解释这个论证，需回忆 Gaitsgory 的构造（[16]）。他考虑了图
\[
\begin{tikzcd}[column sep=large, row sep=large]
\mathrm{GR}^{\PP^1\setminus\{0,\infty\}}_G \arrow[r, hook, "j_{\mathrm{GR}}"] \arrow[d] & \mathrm{GR}_G \arrow[d] & \mathrm{GR}_{G,0}=\mathrm{Fl}_G \arrow[l, hook'] \arrow[d] \\
\PP^1\setminus\{0,\infty\} \arrow[r, hook, "j"] & \Aff^1 & \{0\} \arrow[l, hook']
\end{tikzcd}
\]
以及诱导的邻近闭链函子 $\Psi:\mathrm{Perv}(\mathrm{GR}^{\PP^1\setminus\{0,\infty\}}_G)\to\mathrm{Perv}_{I_0}(\mathrm{Fl}_G)$。他证明了 $\Psi(\mathrm{IC}_V)$ 上的单值作用是幂幺的。

由于映射 $\mathrm{pr}_1:\Hecke^{\Aff^1}_{G(0,2)}\to\Bun_{G(0,2)}$ 局部同构于与纤维 $\mathrm{GR}_{G(0,2)}$ 之积，我们知道
\[
\Psi(\mathrm{pr}_1^*(A)\otimes\mathrm{IC}_V)=\mathrm{pr}_1^*(A)\otimes\Psi(\mathrm{IC}_V).
\]
特别地，由 Gaitsgory 的结果 [16, Theorem 2] 知这个层上的单值作用是幂幺的。因此（[19, Lemma 5.6]）$\mathrm{pr}_{2,!}(\Psi(\mathrm{pr}_1^*(A)\otimes\mathrm{IC}_V))$ 上的单值作用仍是幂幺的。由于取邻近闭链与真推出交换，单值作用在
\begin{equation}
\mathrm{pr}_{2,!}(\mathrm{pr}_1^*(A)\otimes\Psi(\mathrm{IC}_V))=\mathrm{pr}_{2,!}(\Psi(\mathrm{pr}_1^*(A)\otimes\mathrm{IC}_V))=\Psi(A\boxtimes E_V)=A\boxtimes\Psi(E_V)
\end{equation}
上也是幂幺的。按定义，$\Psi(E_V)$ 是 $E_V$ 在 $0$ 的去心形式邻域上方几何点 $\operatorname{Spec}K^{\mathrm{sep}}_0$ 处的茎，带有作为单值的 $\operatorname{Gal}(K^{\mathrm{sep}}_0/K_0)$-作用。把 (4.3) 限制到平凡 $G(0,2)$-丛，于是得到
\begin{equation}
\Kl^V_{\ckG}|_{\operatorname{Spec}K^{\mathrm{sep}}_0}=R\Gamma_c(\mathrm{pr}_1^*A\otimes\Psi(\mathrm{IC}_V)),
\end{equation}
其中惯性群 $I_0\subset\operatorname{Gal}(K^{\mathrm{sep}}_0/K_0)$ 驯顺且幂幺地作用。

我们须证明单值作用由一个主幂幺元给出。回忆仿射旗流形 $\mathrm{Fl}_G=\mathrm{GR}_{G,0}$ 上 $G(0,2)(k[[t]])=I(0)^{\mathrm{opp}}$-轨道由 Iwahori--Weyl 群 $\widetilde W$ 参数化（[32, Prop.~8.1]）。这些轨道闭包的相交上同调层记为 $\mathrm{IC}_{\tilde w}$。与这些层作卷积定义 Hecke 算子 $\Hk_{\tilde w}(A):=\mathrm{pr}_{2,!}(\mathrm{pr}_1^*(A)\otimes\mathrm{IC}_{\tilde w})$。

由 Görtz 与 Haines 的结果 [19, Corollary 1.2]，层 $\Psi(\mathrm{IC}_V)$ 有一个滤链，其伴随分次层同构于 $\mathrm{IC}_w(i)$（$w\in\widetilde W$，$i\in\mathbb{Z}$），且 $\mathrm{IC}_1(i)$ 的重数等于 $\dim H^{2i}(\mathrm{Gr}_G,\mathrm{IC}_V)$。

现在，对任意长度 $l(\tilde w)>0$ 的 $\tilde w\in\widetilde W$，存在单反射 $s$ 使 $\tilde w=\tilde w's$ 且 $l(\tilde w')<l(\tilde w)$。记 $P_s:=I(0)^{\mathrm{opp}}\cup I(0)^{\mathrm{opp}}sI(0)^{\mathrm{opp}}$ 为由 $I(0)^{\mathrm{opp}}$ 与 $s$ 生成的 parahoric 子群，于是投影 $\mathrm{pr}_s:\mathrm{Fl}\to G((t))/P_s$ 是 $\mathbb{P}^1$-丛。由此看出层 $\mathrm{IC}_{\tilde w}$ 形如 $\mathrm{pr}_s^*(\mathrm{IC}_{\tilde w'})$。因此复形 $\Hk_{\tilde w}(A)$ 在更大的 parahoric 子群 $P_s\supset I(0)^{\mathrm{opp}}$ 下不变。由于 $\Hk_{\tilde w}(A)$ 又是 $(I(1)/I(2),\mathrm{AS}_\varphi)$-等变的，这蕴涵 $\Hk_{\tilde w}(A)=0$。

于是只有长度为 $0$ 的 Hecke 算子非平凡地作用在 $A$ 上。因此，由 Görtz 与 Haines 的结果，$\Psi(\Kl^V_{\ckG})=\Kl^V_{\ckG,\eta}$ 的权滤过的维数由 $\dim H^{2i}(\mathrm{Gr}_G,\mathrm{IC}_V)$ 给出。由于单值滤过与权滤过一致 [2]，单值必按主幂幺元作用。这证明了我们的断言。

\begin{remark}
Görtz 与 Haines 文中所需的关键要素是层 $\Psi(\mathrm{IC}_\mu)$ 的迹函数公式。这个迹由 Iwahori--Hecke 代数中心元素的 Bernstein 公式给出（见 [19, §2.7]）。也可以由这个公式直接推出我们的命题。
\end{remark}

% ============================================================
%  ch07 —— §5.1–5.2（原文 p.26–p.31）
% ============================================================
\routenote{p.26--p.31}

\section{Kloosterman 层的上同调性质}

本节设 $G$ 是 $k$ 上分裂几乎单群，视作 $\PP^1$ 上的常群概形。回忆对每个根 $\alpha\in\Phi$，$U_\alpha$ 是 $G$ 中的根群。对每个根 $\alpha$ 固定同构 $u_\alpha:\mathbb{G}_a\xrightarrow{\sim}U_\alpha$。

\medskip\noindent\textbf{记号.} 对 $k$ 上概形 $X$ 与 $X$ 上 $\mathbb{Q}_\ell$-层复形 $F$，$\chi_c(X,F)$ 与 $H^*(X,F)$ 指 $F$ 拉回到 $X\otimes_k\bar k$ 后的 Euler 特征与上同调。

\subsection{Euler 特征、Swan 导子与上同调刚性}

有两个维数较小的 $V\in\Rep(\ckG)$ 的一致选择。其一是 $V=\mathrm{Ad}=\ckg$，即伴随表示；其二是 $V=V_{\theta^\vee}$，即非零权由 $\ckg$ 的短根组成的表示。称 $V_{\theta^\vee}$ 为 $\ckG$ 的\emph{拟极小表示}（当 $\ckG$ 单链时它与伴随表示重合）。关于 $V_{\theta^\vee}$ 的基本事实汇总在附录 §10。记 $G$ 的长（相应地短）单根个数为 $r_\ell(G)$（相应地 $r_s(G)$），称之为 $G$ 的长秩（相应地短秩）。这些数的等价描述见引理 10.1。

记 $\Kl^{\theta^\vee}_{\ckG}(\varphi,\chi)$（相应地 $\Kl^{\mathrm{Ad}}_{\ckG}(\varphi,\chi)$）为与 $\Kl_{\ckG}(\varphi,\chi)$ 及 $\ckG$ 的拟极小表示（相应地伴随表示）相伴的局部系统。

\begin{mainthm}[4]
\begin{enumerate}[label=(\arabic*), leftmargin=1.9em, itemsep=.25em]
\item $-\chi_c(\PP^1\setminus\{0,\infty\},\Kl^{\theta^\vee}_{\ckG}(\varphi,\chi))$ 等于 $G$ 的长单根个数。
\item 设当 $G$ 非单链时 $\operatorname{char}(k)$ 对 $G$ 良好，则 $-\chi_c(\PP^1\setminus\{0,\infty\},\Kl^{\mathrm{Ad}}_{\ckG}(\varphi,\chi))$ 等于 $G$ 的秩。
\end{enumerate}
\end{mainthm}

由 Grothendieck--Ogg--Shafarevich 公式，得到

\begin{corollary}
\begin{enumerate}[label=(\arabic*), leftmargin=1.9em, itemsep=.2em]
\item $\Swan_\infty(\Kl^{\theta^\vee}_{\ckG}(\varphi,\chi))=r_s(\ckG)$。
\item 在与定理 4(2) 相同的 $\operatorname{char}(k)$ 假设下，$\Swan_\infty(\Kl^{\mathrm{Ad}}_{\ckG}(\varphi,\chi))=r(\ckG)$。
\end{enumerate}
\end{corollary}

以下各小节将致力于证明定理 4。先推出一些推论。

\begin{lemma}
设主幂幺元作用在 $V$ 上无平凡 Jordan 块，则 $H^0(\PP^1\setminus\{0,\infty\},\Kl^V_{\ckG}(\varphi))=0$。
\end{lemma}

\begin{proof}
设不然，则 $\Kl^V_{\ckG}(\varphi)$ 以常层为子局部系统。按构造 $\Kl^V_{\ckG}(\varphi)$ 是纯的。由 [3, Th.~5.3.8]，在 $\PP^1\setminus\{0,\infty\}\otimes_k\bar k$ 上 $\Kl^V_{\ckG}(\varphi)$ 是单 perverse 层之直和。故该常子层必是直和项。

另一方面，按假设，$V$ 不含任何使 $\PP^1$ 中 $\{0\}$ 处的驯单值（由定理 1，即主幂幺元）平凡作用的直和项。这导致矛盾。
\end{proof}

设 $i_0:\{0\}\hookrightarrow\PP^1$，$i_\infty:\{\infty\}\hookrightarrow\PP^1$，$j:\PP^1\setminus\{0,\infty\}\hookrightarrow\PP^1$ 为包含。对 $\PP^1\setminus\{0,\infty\}$ 上任意局部系统 $L$，滥用记号，以 $j_{!*}L$ 表示 $j_{!*}(L[1])[-1]$。

如定理 2 之前的讨论那样，局部 Galois 群 $\operatorname{Gal}(K^{\mathrm{sep}}_\infty/K_\infty)$ 与 $\operatorname{Gal}(K^{\mathrm{sep}}_0/K_0)$（从而惯性群 $I_0$ 与 $I_\infty$）作用在 $\Kl^V_{\ckG}(\varphi)$ 的相应几何茎上，定义 $V$ 上在共轭意义下的表示。

\begin{proposition}
在与定理 4(2) 相同的 $\operatorname{char}(k)$ 假设下，有
\begin{enumerate}[label=(\arabic*), leftmargin=1.9em, itemsep=.2em]
\item （上同调刚性）$H^*(\PP^1_k,j_{!*}\Kl^{\mathrm{Ad}}_{\ckG}(\varphi))=0$。
\item $\ckg^{I_\infty}=0$（注意 $\ckg$ 是 $V=\mathrm{Ad}$ 的空间）。
\end{enumerate}
以 $\Kl^{\theta^\vee}_{\ckG}(\varphi)$ 代替 $\Kl^{\mathrm{Ad}}_{\ckG}(\varphi)$ 时同样成立，且对 $\operatorname{char}(k)$ 无限制。
\end{proposition}

\begin{proof}
这些陈述是几何的，故在证明中忽略 Tate 扭转。由于伴随表示 $\ckg$ 自对偶，局部系统 $\Kl^{\mathrm{Ad}}_{\ckG}(\varphi)$ 也自对偶。固定这样的同构 $\Kl^{\mathrm{Ad}}_{\ckG}(\varphi)\xrightarrow{\sim}(\Kl^{\mathrm{Ad}}_{\ckG}(\varphi))^\vee$，得到同构
\begin{equation}
\begin{aligned}
H^0(\PP^1,j_{!*}\Kl^{\mathrm{Ad}}_{\ckG}(\varphi))&\cong H^2(\PP^1,j_{!*}\Kl^{\mathrm{Ad}}_{\ckG}(\varphi))^\vee,\\
H^0(\PP^1\setminus\{0,\infty\},\Kl^{\mathrm{Ad}}_{\ckG}(\varphi))&\cong H^2_c(\PP^1\setminus\{0,\infty\},\Kl^{\mathrm{Ad}}_{\ckG}(\varphi))^\vee.
\end{aligned}
\end{equation}
在 $D^b_c(\PP^1,\mathbb{Q}_\ell)$ 中有 distinguished triangle
\begin{equation}
j_!\Kl^{\mathrm{Ad}}_{\ckG}(\varphi)\to j_{!*}\Kl^{\mathrm{Ad}}_{\ckG}(\varphi)\to H^0i_0^*j_*\Kl^{\mathrm{Ad}}_{\ckG}(\varphi)\oplus H^0i_\infty^*j_*\Kl^{\mathrm{Ad}}_{\ckG}(\varphi)\to,
\end{equation}
它诱导长正合列
\begin{equation}
\begin{aligned}
0=H^0_c(\PP^1\setminus\{0,\infty\},\Kl^{\mathrm{Ad}}_{\ckG}(\varphi))&\to H^0(\PP^1,j_{!*}\Kl^{\mathrm{Ad}}_{\ckG}(\varphi))\to\ckg^{I_0}\oplus\ckg^{I_\infty}\xrightarrow{d}\\
&\xrightarrow{d}H^1_c(\PP^1\setminus\{0,\infty\},\Kl^{\mathrm{Ad}}_{\ckG}(\varphi))\to H^1_c(\PP^1,j_{!*}\Kl^{\mathrm{Ad}}_{\ckG}(\varphi))\to 0\to\\
&\to H^2_c(\PP^1\setminus\{0,\infty\},\Kl^{\mathrm{Ad}}_{\ckG}(\varphi))\to H^2_c(\PP^1,j_{!*}\Kl^{\mathrm{Ad}}_{\ckG}(\varphi))\to 0.
\end{aligned}
\end{equation}
由第二个正合列先得
\[
\begin{aligned}
H^2(\PP^1,j_{!*}\Kl^{\mathrm{Ad}}_{\ckG}(\varphi))&\cong H^2_c(\PP^1\setminus\{0,\infty\},\Kl^{\mathrm{Ad}}_{\ckG}(\varphi))\\
&\cong H^0(\PP^1\setminus\{0,\infty\},\Kl^{\mathrm{Ad}}_{\ckG}(\varphi))^\vee=0,
\end{aligned}
\]
最后一项的消没来自把引理 5.2 用于 $V=\ckg$。由对偶性，$H^0(\PP^1,j_{!*}\Kl^{\mathrm{Ad}}_{\ckG}(\varphi))=0$。由第一个正合列与 $H^0$ 的消没，连接同态 $d$ 是单射。

一方面，由定理 1(2)，$\dim\ckg^{I_0}=r(\ckG)$，因为 $I_0$ 通过主幂幺元作用在 $\ckg$ 上。另一方面，由定理 4，$\dim H^1_c(\PP^1\setminus\{0,\infty\},\Kl^{\mathrm{Ad}}_{\ckG}(\varphi))=-\chi_c(\PP^1\setminus\{0,\infty\},\Kl^{\mathrm{Ad}}_{\ckG}(\varphi))=r(\ckG)$（这里再次用了 $H^i_c(\PP^1\setminus\{0,\infty\},\Kl^{\mathrm{Ad}}_{\ckG}(\varphi))$ 对 $i=0,2$ 的消没）。因此 $d$ 必是同构。这蕴涵 $\ckg^{I_\infty}=0$。由第二个正合列，得到 $H^1(\PP^1,j_{!*}\Kl^{\mathrm{Ad}}_{\ckG}(\varphi))=0$。

对 $\Kl^{\theta^\vee}_{\ckG}(\varphi)$ 的陈述同样证明。由引理 10.1 证明中的等式 (10.2)，可以把引理 5.2 用于 $V_{\theta^\vee}$。
\end{proof}

\begin{remark}
沿用 Katz [27, §5.0]，称 $\PP^1\setminus\{0,\infty\}\otimes_k\bar k$ 上的 $\ckG$-局部系统 $L$ 为\emph{上同调刚性的}，若 $H^1(\PP^1,j_{!*}L^{\mathrm{Ad}})=0$。我们把 $H^1(\PP^1,j_{!*}L^{\mathrm{Ad}})$ 设想为 $L$ 的、在 $0$ 与 $\infty$ 周围形式去心圆盘上同构型固定的无穷小变形空间，尽管这种变形的概念尚未定义。命题 5.3 蕴涵 $\Kl_{\ckG}(\varphi)$ 是上同调刚性的，这为其物理刚性提供了证据。更精确的猜想见 §7。
\end{remark}

\subsection{一般计算方法}

以下计算中只考虑 $\Bun_{G(1,2)}$ 的中性分支。故可假设 $G$ 单连通。记 $\star\in\Bun_{G(1,2)}$ 为平凡丛对应的点。为计算层 $\Kl^V_{\ckG}(\varphi,\chi)$，把卷积图 (2.6) 限制到 $\star\times\PP^1\setminus\{0,\infty\}\subset\Bun_{G(1,2)}\times\PP^1\setminus\{0,\infty\}$：
\[
\begin{tikzcd}[column sep=large, row sep=large]
& \mathrm{GR} \arrow[dl, "\mathrm{pr}_1"'] \arrow[dr, "\mathrm{pr}_2"] & \\
\Bun_{G(1,2)} & & \star\times\PP^1\setminus\{0,\infty\}.
\end{tikzcd}
\]
由真基变换，$\Kl^V_{\ckG}(\varphi,\chi)=\mathrm{pr}_{2,!}(\mathrm{pr}_1^*A\otimes\mathrm{IC}_V)$。此外，$A$ 只支撑在大胞腔 $j:T\times I(1)/I(2)\subset\Bun_{G(1,2)}$ 上。记 $\mathrm{GR}^\circ\subset\mathrm{GR}$ 为大胞腔的原像，并写 $I(1)/I(2)\cong U/[U,U]\times U_{-\theta}$。按这个分解写 $\varphi=(\varphi_+,\varphi_0)$。

得到图
\[
\begin{tikzcd}[column sep=large, row sep=large]
\mathrm{GR}^\circ \arrow[r, hook] \arrow[dr, "\mathrm{pr}^\circ_2"'] \arrow[d, "{(f_T,f_0,f_+)}"'] & \mathrm{GR} \arrow[dl, "\mathrm{pr}_1"'] \arrow[dr, "\mathrm{pr}_2"] & \\
T\times U_{-\theta}\times U/[U,U] \arrow[r, hook] & \Bun_{G(1,2)} & \PP^1\setminus\{0,\infty\}
\end{tikzcd}
\]
由真基变换，
\begin{equation}
\Kl^V_{\ckG}(\varphi,\chi)=\mathrm{pr}^\circ_{2,!}\big(f_T^*\mathrm{Kum}_\chi\,f_0^*\mathrm{AS}_{\varphi_0}\otimes f_+^*\mathrm{AS}_{\varphi_+}\otimes\mathrm{IC}_V\big).
\end{equation}
故需描述子空间 $\mathrm{GR}^\circ$ 并计算映射 $f_T,f_0,f_+$。

首先记 $\mathrm{GR}^{\mathrm{triv}}\subset\mathrm{GR}$ 为遗忘映射 $\mathrm{GR}\to\Bun_G$ 下平凡 $G$-丛的原像。固定点 $x\in\PP^1\setminus\{0,\infty\}$，取 $t_x:=1-\tfrac tx$ 作为 $x$ 处的局部参数。于是 $t_x=1$ 对应 $t=0$，$t_x=\infty$ 对应 $t=\infty$。ind-概形 $\mathrm{Gr}^{\mathrm{triv}}_x$ 分类同构类对 $(E,\phi)$，其中 $E$ 是 $\PP^1$ 上的平凡 $G$-丛，$\phi:E|_{\PP^1\setminus\{x\}}\xrightarrow{\sim}G\times(\PP^1\setminus\{x\})$ 是同构。我们可以刚性化这个模问题，使 $\mathrm{Gr}^{\mathrm{triv}}_x$ 分类 $\PP^1\setminus\{x\}$ 上平凡 $G$-丛的、在 $\infty$ 处为恒等的自同构。故可等同
\begin{equation}
\mathrm{Gr}^{\mathrm{triv}}_x\xrightarrow{\ \sim\ }G[t_x^{-1}]_1:=\ker\big(G[t_x^{-1}]\xrightarrow{\mathrm{ev}(t_x=\infty)}G\big),
\end{equation}
其中 $G[t_x^{-1}]$ 是其 $R$-点为 $G(R[t_x^{-1}])$ 的 ind-概形。

\begin{remark}
有一族膨胀作用 $\mathbb{G}_m^{\mathrm{rot}}$ 作用在 $\PP^1\setminus\{0,\infty\}$ 上，它扩充到 $\mathrm{GR}$：$\lambda\in\mathbb{G}_m^{\mathrm{rot}}$ 把 $(x,E,\phi)$ 送到 $(\lambda x,\lambda^{-1,*}E,\lambda^{-1,*}\phi)$，并稳定 $\mathrm{GR}^{\mathrm{triv}}$。注意局部坐标 $t_x$ 在 $t$ 与 $x$ 的同时膨胀下不变。
\end{remark}

设 $\mathrm{Gr}^{\mathrm{triv}}\subset\mathrm{Gr}=G((\tau))/G[[\tau]]$ 是仿射 Grassmannian 的一个抽象拷贝，不涉及 $\PP^1\setminus\{0,\infty\}$ 上的任何点。膨胀作用连同 (5.9) 给出族 $\mathrm{GR}^{\mathrm{triv}}$ 在 $\PP^1\setminus\{0,\infty\}$ 上的平凡化：
\begin{equation}
\mathrm{GR}^{\mathrm{triv}}\cong(\PP^1\setminus\{0,\infty\})\times G[\tau^{-1}]_1\cong(\PP^1\setminus\{0,\infty\})\times\mathrm{Gr}^{\mathrm{triv}},\quad
(x,g(t_x^{-1}))\mapsto(x,g(\tau^{-1})).
\end{equation}
按定义，$\Bun_{G(1,2)}$ 中的大胞腔由 $T\times I(1)/I(2)$ 作用在平凡丛上得到。记 $\mathrm{ev}_0:G[t_x^{-1}]\to G[[t]]$ 与 $\mathrm{ev}_\infty:G[t_x^{-1}]\to G[[t^{-1}]]$ 为在 $0$ 与 $\infty$ 处的展开，$\mathrm{ev}_0,\mathrm{ev}_\infty:G[t_x^{-1}]\to G$ 为在 $t=0$ 与 $t=\infty$ 处的取值。

$G$-丛 $(E,\phi)\in\mathrm{Gr}_x$ 上的 $G(1,2)$-水平结构由平凡化 $\phi$ 得到。对 $g\in G[t_x^{-1}]$，复合 $E|_{\PP^1\setminus\{x\}}\xrightarrow{\phi}G\times\PP^1\setminus\{x\}\xrightarrow{g}G\times\PP^1\setminus\{x\}$ 扩充为同构 $E\to G\times\PP^1$。故 $g$ 定义了平凡丛上由 $\mathrm{ev}_0(g),\mathrm{ev}_\infty(g)$ 给出的水平结构。有交换图
\[
\begin{tikzcd}[column sep=large, row sep=large]
G[t_x^{-1}] \arrow[r, "{(\mathrm{ev}_0,\mathrm{ev}_\infty)}"] \arrow[dr, "{g\mapsto(\mathrm{ev}_0(g),\mathrm{ev}_\infty(g))}"'] & G\backslash\big(G[[t]]\times G[[t^{-1}]]\big) \arrow[d] \\
& \Bun_{G(1,2)}.
\end{tikzcd}
\]
这里 $G=\operatorname{Aut}(G\times\PP^1)$ 对角地作用在 $G[[t]]\times G[[t^{-1}]]$ 上。

于是我们发现 $g\in G[t_x^{-1}]$ 落在 $\mathrm{Gr}^\circ_x$ 中，当且仅当 $(\mathrm{ev}_0(g),\mathrm{ev}_\infty(g))\in G\backslash G(B^{\mathrm{opp}}\times U)$。利用等同 (5.9) 与 (5.11)，得到
\begin{equation}
\begin{aligned}
\mathrm{Gr}^\circ_x&\xrightarrow{\ \sim\ }\big\{g(t_x^{-1})\in G[t_x^{-1}]_1\ \big|\ \mathrm{ev}_0(g)\in UB^{\mathrm{opp}}\big\},\\
\mathrm{GR}^\circ&\xrightarrow{\ \sim\ }(\PP^1\setminus\{0,\infty\})\times\big\{g(\tau^{-1})\in G[\tau^{-1}]_1\ \big|\ g(1)\in UB^{\mathrm{opp}}\big\}
\end{aligned}
\end{equation}
包含在 $\mathrm{GR}^{\mathrm{triv}}=(\PP^1\setminus\{0,\infty\})\times G[\tau^{-1}]_1$ 中。

\begin{lemma}
对 (5.12) 参数化下的 $(x,g(\tau^{-1}))\in\mathrm{GR}^\circ$，写 $g(1)=ub^{\mathrm{opp}}$，其中 $u\in U$，$b^{\mathrm{opp}}\in B^{\mathrm{opp}}$。则
\[
\begin{aligned}
f_T(x,g)&=b^{\mathrm{opp}}\bmod U^{\mathrm{opp}}\in T;\qquad
f_+(x,g)=u^{-1}\bmod[U,U]\in U/[U,U];\\
f_0(x,g)&=x\,a_{-\theta}(g)\in U_{-\theta}\cong\ckg_{-\theta},
\end{aligned}
\]
其中 $a_{-\theta}:G[\tau^{-1}]\to\ckg_{-\theta}$ 把 $g$ 送到切向量 $\frac{dg(\tau^{-1})}{d(\tau^{-1})}\big|_{\tau^{-1}=0}\in\ckg$ 的 $\ckg_{-\theta}$-分量。
\end{lemma}

\begin{proof}
$f_T,f_+$ 的公式由 (5.12) 中的描述推出。按定义，$f_0(x,g)$ 由用局部参数 $t^{-1}$ 在 $t=\infty$ 处展开 $g(t_x^{-1})$ 并取 $t^{-1}$ 系数的 $\ckg_{-\theta}$-分量得到。注意
\[
\frac{dg(t_x^{-1})}{d(t^{-1})}\Big|_{t^{-1}=0}=x\,\frac{dg(t_x^{-1})}{d(t_x^{-1})}\Big|_{t_x^{-1}=0}\in\ckg.
\]
此外，在等同 (5.11) 下参数 $t_x$ 对应 $\tau$，故 $f_0(x,g)=x\,a_{-\theta}(g)$。引理得证。
\end{proof}

设 $\mathrm{AS}_{U,-\varphi_+}$ 为 $\mathrm{AS}_\psi$ 沿 $U\to U/[U,U]\xrightarrow{-\varphi_+}\mathbb{G}_a$ 的拉回。设 $\mathrm{Kum}_{B^{\mathrm{opp}},\chi}$ 为 $\mathrm{Kum}_\chi$ 沿 $B^{\mathrm{opp}}\to T$ 的拉回。设 $j_{UB^{\mathrm{opp}}}:UB^{\mathrm{opp}}\hookrightarrow G$ 为包含，并记
\[
J=J_{-\varphi_+,\chi}:=j_{UB^{\mathrm{opp}},!}\big(\mathrm{AS}_{U,-\varphi_+}\boxtimes\mathrm{Kum}_{B^{\mathrm{opp}},\chi}\big)\in D^b_c(G,\mathbb{Q}_\ell).
\]

\begin{remark}
由 [4]，或本文引理 2.3 的论证，$J$ 有 clean 性：
\[
j_{UB^{\mathrm{opp}},!}\big(\mathrm{AS}_{U,-\varphi_+}\boxtimes\mathrm{Kum}_{B^{\mathrm{opp}},\chi}\big)\xrightarrow{\ \sim\ }j_{UB^{\mathrm{opp}},*}\big(\mathrm{AS}_{U,-\varphi_+}\boxtimes\mathrm{Kum}_{B^{\mathrm{opp}},\chi}\big).
\]
我们可以把 $J=J_{-\varphi_+,\chi}\in D^b_c(G,\mathbb{Q}_\ell)$ 视作定义 2.4 与注 2.5 中自守层 $A_{\varphi,\chi}$ 的有限域类比。\zhnote{原文此处引作 ``Definition 2.2''，但该编号在 §2 中已被注 2.2 占用，所指显然是定义 2.4 的自守层；按 2.4 译。[来源：原文 p.30 与 §2.2 的编号顺序]}
\end{remark}

有了这些记号并利用等同 (5.11)，公式 (5.8) 变为
\begin{equation}
\Kl^V_{\ckG}(\varphi,\chi)\xrightarrow{\ \sim\ }\pi^{\mathrm{triv}}_!\big(\mathrm{IC}^{\mathrm{triv}}_V\otimes f_0^*\mathrm{AS}_{\varphi_0}\otimes\mathrm{ev}^*_{\tau=1}J\big),
\end{equation}
其中 $\pi^{\mathrm{triv}}:\mathrm{GR}^{\mathrm{triv}}=\PP^1\setminus\{0,\infty\}\times G[\tau^{-1}]_1\to\PP^1\setminus\{0,\infty\}$ 是投影。以下常把 $\mathrm{ev}_{\tau=1}$ 简写作 $\mathrm{ev}$。

本节最终目标是计算 $V=V_{\theta^\vee}$ 与 $\ckg$ 时 $\Kl^V_{\ckG}(\varphi,\chi)$ 的 Euler 特征。这里先对一般 $V$ 作几步化归。由 (5.15)，需计算
\begin{equation}
\chi_c(\PP^1\setminus\{0,\infty\},\Kl^V_{\ckG}(\varphi,\chi))=\chi_c\big(\PP^1\setminus\{0,\infty\}\times\mathrm{Gr}^{\mathrm{triv}},\mathrm{IC}^{\mathrm{triv}}_V\otimes f_0^*\mathrm{AS}_{\varphi_0}\otimes\mathrm{ev}^*J\big).
\end{equation}
按 $a_{-\theta}$ 是否消没，把 $\mathrm{Gr}^{\mathrm{triv}}\cong G[\tau^{-1}]_1$ 分解为 $\mathrm{Gr}^{\mathrm{triv},a\neq 0}$ 与 $\mathrm{Gr}^{\mathrm{triv},a=0}$。在 $\PP^1\setminus\{0,\infty\}\times\mathrm{Gr}^{\mathrm{triv},a=0}$ 上，(5.16) 中的复形沿 $\PP^1\setminus\{0,\infty\}$ 常值，故 Euler 特征为 $0$。

另一方面，有变量替换同构
\[
\PP^1\setminus\{0,\infty\}\times\mathrm{Gr}^{\mathrm{triv},a\neq 0}\ni(x,g)\mapsto(xa_{-\theta}(g),g)\in\mathbb{G}_m\times\mathrm{Gr}^{\mathrm{triv},a\neq 0}.
\]
在此同构下，$f_0=xa_{-\theta}$ 变成到 $\mathbb{G}_m$-因子的投影，于是可应用 Künneth 公式。汇总这些步骤，得到
\begin{align}
\chi_c(\PP^1\setminus\{0,\infty\},\Kl^V_{\ckG}(\varphi,\chi))
&=\chi_c\big(\PP^1\setminus\{0,\infty\}\times\mathrm{Gr}^{\mathrm{triv},a\neq 0},\mathrm{IC}^{\mathrm{triv}}_V\otimes f_0^*\mathrm{AS}_{\varphi_0}\otimes\mathrm{ev}^*J\big)\notag\\
&\qquad=\chi_c(\mathbb{G}_m,\mathrm{AS}_{\varphi_0})\,\chi_c(\mathrm{Gr}^{\mathrm{triv},a\neq 0},\mathrm{IC}^{\mathrm{triv}}_V\otimes\mathrm{ev}^*J)\notag\\
&=-\chi_c(\mathrm{Gr}^{\mathrm{triv},a\neq 0},\mathrm{IC}^{\mathrm{triv}}_V\otimes\mathrm{ev}^*J)\notag\\
&=-\chi_c(\mathrm{Gr}^{\mathrm{triv}}\setminus\star,\mathrm{IC}^{\mathrm{triv}}_V\otimes\mathrm{ev}^*J)\notag\\
&\qquad+\chi_c(\mathrm{Gr}^{\mathrm{triv},a=0}\setminus\star,\mathrm{IC}^{\mathrm{triv}}_V\otimes\mathrm{ev}^*J).
\end{align}
最后一个等式是因为基点 $\star\in\mathrm{Gr}$ 属于 $\mathrm{Gr}^{\mathrm{triv},a=0}$。

\begin{lemma}
$\chi_c(\mathrm{Gr}^{\mathrm{triv}}\setminus\star,\mathrm{IC}^{\mathrm{triv}}_V\otimes\mathrm{ev}^*J)=0$。
\end{lemma}

\begin{proof}
由于 $\mathrm{IC}_V$ 沿层 $\mathrm{Gr}_\lambda$ 常值，只需证明对支配余权 $\lambda\neq 0$ 有 $\chi_c(\mathrm{Gr}^{\mathrm{triv}}_\lambda,\mathrm{ev}^*J)=0$。

记 $D^b_c(G,\mathbb{Q}_\ell)$ 上的卷积积为：对 $K_1,K_2\in D^b_c(G,\mathbb{Q}_\ell)$，
\[
K_1*_G K_2:=m_!(K_1\boxtimes K_2),
\]
其中 $m:G\times G\to G$ 是乘法。令 $K_\lambda:=\mathrm{ev}_!\mathbb{Q}_{\ell,\mathrm{Gr}^{\mathrm{triv}}_\lambda}\in D^b_c(G)$。则 $R\Gamma_c(\mathrm{Gr}^{\mathrm{triv}}_\lambda,\mathrm{ev}^*J)$ 是 $\mathrm{AS}_{U,\varphi_+}*_G K_\lambda*_G\mathrm{Kum}_{B^{\mathrm{opp}},\chi^{-1}}$ 在 $e\in G$ 处的茎。由于 $\mathrm{ev}=\mathrm{ev}_{\tau=1}$ 是 $G$-等变的（在共轭下），$K_\lambda$ 带有自然的 $G$-等变结构。故
\[
K_\lambda*_G\mathrm{Kum}_{B^{\mathrm{opp}},\chi^{-1}}\cong\mathrm{Kum}_{B^{\mathrm{opp}},\chi^{-1}}*_G K_\lambda.
\]
特别地，在每个 Bruhat 层 $B^{\mathrm{opp}}wB^{\mathrm{opp}}$ 上，$K_\lambda*_G\mathrm{Kum}_{B^{\mathrm{opp}},\chi^{-1}}$ 形如 $S_w\otimes L_w$，其中 $L_w$ 是 $U^{\mathrm{opp}}wB^{\mathrm{opp}}$ 上的局部系统，而 $S_w$（$\mathbb{Q}_\ell$-向量空间复形）是 $K_\lambda*_G\mathrm{Kum}_{B^{\mathrm{opp}},\chi^{-1}}$ 在 $\dot w\in N_G(T)$（$w\in W$ 的代表元）处的茎。

因此 $\mathrm{AS}_{U,\varphi_+}*_G K_\lambda*_G\mathrm{Kum}_{B^{\mathrm{opp}},\chi^{-1}}$ 是诸 $S_w\otimes(\mathrm{AS}_{U,\varphi_+}*_G L_w)$（$w\in W$）的逐次扩张。为证明 $\mathrm{AS}_{U,\varphi_+}*_G K_\lambda*_G\mathrm{Kum}_{B^{\mathrm{opp}},\chi^{-1}}$ 的茎的 Euler 特征消没，只需证明对所有 $w$ 有 $\chi(S_w)=0$。

按定义，
\begin{equation}
\begin{split}
S_w&=\big(K_\lambda*_G\mathrm{Kum}_{B^{\mathrm{opp}},\chi^{-1}}\big)_{\dot w}=R\Gamma_c(wB^{\mathrm{opp}},K_\lambda\otimes\mathrm{Kum}_{B^{\mathrm{opp}},\chi})\\
&=R\Gamma_c\big(\mathrm{Gr}^{\mathrm{triv}}_\lambda\cap\mathrm{ev}^{-1}(wB^{\mathrm{opp}}),\mathrm{ev}^*\mathrm{Kum}_{B^{\mathrm{opp}},\chi}\big).
\end{split}
\end{equation}
$T$ 通过共轭作用在 $\mathrm{Gr}^{\mathrm{triv}}=G[\tau^{-1}]_1$ 上，并保持 $\mathrm{Gr}^{\mathrm{triv}}_\lambda\cap\mathrm{ev}^{-1}(wB^{\mathrm{opp}})$。$\mathrm{Gr}$ 上唯一的 $T$-不动点是 $\tau^\mu$（$\mu\in X_*(T)$），它们不属于任何 $\mathrm{Gr}^{\mathrm{triv}}_\lambda$（$\lambda\neq 0$）。此外，局部系统 $\mathrm{ev}^*\mathrm{Kum}_{B^{\mathrm{opp}},\chi}$ 在 $T$-共轭下是单值的：存在 $m\ge 1$（与 $p$ 互素）使 $\mathrm{ev}^*\mathrm{Kum}_{B^{\mathrm{opp}},\chi}$ 在 $T$-共轭的 $m$ 次幂下等变。由于 $T$-共轭的 $m$ 次幂在 $\mathrm{Gr}^{\mathrm{triv}}_\lambda\cap\mathrm{ev}^{-1}(wB^{\mathrm{opp}})$ 上仍无不动点，(5.19) 中的 Euler 特征为零。引理得证。
\end{proof}

\begin{corollary}
设 $X_*(T)^+\subset X_*(T)$ 为支配余权，$V(\lambda)\subset V$ 为权 $\lambda\in X_*(T)$ 的权空间。则
\[
\chi_c(\PP^1\setminus\{0,\infty\},\Kl^V_{\ckG}(\varphi,\chi))=\sum_{\lambda\in X_*(T)^+,\ \lambda\neq 0}\dim V(\lambda)\ \chi_c\big(\mathrm{Gr}^{\mathrm{triv},a=0}_\lambda,\mathrm{ev}^*J\big).
\]
\end{corollary}

\begin{proof}
由 [29, Theorem 6.1]，$\mathrm{IC}_V$ 沿 $\mathrm{Gr}_\lambda$ 的茎的 Euler 特征为 $\dim V(\lambda)$。因此
\[
\chi_c\big(\mathrm{Gr}^{\mathrm{triv},a=0}\setminus\star,\mathrm{IC}^{\mathrm{triv}}_V\otimes\mathrm{ev}^*J\big)=\sum_{\lambda\in X_*(T)^+,\ \lambda\neq 0}\dim V(\lambda)\ \chi_c\big(\mathrm{Gr}^{\mathrm{triv},a=0}_\lambda,\mathrm{ev}^*J\big).
\]
把它与 (5.17) 及引理 5.18 结合即得所求等式。
\end{proof}

% ============================================================
%  ch08 —— §5.3–5.5（原文 p.31–p.37）
% ============================================================
\routenote{p.31--p.37}

\subsection{拟极小 Schubert 簇}

对余权 $\lambda\in X_*(T)$，记 $P_\lambda\subset G$ 为由 $T$ 与满足 $\langle\alpha,\lambda\rangle\ge 0$ 的根空间 $U_\alpha$ 生成的抛物子群。对每个根 $\alpha$ 与 $i\in\mathbb{Z}$，记 $U_{\alpha,\ge i}\subset U_\alpha((\tau))$ 为其 $R$-点为 $U_\alpha(\tau^iR[[\tau]])$ 的子群，并令 $U_{\alpha,i}=U_{\alpha,\ge i}/U_{\alpha,\ge i+1}$。

现设 $\lambda\in X_*(T)^+$，考虑（开）Schubert 簇 $\mathrm{Gr}_\lambda\subset\mathrm{Gr}$。由 [31, Lemme 2.3 及其前的讨论]，有（固定 $\Phi^+$ 的一个全序）
\begin{equation}
\frac{G}{P_{-\lambda}}\times\prod_{\langle\alpha,\lambda\rangle\ge 2}U_{\alpha,\ge 1}\big/U_{\alpha,\ge\langle\alpha,\lambda\rangle}\xrightarrow{\ \sim\ }\mathrm{Gr}_\lambda,\qquad (g,u)\mapsto gu\tau^\lambda.
\end{equation}
$P_{-\lambda}$ 在 (5.21) 中乘积上的作用如下：先经 $P_{-\lambda}$ 在 $\prod_\alpha U_{\alpha,\ge 1}$ 上的伴随作用，再投影到满足 $\langle\alpha,\lambda\rangle\ge 2$ 的因子 $U_{\alpha,\ge 1}$。

现考虑特殊情形 $\lambda=\theta^\vee$，即支配短余根。我们需要回忆 [31] 中拟极小 Schubert 簇 $\mathrm{Gr}_{\theta^\vee}$ 的描述。记 $P_{-\theta}:=P_{-\theta^\vee}$。

\begin{lemma}
开子簇 $\mathrm{Gr}^{\mathrm{triv}}_{\theta^\vee}\subset\mathrm{Gr}_{\theta^\vee}$ 由单个 $G$-轨道组成。更精确地，态射
\[
\frac{G}{P_{-\theta}}\times U^\times_{-\theta,-1}\to G[\tau^{-1}]_1=\mathrm{Gr}^{\mathrm{triv}},\qquad (g,u_{-\theta}(c\tau^{-1}))\mapsto\mathrm{Ad}(g)u_{-\theta}(c\tau^{-1})
\]
给出到 $\mathrm{Gr}^{\mathrm{triv}}_{\theta^\vee}$ 的同构（这里 $P_{-\theta}$ 经伴随作用作用在 $U_{-\theta,-1}$ 上，而 $U^\times_{\alpha,i}=U_{\alpha,i}\setminus\{1\}$）。
\end{lemma}

\begin{proof}
这本质上是 [31, Lemme 7.2]。把 (5.21) 用于 $\lambda=\theta^\vee$，发现 (5.21) 中的乘积只剩一项 $U_{\theta,1}$（因为对任何根 $\alpha\neq\theta$ 有 $\langle\alpha,\theta^\vee\rangle\le 1$）。因此
\[
\mathrm{Gr}^{\mathrm{triv}}_{\theta^\vee}\cong\frac{G}{P_{-\theta}}\times U^{\mathrm{triv}}_{\theta,1}\cdot\tau^{\theta^\vee},
\]
其中 $U^{\mathrm{triv}}_{\theta,1}\subset U_{\theta,1}$ 是某个 $P_{-\theta}$-稳定的真开子集（$\mathrm{Gr}^{\mathrm{triv}}_{\theta^\vee}$ 不可能等于 $\mathrm{Gr}_{\theta^\vee}$，因为 $\tau^{\theta^\vee}\notin\mathrm{Gr}^{\mathrm{triv}}$）。由于 $P_{-\theta}$ 经伸缩作用在 $U_{\theta,1}$ 上，$U^{\mathrm{triv}}_{\theta,1}$ 必为 $U^\times_{\theta,1}$。故
\[
\mathrm{Gr}^{\mathrm{triv}}_{\theta^\vee}\cong\frac{G}{P_{-\theta}}\times U^\times_{\theta,1}\cdot\tau^{\theta^\vee}.
\]
另一方面，在由 $\theta$ 决定的 $SL_2$-子群中有如下计算：
\begin{equation}
\begin{pmatrix}1&c\tau\\&1\end{pmatrix}\begin{pmatrix}\tau&&\\&&\\&&\tau^{-1}\end{pmatrix}
=\begin{pmatrix}\tau&c\\&\tau^{-1}\end{pmatrix}
=\begin{pmatrix}1&c^{-1}\tau^{-1}\\&1\end{pmatrix}\begin{pmatrix}\tau&c\\&-c^{-1}\end{pmatrix}
\end{equation}
它表明：当 $c$ 可逆时，$u_\theta(c\tau)\tau^{\theta^\vee}\in\mathrm{Gr}$ 等于 $u_{-\theta}(c^{-1}\tau^{-1})\in U^\times_{-\theta,-1}\subset G[\tau^{-1}]_1$。引理得证。
\end{proof}

$G$ 的 Bruhat 分解 $G=\bigsqcup_{w\in W/W_\theta}UwP_{-\theta}$ 给出分解
\[
\frac{G}{P_{-\theta}}\times U^\times_{-\theta,-1}=\bigsqcup_{w\in W/W_\theta}\frac{UwP_{-\theta}}{P_{-\theta}}\times U^\times_{-\theta,-1}.
\]
层 $UwP_{-\theta}/P_{-\theta}\times U^\times_{-\theta,-1}$ 在 $G[\tau^{-1}]_1$ 中的像是 $\mathrm{Ad}(U)U^\times_{-w\theta,-1}$。由于 $w\mapsto-w\theta$ 给出 $W/W_\theta$ 与 $G$ 的长根集合之间的双射，可把上述分解改写为
\[
\mathrm{Gr}^{\mathrm{triv}}_{\theta^\vee}=\bigsqcup_{\beta\ \text{长根}}\mathrm{Ad}(U)U^\times_{\beta,-1}.
\]
对每个 $w\in W/W_\theta$，有同构（固定 $\Phi^+$ 的全序）
\[
\prod_{\alpha\in\Phi^+,\ \langle w^{-1}\alpha,\theta^\vee\rangle>0}U_\alpha\xrightarrow{\ \sim\ }UwP_{-\theta}/P_{-\theta}.
\]
因此，对 $\beta=-w\theta$，层 $\mathrm{Ad}(U)U^\times_{\beta,-1}$ 可写成
\begin{equation}
\begin{aligned}
\mathrm{Ad}(U)U^\times_{\beta,-1}&\xrightarrow{\ \sim\ }\prod_{\alpha\in\Phi^+,\ \langle\alpha,\beta^\vee\rangle<0}U_\alpha\times U^\times_{\beta,-1},\\
\mathrm{Ad}\Big(\prod u_\alpha\Big)u_\beta(c\tau^{-1})&\leftrightarrow\Big(\prod u_\alpha,\ u_\beta(c\tau^{-1})\Big).
\end{aligned}
\end{equation}

\begin{lemma}
引理 5.13 中的函数 $a_{-\theta}:\mathrm{Gr}^{\mathrm{triv}}_{\theta^\vee}\to\ckg_{-\theta}$ 在每个层 $\mathrm{Ad}(U)U^\times_{\beta,-1}\subset\mathrm{Gr}^{\mathrm{triv}}_{\theta^\vee}$ 上由
\[
a_{-\theta}\big(\mathrm{Ad}(u)u_\beta(c\tau^{-1})\big)=
\begin{cases}
c\,x_{-\theta},&\beta=-\theta,\\
0,&\text{否则}
\end{cases}
\]
给出，其中 $x_{-\theta}\in\ckg_{-\theta}$ 对应 $u_{-\theta}(1)\in U_{-\theta}$。特别地，
\begin{equation}
\mathrm{Gr}^{\mathrm{triv},a=0}_{\theta^\vee}=\bigsqcup_{\beta\ \text{长根},\ \beta\neq-\theta}\mathrm{Ad}(U)U^\times_{\beta,-1}.
\end{equation}
\end{lemma}

\begin{proof}
按定义，$a_{-\theta}$ 在 $\mathrm{Ad}(u)u_\beta(c\tau^{-1})$ 处的值为
\[
\frac{d}{d(\tau^{-1})}\mathrm{Ad}(u)u_\beta(c\tau^{-1})\Big|_{\tau^{-1}=0}\in\mathrm{Ad}(U)\ckg_\beta.
\]
若 $\beta\neq-\theta$，$\mathrm{Ad}(u)\ckg_\beta$ 只涉及 $\ge\beta$ 的根。若 $\beta=-\theta$，上述导数等于 $\mathrm{Ad}(u)c\,x_{-\theta}\in\ckg$，其 $\ckg_{-\theta}$-分量为 $cx_{-\theta}\in\ckg_{-\theta}$。
\end{proof}

\subsection{定理 4(1) 的证明}

把推论 5.20 用于 $V=V_{\theta^\vee}$（它只有一个非零支配权 $\theta^\vee$），得
\[
\chi_c(\PP^1\setminus\{0,\infty\},\Kl^{\theta^\vee}_{\ckG}(\varphi,\chi))=\chi_c\big(\mathrm{Gr}^{\mathrm{triv},a=0}_{\theta^\vee},\mathrm{ev}^*J\big).
\]
由引理 5.25 中的分解 (5.26)，只需对长根 $\beta\neq-\theta$ 计算 $\chi_c(\mathrm{Ad}(U)U^\times_{\beta,-1},\mathrm{ev}^*J)$。于是定理 4(1) 由下述断言推出。

\medskip\noindent\textbf{断言.} 设 $\beta\neq-\theta$ 是长根，则
\[
\chi_c\big(\mathrm{Ad}(U)U^\times_{\beta,-1},\mathrm{ev}^*J\big)=
\begin{cases}
-1,&\beta\ \text{是长单根},\\
0,&\text{否则}.
\end{cases}
\]

为证该断言，分三种情形。

\emph{情形 I：$\beta$ 为正但非单根。} 由于 $\beta>0$，有 $\mathrm{ev}(\mathrm{Ad}(U)U^\times_{\beta,-1})\subset U$。由于 $\beta$ 非单根，$\mathrm{ev}(\mathrm{Ad}(U)U^\times_{\beta,-1})$ 在 $U/[U,U]$ 中的像是平凡的，故 $\mathrm{ev}^*J$ 是 $\mathrm{Ad}(U)U^\times_{\beta,-1}$ 上的常层。由 (5.24)，$\mathrm{Ad}(U)U^\times_{\beta,-1}$ 有一个因子 $U^\times_{\beta,-1}\cong\mathbb{G}_m$，故
\[
\chi_c\big(\mathrm{Ad}(U)U^\times_{\beta,-1},\mathrm{ev}^*J\big)=\chi_c\big(\mathrm{Ad}(U)U^\times_{\beta,-1}\big)=0.
\]

\emph{情形 II：$\beta=\alpha_i$ 是长单根。} 对 $u\in U$，$\mathrm{ev}(\mathrm{Ad}(u)u_\beta(c\tau^{-1}))$ 在 $U/[U,U]$ 中的像是 $u_{\alpha_i}(c)$。按 (5.24) 的坐标，$\mathrm{ev}^*J$ 是从 $U^\times_{\beta,-1}\cong U^\times_{\alpha_i}$-因子拉回的 $\mathrm{AS}_{\varphi_i}$。故
\[
\chi_c\big(\mathrm{Ad}(U)U^\times_{\beta,-1},\mathrm{ev}^*J\big)=\chi_c\big(U^\times_{\alpha_i},\mathrm{AS}_{\varphi_i}\big)=-1.
\]

\emph{情形 III：$\beta$ 为负且 $\beta\neq-\theta$。} 由于 $\beta\neq-\theta$，存在单根 $\alpha_i$ 使 $\beta-\alpha_i$ 仍是根。则 $\alpha=-\beta+\alpha_i$ 是正根。又因 $\alpha_i\neq\beta$ 蕴涵 $\langle\alpha_i,\beta^\vee\rangle\le 1$，故
\[
\langle\alpha,\beta^\vee\rangle=\langle-\beta+\alpha_i,\beta^\vee\rangle=-2+\langle\alpha_i,\beta^\vee\rangle<0.
\]
于是 $U_\alpha$ 出现在分解 (5.24) 中。

利用 (5.24)，把 $\mathrm{Ad}(U)U^\times_{\beta,-1}$ 中的元素写成 $\mathrm{Ad}(u)u_\beta(c_\beta\tau^{-1})$，其中 $u=u_\alpha u_{\underline\alpha}$，$u_\alpha\in U_\alpha$ 而 $u_{\underline\alpha}\in A:=\prod_{\alpha'>0,\ \langle\alpha',\beta^\vee\rangle<0,\ \alpha'\neq\alpha}U_{\alpha'}$。注意
\begin{equation}
\mathrm{ev}\big(\mathrm{Ad}(u)u_\beta(c_\beta\tau^{-1})\big)=u\,u_\beta(c_\beta)\,u^{-1}=u_\alpha[u_{\underline\alpha}(c_\beta),u_\beta(c_\beta)]\,u_\beta(c_\beta)\,u_{\underline\alpha}^{-1}.
\end{equation}
由于 $\alpha\neq\pm\beta$，可应用 Chevalley 交换子关系 [9, p.36,(4)]，断言 $[u_{\underline\alpha}(c_\beta),u_\beta(c_\beta)]$ 是根群 $U_{i\underline\alpha+j\beta}$（$i,j\in\mathbb{Z}_{>0}$）中元素之积。我们的假设 (i) $\alpha+\beta$ 是单根、(ii) $\beta$ 是长根蕴涵：任何这样的 $i\underline\alpha+j\beta$（若为根）都是正根，且此形式中唯一的单根是 $\alpha+\beta$。因此 $[u_{\underline\alpha}(c_{\underline\alpha}),u_\beta(c_\beta)]\in U$，且其在 $U/[U,U]$ 中的像是 $u_{\alpha_i}(\epsilon c_{\underline\alpha}c_\beta)$，其中 $\epsilon=\pm 1$。

由 (5.27)，$\mathrm{Ad}(u)u_\beta(c_\beta)\in UB^{\mathrm{opp}}$ 当且仅当 $u_\beta(c_\beta)u_{\underline\alpha}^{-1}\in UB^{\mathrm{opp}}$。此外，当 $\mathrm{Ad}(u)u_\beta(c_\beta)\in UB^{\mathrm{opp}}$ 时，它在 $U/[U,U]$ 中的像是 $u_{\alpha_i}(\epsilon c_{\underline\alpha}c_\beta)$ 乘以另一个只依赖于 $c_\beta$ 与 $u_{\underline\alpha}$ 的元素。

在分解 (5.24) 下，$\mathrm{Ad}(U)U^\times_{\beta,-1}\cong U_\alpha\times A$。设 $\mathrm{pr}_A:\mathrm{Ad}(U)U^\times_{\beta,-1}\to A$ 为投影。由上述讨论，$\mathrm{ev}^*J$ 限制在 $\mathrm{pr}_A$ 的每个纤维上同构于 $U_\alpha$ 上的 Artin--Schreier 层。因此 $\mathrm{pr}_{A,!}\mathrm{ev}^*J=0$，从而
\[
H^*_c\big(\mathrm{Ad}(U)U^\times_{\beta,-1},\mathrm{ev}^*J\big)=H^*_c\big(A,\mathrm{pr}_{A,!}\mathrm{ev}^*J\big)=0.
\]
断言得证，定理 4(1) 的证明完成。

\subsection{伴随 Schubert 簇}

本小节设 $G$ 非单链。始终假设 $\operatorname{char}(k)$ 对 $G$ 良好，即当 $G$ 型为 $B_n,C_n$ 时 $\operatorname{char}(k)>2$，当 $G$ 型为 $F_4,G_2$ 时 $\operatorname{char}(k)>3$。

设 $\gamma$ 为 $G$ 的短支配根，$\gamma^\vee$ 为相应的长余根。完全类似 $P_{-\theta},W_\theta$ 等来定义 $P_{-\gamma},W_\gamma$ 等。按 $\langle\alpha,\gamma^\vee\rangle$（$\alpha\in\Phi$）的可能取值，有两种情形：
\begin{itemize}[leftmargin=1.6em, itemsep=.15em]
\item 型 $(B_n,C_n,F_4)$：$|\langle\alpha,\gamma^\vee\rangle|=0,1,2$；
\item 型 $(G_2)$：$|\langle\alpha,\gamma^\vee\rangle|=0,1,2,3$。
\end{itemize}
记 $\Phi^\gamma_n:=\{\alpha\in\Phi\mid\langle\alpha,\gamma^\vee\rangle=n\}$。

\begin{lemma}
设 $G$ 型为 $B_n,C_n$ 或 $F_4$，则 $\mathrm{Gr}^{\mathrm{triv}}_{\gamma^\vee}$ 由单个 $G$-轨道组成。更精确地，令
\[
V_{-\gamma}:=\prod_{\langle\alpha,\gamma^\vee\rangle=2}U_{-\alpha,-1},
\]
并把 $V_{-\gamma}$ 与其 Lie 代数等同。$P_{-\gamma}$ 在 $V_{-\gamma}$ 上的伴随作用（经其 Levi 因子 $L_\gamma$ 分解）稳定由 $V_{-\gamma}$ 上一个非退化二次型 $q_{-\gamma}$ 定义的唯一四元超曲面 $Q_{-\gamma}$，使 $L_\gamma$ 的作用经 $\mathrm{GO}(V_{-\gamma},q_{-\gamma})$ 分解\footnote{$\mathrm{GO}(V_{-\gamma},q_{-\gamma})$ 由 $V_{-\gamma}$ 的把 $q_{-\gamma}$ 保持到相差一个纯量的可逆线性自同构组成。}，并以 $V_{-\gamma}-Q_{-\gamma}$ 为稠密轨道。自然态射
\begin{equation}
\frac{G}{P_{-\gamma}}\times(V_{-\gamma}-Q_{-\gamma})\ni(g,v)\mapsto\mathrm{Ad}(g)v\in G[\tau^{-1}]_1
\end{equation}
给出同构 $\frac{G}{P_{-\gamma}}\times(V_{-\gamma}-Q_{-\gamma})\cong\mathrm{Gr}^{\mathrm{triv}}_{\gamma^\vee}$。
\end{lemma}

\begin{proof}
当 $G$ 型为 $B,C$ 或 $F_4$ 时，(5.21) 中的乘积化为 $\Lambda=\prod_{\langle\alpha,\gamma^\vee\rangle=2}U_{\alpha,1}$（一个交换幂幺群），可与其 Lie 代数等同。由于 $\mathrm{Gr}^{\mathrm{triv}}_{\gamma^\vee}$ 在 $G$-共轭下稳定，它形如 $\frac{G}{P_{-\gamma}}\times\Lambda^{\mathrm{triv}}$，其中 $\Lambda^{\mathrm{triv}}\subset\Lambda$ 是某个 $P_{-\gamma}$-稳定开子集。$P_{-\gamma}$ 在 $\Lambda$ 上的作用经 Levi 商 $L_\gamma$ 分解。

向量空间 $\Lambda$ 带双线性型
\begin{equation}
(x,y)_\Lambda:=(x,\mathrm{Ad}(w_\gamma)y)_{\ckg},
\end{equation}
其中 $w_\gamma$ 是 $\begin{pmatrix}0&1\\-1&0\end{pmatrix}$ 在对应根 $\gamma$ 的同态 $SL_2\to G$ 下的像，$(\cdot,\cdot)_{\ckg}$ 是 $\ckg$ 上 $\mathrm{Ad}(G)$-不变的非退化对称双线性配对（当 $\operatorname{char}(k)$ 良好时存在，见 [8, §1.16]）。容易验证 $(\cdot,\cdot)_\Lambda$ 是非退化对称双线性配对，且 $L_\gamma$ 在 $\Lambda$ 上的作用把该配对保持到相差纯量。记 $q_\Lambda$ 为 $(\cdot,\cdot)_\Lambda$ 相伴的二次型。

经快速的逐型验证，可知作用映射 $L_\gamma\to\mathrm{GO}(\Lambda,q_\Lambda)$ 是满的。因此 $\Lambda$ 含唯一的 $L_\gamma$-稳定不可约除子 $Q_\Lambda=\{q_\Lambda=0\}$，其补集是单个 $L_\gamma$-轨道。

另一方面，补集 $\mathrm{Gr}_{\le\gamma^\vee}-\mathrm{Gr}^{\mathrm{triv}}_{\le\gamma^\vee}$ 是代表行列式线丛类的丰沛除子，故余维为 $1$。由于 $\mathrm{Gr}_{\le\gamma^\vee}-\mathrm{Gr}_{\gamma^\vee}$ 的余维至少为 $2$，$\mathrm{Gr}_{\gamma^\vee}-\mathrm{Gr}^{\mathrm{triv}}_{\gamma^\vee}$ 在 $\mathrm{Gr}_{\gamma^\vee}$ 中余维也是 $1$。于是 $\Lambda-\Lambda^{\mathrm{triv}}$ 在 $\Lambda$ 中余维为 $1$，故必为该不可约除子 $Q_\Lambda$。这蕴涵 $\Lambda^{\mathrm{triv}}=\Lambda-Q_\Lambda$ 是单个 $P_{-\gamma}$-轨道，因此 $\mathrm{Gr}^{\mathrm{triv}}_{\gamma^\vee}$ 是单个 $G$-轨道。

与 (5.23) 相同的 $SL_2$-计算给出
\[
u_\gamma(\tau)\tau^{\gamma^\vee}=u_{-\gamma}(\tau^{-1})\in\mathrm{Gr}^{\mathrm{triv}}.
\]
因此 $\Lambda^{\mathrm{triv}}\tau^{\gamma^\vee}=\mathrm{Ad}(P_{-\gamma})u_\gamma(\tau)\tau^{\gamma^\vee}=\mathrm{Ad}(P_{-\gamma})u_{-\gamma}(\tau^{-1})\subset\mathrm{Gr}^{\mathrm{triv}}$。与上述类似的论证表明 $\mathrm{Ad}(P_{-\gamma})u_{-\gamma}(\tau^{-1})$ 是 $V_{-\gamma}=\prod_{\langle\alpha,\gamma^\vee\rangle=2}U_{-\alpha,-1}$ 中由某四元超曲面 $Q_{-\gamma}$ 的补集所定义的开子集。因此
\[
\mathrm{Gr}^{\mathrm{triv}}_{\gamma^\vee}=\frac{G}{P_{-\gamma}}\times(\Lambda-Q_\Lambda)\tau^{\gamma^\vee}=\frac{G}{P_{-\gamma}}\times(V_{-\gamma}-Q_{-\gamma})\subset G[\tau^{-1}]_1.
\]
\end{proof}

\begin{lemma}
态射 (5.29) 扩充为消解
\[
\nu:\frac{G}{P_{-\gamma}}\times V_{-\gamma}\to\mathrm{Gr}^{\mathrm{triv}}_{\le\gamma^\vee},
\]
它由引理 5.28 在 $\mathrm{Gr}^{\mathrm{triv}}_{\gamma^\vee}$ 上是同构。且
\begin{enumerate}[label=(\arabic*), leftmargin=1.9em, itemsep=.2em]
\item 纤维 $\nu^{-1}(\star)\cong G/P_{-\gamma}$；
\item 在 $\mathrm{Gr}^{\mathrm{triv}}_{\theta^\vee}$ 上的纤维同构于 $L_\theta/L_\theta\cap P_{-\gamma}$。
\end{enumerate}
\end{lemma}

\begin{proof}
由于 $\frac{G}{P_{-\gamma}}\times(V_{-\gamma}-Q_{-\gamma})$ 在 $\frac{G}{P_{-\gamma}}\times V_{-\gamma}$ 中稠密，而 $\mathrm{Gr}^{\mathrm{triv}}_{\le\gamma^\vee}$ 是 $\mathrm{Gr}^{\mathrm{triv}}_{\gamma^\vee}$ 在 $\mathrm{Gr}^{\mathrm{triv}}$ 中的闭包，态射 (5.29) 扩充为 $\nu$。由引理 5.28，$V_{-\gamma}$ 上有三个 $\mathrm{Ad}(P_{-\gamma})$（或 $\mathrm{Ad}(L_\gamma)$）-轨道：$V_{-\gamma}-Q_{-\gamma}$、$Q^\times_{-\gamma}:=Q_{-\gamma}-\{0\}$ 与 $\{0\}$，它们给出 $\frac{G}{P_{-\gamma}}\times V_{-\gamma}$ 的三个 $G$-轨道。

轨道 $\frac{G}{P_{-\gamma}}\times\{0\}$ 映到 $\mathrm{Gr}^{\mathrm{triv}}_0=\star$，这证明 (1)。

于是轨道 $\frac{G}{P_{-\gamma}}\times Q^\times_{-\gamma}$ 必映到 $\mathrm{Gr}^{\mathrm{triv}}_{\theta^\vee}$。$u_{-\theta}(\tau^{-1})$ 在 $\frac{G}{P_{-\gamma}}\times Q^\times_{-\gamma}$ 中与在 $\mathrm{Gr}^{\mathrm{triv}}_{\theta^\vee}=\frac{G}{P_{-\theta}}\times U^\times_{-\theta,-1}$ 中的 $G$-稳定子分别是 $P_{-\gamma}\cap P^1_{-\theta}$ 与 $P^1_{-\theta}$，其中 $P^1_{-\gamma}=\ker(\theta:P_{-\theta}\to\mathbb{G}_m)$。故纤维 $\nu^{-1}(u_{-\theta}(\tau^{-1}))$ 等于 $P^1_{-\theta}/P^1_{-\theta}\cap P_{-\gamma}$。容易验证包含 $L_\theta\hookrightarrow P_{-\theta}$ 与 $P^1_{-\theta}\hookrightarrow P_{-\theta}$ 诱导同构
\[
L_\theta/L_\theta\cap P_{-\gamma}\xrightarrow{\ \sim\ }P_{-\theta}/P_{-\theta}\cap P_{-\gamma}\xleftarrow{\ \sim\ }P^1_{-\theta}/P^1_{-\theta}\cap P_{-\gamma}.
\]
由引理 5.22，$\mathrm{Gr}^{\mathrm{triv}}_{\theta^\vee}$ 是单个 $G$-轨道，故 $\nu$ 在 $\mathrm{Gr}^{\mathrm{triv}}_{\theta^\vee}$ 上的所有纤维都同构于 $\nu^{-1}(u_{-\theta}(\tau^{-1}))\cong L_\theta/L_\theta\cap P_{-\gamma}$。
\end{proof}

$G$ 的 Bruhat 分解 $G=\bigsqcup_{W/W_\gamma}UwP_{-\gamma}$ 给出分解
\[
\frac{G}{P_{-\gamma}}\times V_{-\gamma}=\bigsqcup_{W/W_\gamma}\frac{UwP_{-\gamma}}{P_{-\gamma}}\times V_{-\gamma}.
\]
映射 $w\mapsto-w\gamma$ 给出 $W/W_\gamma$ 与 $G$ 的短根集合之间的双射。对 $\beta=-w\gamma$，令 $V_\beta=\mathrm{Ad}(w)V_{-\gamma}$。则上述分解可改写为
\begin{equation}
\frac{G}{P_{-\gamma}}\times V_{-\gamma}=\bigsqcup_{\beta\ \text{短根}}\mathrm{Ad}(U)V_\beta.
\end{equation}
与拟极小情形一样，可进一步把每个层写成
\begin{equation}
\begin{aligned}
\mathrm{Ad}(U)V_\beta&\xrightarrow{\ \sim\ }\prod_{\alpha\in\Phi^+,\ \langle\alpha,\beta^\vee\rangle<0}U_\alpha\times\prod_{\langle\beta',\beta^\vee\rangle=2}U_{\beta',-1},\\
\mathrm{Ad}\Big(\prod u_\alpha\Big)\prod u_{\beta'}(c_{\beta'}\tau^{-1})&\leftrightarrow\Big(\prod u_\alpha,\ \prod u_{\beta'}(c_{\beta'}\tau^{-1})\Big).
\end{aligned}
\end{equation}

\begin{remark}
\begin{enumerate}[label=(\arabic*), leftmargin=1.9em, itemsep=.2em]
\item 对每个短根 $\beta$，集合 $\Phi^\beta_2$（即出现在 $V_\beta$ 各因子中的那些根）按高度全序。为此只需证明：对不同的 $\beta',\beta''\in\Phi^\beta_2$，有 $\langle\rho^\vee,\beta'-\beta''\rangle\neq 0$。在引理 5.28 的证明中我们已注意到 $L_\gamma\twoheadrightarrow\mathrm{GO}(\Lambda,q_\Lambda)$ 是满的，这意味着对任意两个不同的根 $\alpha',\alpha''\in\Phi^\gamma_2$（即出现在 $\Lambda$ 各因子中的根），差 $\alpha'-\alpha''$ 是 $L_\gamma$ 的某个根的非零倍数。由于当 $\beta=w\gamma$ 时 $\Phi^\beta_2=w\Phi^\gamma_2$，差 $\beta'-\beta''$ 也是 $L_\beta$ 的某个根的非零倍数，故 $\langle\rho^\vee,\beta'-\beta''\rangle\neq 0$。

集合 $\Phi^\beta_2$ 带对合 $\beta'\mapsto 2\beta-\beta'$，其唯一不动点是 $\beta$。该对合反序。
\item 仍由引理 5.28 证明中的 $L_\gamma\twoheadrightarrow\mathrm{GO}(\Lambda,q_\Lambda)$，所有根 $\gamma'\in\Phi^\gamma_2\setminus\{\gamma\}$ 都被 $W_\gamma$ 置换。把 $\gamma$ 换成任一短根 $\beta$ 也有类似结论。特别地，$\Phi^\beta_2$ 中除 $\beta$ 自身外都是长根。
\end{enumerate}
\end{remark}

\begin{lemma}
对短根 $\beta$、$\prod u_{\beta'}(c_{\beta'}\tau^{-1})\in V_\beta$（乘积跑遍 $\Phi^\beta_2$）与 $u\in U$，有
\[
a_{-\theta}\Big(\mathrm{Ad}(u)\big(\prod u_{\beta'}(c_{\beta'}\tau^{-1})\big)\Big)=
\begin{cases}
c_{-\theta}\,x_{-\theta},&-\theta\in\Phi^\beta_2,\\
0,&\text{否则}.
\end{cases}
\]
\end{lemma}

\begin{proof}
证明与引理 5.25 类似。
\end{proof}

% ============================================================
%  ch09 —— §5.6（原文 p.37–p.40）
% ============================================================
\routenote{p.37--p.40}

\subsection{定理 4(2) 的证明}

本节假设 $G$ 型为 $B,C$ 或 $F_4$。对型 $A,D,E$ 的 $G$，定理 4 的陈述 (2) 与 (1) 相同；$G=G_2$ 的情形将在 §11 中证明。

把推论 5.20 用于 $V=\ckg$（它有两个非零支配权 $\theta^\vee$ 与 $\gamma^\vee$，重数均为 $1$），得
\[
\chi_c(\PP^1\setminus\{0,\infty\},\Kl^{\mathrm{Ad}}_{\ckG}(\varphi,\chi))=\chi_c\big(\mathrm{Gr}^{\mathrm{triv},a=0}_{\theta^\vee},\mathrm{ev}^*J\big)+\chi_c\big(\mathrm{Gr}^{\mathrm{triv},a=0}_{\gamma^\vee},\mathrm{ev}^*J\big).
\]
由 (1) 已知 $\chi_c(\mathrm{Gr}^{\mathrm{triv},a=0}_{\theta^\vee},\mathrm{ev}^*J)=r_\ell(G)$，只需证明
\[
\chi_c\big(\mathrm{Gr}^{\mathrm{triv},a=0}_{\gamma^\vee},\mathrm{ev}^*J\big)=r_s(G).
\]

对任意子集 $S\subset\Phi$，令
\[
V^S_\beta=\Big\{\prod u_{\beta''}(c_{\beta''}\tau^{-1})\in V_\beta\ \Big|\ c_{\beta''}=0\ \text{对}\ \beta''\notin S\Big\}.
\]
利用此定义与 $\Phi$ 上的标准偏序（由注 5.34(1)，它限制为 $\Phi^\beta_2$ 上的全序），$V^{>-\theta}_\beta$、$V^{>0}_\beta$ 等记号的含义是显然的。

利用引理 5.31、引理 5.35 与分解 (5.32)，消解 $\nu$ 限制到 $\mathrm{Gr}^{\mathrm{triv},a=0}_{\le\gamma^\vee}$ 上可读作
\[
\nu_{a=0}:\bigsqcup_{\beta\ \text{短根}}\mathrm{Ad}(U)\big(V^{>-\theta}_\beta\big)\to\mathrm{Gr}^{\mathrm{triv},a=0}_{\le\gamma^\vee}.
\]
由引理 5.31，$\mathrm{ev}^*J$ 在 $\nu_{a=0}$ 的目标上的 Euler 特征有来自 $\star$、$\mathrm{Gr}^{\mathrm{triv},a=0}_{\theta^\vee}$ 与 $\mathrm{Gr}^{\mathrm{triv},a=0}_{\gamma^\vee}$ 的贡献：
\begin{equation}
\begin{aligned}
\sum_{\beta\ \text{短根}}\chi_c\big(\mathrm{Ad}(U)V^{>-\theta}_\beta,\mathrm{ev}^*J\big)
&=\chi_c(G/P_{-\gamma})+\chi_c\big(\mathrm{Gr}^{\mathrm{triv},a=0}_{\theta^\vee},\mathrm{ev}^*J\big)\chi_c(L_\theta/L_\theta\cap P_{-\gamma})\\
&\quad+\chi_c\big(\mathrm{Gr}^{\mathrm{triv},a=0}_{\gamma^\vee},\mathrm{ev}^*J\big).
\end{aligned}
\end{equation}
Bruhat 分解给出
\begin{equation}
\chi_c(G/P_{-\gamma})=\#(W/W_\gamma)=\#\{\Phi\ \text{中的短根}\};
\end{equation}
\begin{equation}
\chi_c(L_\theta/L_\theta\cap P_{-\gamma})=\#(W_\theta/W_\theta\cap W_\gamma).
\end{equation}

\medskip\noindent\textbf{断言.} 对短根 $\beta$，
\[
\chi_c\big(\mathrm{Ad}(U)(V^{>-\theta}_\beta),\mathrm{ev}^*J\big)=
\begin{cases}
0,&\Phi^\beta_2\ \text{含单根},\\
1,&\text{否则}.
\end{cases}
\]

先承认该断言来完成证明。结合 (5.36)、(5.37)、(5.38) 与断言，得
\begin{equation}
\begin{aligned}
-\chi_c\big(\mathrm{Gr}^{\mathrm{triv},a=0}_{\gamma^\vee},\mathrm{ev}^*J\big)
&=\#\{\text{短根}\}-r_\ell(G)\#(W_\theta/W_\theta\cap W_\gamma)\\
&\quad-\#\{\text{短根}\ \beta\mid\Phi^\beta_2\ \text{不含单根}\}\\
&=\#\{\text{短根}\ \beta\mid\Phi^\beta_2\ \text{含单根}\}-r_\ell(G)\#(W_\theta/W_\theta\cap W_\gamma).
\end{aligned}
\end{equation}
由注 5.34(1)，$\Phi^\beta_2$ 按高度全序，故至多含一个单根。因此
\[
\#\{\text{短根}\ \beta\mid\Phi^\beta_2\ \text{含单根}\}=N_s+N_\ell,
\]
其中 $N_s$（相应地 $N_\ell$）是使 $\Phi^\beta_2$ 含短（相应地长）单根的短根 $\beta$ 的个数。
\begin{itemize}[leftmargin=1.6em, itemsep=.2em]
\item 若 $\Phi^\beta_2$ 含短单根，则由注 5.34(2)，该单根必为 $\beta$，因为 $\beta$ 是 $\Phi^\beta_2$ 中唯一的短根。故 $N_s=r_s(G)$。
\item 任一长单根都在 $\theta$ 的 $W$-轨道中，故 $N_\ell=r_\ell\cdot N_\theta$，其中 $N_\theta=\#\{\text{短根}\ \beta\mid\theta\in\Phi^\beta_2\}$。这样的短根 $\beta$ 在 $W_\theta$-轨道中（这由把注 5.34(2) 用于对偶根系得到），而 $\gamma$ 是其中之一，故 $N_\theta=\#(W_\theta/W_\theta\cap W_\gamma)$。于是 $N_\ell=r_\ell(G)\#(W_\theta/W_\theta\cap W_\gamma)$。
\end{itemize}
结合这些计算与 (5.39)，即得 $-\chi_c(\mathrm{Gr}^{\mathrm{triv},a=0}_{\gamma^\vee},\mathrm{ev}^*J)=r_s(G)$，从而证明定理 4(2)。

剩下证明断言。

\emph{情形 I：$\Phi^\beta_2$ 含单根 $\alpha_i$。} 对 $u\in U$ 与 $v=u_{\alpha_i}(c_{\alpha_i}\tau^{-1})v_{\alpha_i}\in V_\beta$（其中 $v_{\alpha_i}\in\prod_{\beta'\neq\alpha_i}U_{\beta',1}$），有
\[
\mathrm{ev}(\mathrm{Ad}(u)v)=\mathrm{Ad}(u)u_{\alpha_i}(c_{\alpha_i})\cdot\mathrm{Ad}(u)\mathrm{ev}(v_{\alpha_i}).
\]
这意味着 $\mathrm{Ad}(u)v\in UB^{\mathrm{opp}}$ 当且仅当 $\mathrm{Ad}(u)\mathrm{ev}(v_{\alpha_i})\in UB^{\mathrm{opp}}$；且当这发生时，它在 $U/[U,U]$ 中的像是 $u_{\alpha_i}(c_{\alpha_i})$ 乘以 $\mathrm{Ad}(u)\mathrm{ev}(v_{\alpha_i})$ 在 $U/[U,U]$ 中的像。

利用分解 (5.33)，可把 $\mathrm{Ad}(U)V^{>-\theta}_\beta$ 写成 $U_{\alpha_i,-1}\times A$，其中 $A$ 是某个仿射空间（其余因子之积），且 $\mathrm{ev}^*J$ 是外积 $\mathrm{ev}^*\mathrm{AS}_{\varphi_i}\boxtimes L$，$L$ 是 $A$ 上某局部系统。由 Künneth 公式，
\[
H^*_c\big(\mathrm{Ad}(U)V^{>-\theta}_\beta,\mathrm{ev}^*J\big)=H^*_c\big(U_{-\alpha_i,-1},\mathrm{ev}^*\mathrm{AS}_{\varphi_i}\big)\otimes H^*_c(A,L)=0.
\]

\emph{情形 II：$\Phi^\beta_2$ 不含任何单根。} 把向量空间 $V^{>-\theta}_\beta$ 分层为
\[
V^{>-\theta}_\beta=V^{>0}_\beta\ \sqcup\ \Big(\bigsqcup_{\beta'\in\Phi^\beta_2,\ -\theta<\beta'<0}\big(V^{\ge\beta'}_\beta-V^{>\beta'}_\beta\big)\Big).
\]
首先证明 $\chi_c(\mathrm{Ad}(U)V^{>0}_\beta,\mathrm{ev}^*J)=1$。事实上，由于 $\Phi^\beta_2\cap\Phi^+$ 不含单根，$\mathrm{ev}(\mathrm{Ad}(U)V^{>0}_\beta)\subset[U,U]$，故 $\mathrm{ev}^*J$ 是 $\mathrm{Ad}(U)V^{>0}_\beta$ 上的常层。又由 (5.33)，$\mathrm{Ad}(U)V^{>0}_\beta$ 是仿射空间，故得结论。

其次证明对每个 $\beta'\in\Phi^\beta_2$、$-\theta<\beta'<0$，有 $\chi_c(\mathrm{Ad}(U)(V^{\ge\beta'}_\beta-V^{>\beta'}_\beta),\mathrm{ev}^*J)=0$。由于 $\beta'\neq-\theta$，存在单根 $\alpha_i$ 使 $\beta'-\alpha_i$ 仍是根。则 $\alpha=\alpha_i-\beta'$ 是正根。按假设 $\alpha_i\notin\Phi^\beta_2$，故 $\langle\alpha,\beta^\vee\rangle=\langle\alpha_i,\beta^\vee\rangle-2<0$，即 $\alpha$ 出现在 (5.33) 的第一个乘积中。

\begin{lemma}
设 $a,b\in\mathbb{Z}_{>0}$，$\beta'\le\beta_1,\dots,\beta_b$ 是 $\Phi^\beta_2$ 中的根（不必互异），则 $a\alpha+\sum_{i=1}^b\beta_i\neq 0$。若 $a\alpha+\sum_{i=1}^b\beta_i$ 是根，则下列情形之一发生：
\begin{enumerate}[label=(\arabic*), leftmargin=1.9em, itemsep=.2em]
\item $a\alpha+\sum_{i=1}^b\beta_i$ 是负根。此时 $a=b-1$，$a\alpha+\sum_{i=1}^b\beta_i\in\Phi^\beta_2$，且（在 $\Phi^\beta_2$ 的全序中）大于任一 $\beta_i$；
\item $a\alpha+\sum_{i=1}^b\beta_i$ 为正但非单根；
\item $a=b=1$，$\beta_1=\beta'$，且 $a\alpha+\sum_{i=1}^b\beta_i=\alpha_i$。
\end{enumerate}
\end{lemma}

\begin{proof}
若 $a\ge b$。由于 $\alpha+\beta'>0$ 且 $\beta_i\ge\beta'$，$\alpha+\beta_i$ 必有正高度。因此
\[
a\alpha+\sum_{i=1}^b\beta_i=(a-b)\alpha+\sum_{i=1}^b(\alpha+\beta_i)>0.
\]
由于上式右端每一项都有正高度，当只有一个被加项时它必是单根，而这只能发生于情形 (3)。

若 $a<b$，则 $\langle a\alpha+\sum_{i=1}^b\beta_i,\beta^\vee\rangle\ge 2b-2a\ge 2$。因此若 $a\alpha+\sum_{i=1}^b\beta_i$ 是根，必有 $a=b-1$ 且 $a\alpha+\sum_{i=1}^b\beta_i\in\Phi^\beta_2$。由于 $\Phi^\beta_2$ 不含单根，得情形 (1) 或 (2)。

无论如何有 $a\alpha+\sum_{i=1}^b\beta_i\neq 0$。
\end{proof}

对 $u=\prod u_{\alpha'}(c_{\alpha'})$（乘积跑遍 $\alpha'\in\Phi^+$、$\langle\alpha,\beta^\vee\rangle<0$），记 $u_{\underline\alpha}=\prod_{\alpha'\neq\alpha}u_{\alpha'}(c_{\alpha'})$。类似地，对 $v\in V^{\ge\beta'}_\beta-V^{>\beta'}_\beta$，写 $v=u_{\beta'}(c_{\beta'}\tau^{-1})v_{\beta'}\prod u_{\beta''}(c_{\beta''}\tau^{-1})$（乘积跑遍 $\beta'<\beta''\in\Phi^\beta_2$，$c_{\beta'}$ 可逆）。则
\begin{equation}
\mathrm{ev}(\mathrm{Ad}(u)v)=u_\alpha\big[u_{\underline\alpha}(c_\alpha),u_{\beta'}(c_{\beta'})\big]u_{\beta'}(c_{\beta'})\Big(\prod_{\beta''>\beta'}\big[u_{\underline\alpha}(c_\alpha),u_{\beta''}(c_{\beta''})\big]u_{\beta''}(c_{\beta''})\Big)u_{\underline\alpha}^{-1}.
\end{equation}
为计算 (5.41)，用 Chevalley 交换子关系把每个 $[u_{\underline\alpha}(c_\alpha),u_{\beta''}(c_{\beta''})]$ 写成正根与负根因子之积，并设法把负根因子（由引理 5.40(1)，它们必出现在 $\Phi^\beta_2$ 中）移到右边。由 Chevalley 交换子关系，每一步都会产生形如 $a\alpha+\beta_1+\cdots+\beta_b$（$\beta_i\in\Phi^\beta_2$，$a,b>0$）的新因子。我们保留正因子，把负因子继续往右移。此过程中唯一需保证的是：在做交换子 $[U_{\alpha'},U_{\gamma'}]$ 时总有 $\alpha'$ 与 $\gamma'$ 线性无关，从而 Chevalley 交换子关系可用。事实上，过程中只出现 $\alpha'\in\Phi^+$、$\gamma'\in\Phi^-$ 且 $\alpha'+\gamma'$ 形如 $a\alpha+\beta_1+\cdots+\beta_b$（$\beta'\le\beta_i\in\Phi^\beta_2$，$a,b\in\mathbb{Z}_{>0}$）的情形。$\alpha'$ 与 $\gamma'$ 线性相关的唯一可能是 $\alpha'+\gamma'=0$，而这已被引理 5.40 排除。

过程结束时得
\begin{equation}
\mathrm{ev}(\mathrm{Ad}(u)v)=u_\alpha u_{\alpha_i}(\epsilon c_\alpha c_{\beta'})u_+\Big(\prod_{\beta''\ge\beta',\ \beta''\in\Phi^\beta_2}u_{\beta''}(c_{\beta''}+e c_{\beta''})\Big)u_{\underline\alpha}^{-1}.
\end{equation}
其中项 $u_{\alpha_i}(\epsilon c_\alpha c_{\beta'})$（$\epsilon=\pm 1$）来自 $[u_\alpha(c_\alpha),u_{\beta'}(c_{\beta'})]$。项 $u_+$ 是 $U_{a\alpha+\beta_1+\cdots+\beta_b}$（即 $a\alpha+\beta_1+\cdots+\beta_b\in\Phi^+$）中所有其他正因子之积。由引理 5.40(2)(3)，这些 $a\alpha+\beta_1+\cdots+\beta_b$ 从不是单根，故 $u_+\in[U,U]$。最后，额外系数 $e c_{\beta''}$ 来自 $U_{a\alpha+\beta_1+\cdots+\beta_b}$ 中的负因子，它是 $c_\alpha$ 与 $c_{\beta'''}$（$\beta'''\in\Phi^\beta_2$）的多项式函数。由引理 5.40(1)，$c_{\beta''}$ 只涉及那些满足 $\beta'''<\beta''$ 的 $c_{\beta'''}$。

因此可作变量替换
\begin{equation}
\begin{aligned}
\mathrm{Ad}(U)\big(V^{\ge\beta'}_\beta-V^{>\beta'}_\beta\big)&\xrightarrow{\ \sim\ }U_\alpha\times U^\times_{\beta',-1}\times\prod_{\langle\alpha',\beta^\vee\rangle<0,\ \alpha'\neq\alpha}U_{\alpha'}\times\prod_{\beta''\in\Phi^\beta_2,\ \beta''>\beta'}U_{\beta'',-1},\\
\mathrm{Ad}(u)(v)&\leftrightarrow\Big(u_\alpha(c_\alpha),\ u_{\beta'}(c_{\beta'}\tau^{-1}),\ u_{\underline\alpha},\ \prod u_{\beta''}(c_{\beta''}+ec_{\beta''})\Big).
\end{aligned}
\end{equation}
设 $A$ 为 (5.43) 中后三项之积，$\mathrm{pr}_A:\mathrm{Ad}(U)(V^{\ge\beta'}_\beta-V^{>\beta'}_\beta)\to A$ 为投影。由 (5.42)，$\mathrm{ev}^*J$ 限制在 $\mathrm{pr}_A$ 的纤维上同构于 $U_\alpha$ 上的 Artin--Schreier 层，故 $\mathrm{pr}_{A,!}\mathrm{ev}^*J=0$，从而
\[
H^*_c\big(\mathrm{Ad}(U)U^\times_{\beta,-1},\mathrm{ev}^*J\big)=H^*_c\big(A,\mathrm{pr}_{A,!}\mathrm{ev}^*J\big)=0.
\]
这完成了断言的情形 II。

% ============================================================
%  ch10 —— §6（原文 p.40–p.44）
% ============================================================
\routenote{p.40--p.44}

\section{整体单值性}

本节致力于证明定理 3。设 $G$ 是 $k$ 上分裂几乎单群。

\subsection{对加性特征的依赖}

回忆注 2.7(3)：$T\rtimes\operatorname{Aut}^\dagger(G)\times\mathbb{G}_m^{\mathrm{rot}}$ 作用在 $\Bun_{G(0,2)}$ 与 Hecke 对应 (2.6) 上。群 $T(k)\rtimes\operatorname{Aut}^\dagger(G)\times\mathbb{G}_m^{\mathrm{rot}}(k)$ 也作用在 $I^\infty(1)/I^\infty(2)$ 上，从而作用在一般加性特征构成的空间上。记 $S_\varphi$ 为 $\varphi$ 在 $T\rtimes\operatorname{Aut}^\dagger(G)\times\mathbb{G}_m^{\mathrm{rot}}$ 作用下的稳定子。这是 $k$ 上的有限群概形。

当 $G$ 为伴随型时，对每个 $\sigma\in\operatorname{Aut}^\dagger(G)$，存在唯一的 $(t,s)\in(T\times\mathbb{G}_m^{\mathrm{rot}})(k)$ 使 $(t,\sigma,s)$ 固定 $\varphi$。因此在此情形，投影 $S_\varphi\to\operatorname{Aut}^\dagger(G)\xrightarrow{\sim}\operatorname{Out}(G)$ 是同构（作为 $k$ 上的离散群）。一般地，$S_\varphi\to\operatorname{Aut}^\dagger(G)$ 是 $Z_G$-挠子，但在 $k$-点上未必满。

下述引理直接由几何 Hecke 算子的定义推出。

\begin{lemma}
定义 $\Kl_{\ckG}(\varphi)$ 的张量函子
\[
\Hk_\varphi:\mathcal S\ni\mathrm{IC}_V\mapsto\Hk_V(A_\varphi)|_{\star\times\PP^1\setminus\{0,\infty\}}\in\Loc(\PP^1\setminus\{0,\infty\})
\]
带自然的 $S_\varphi$-等变结构。这里 $S_\varphi(k)$ 经其在 $\operatorname{Aut}^\dagger(G)$ 中的像作用在 $\mathcal S$ 上，并经 $S_\varphi\to\mathbb{G}_m^{\mathrm{rot}}$ 对 $\PP^1\setminus\{0,\infty\}$ 的作用作用在 $\Loc(\PP^1\setminus\{0,\infty\})$ 上。
\end{lemma}

在引理 9.3(1) 中我们将验证：在等价 $\Rep(\ckG)\cong\mathcal S$ 下，$\operatorname{Aut}^\dagger(G)$ 在 $\mathcal S$ 上的作用与 $\operatorname{Aut}^\dagger(G)\cong\operatorname{Aut}^\dagger(\ckG)$（相对于对偶群上选定的钉住）的作用一致。

记 $S^1_\varphi(k)=\ker(S_\varphi(k)\to\mathbb{G}_m^{\mathrm{rot}}(k))$，并记 $S^{\mathrm{rot}}_\varphi(k)$ 为 $S_\varphi(k)$ 在 $\mathbb{G}_m^{\mathrm{rot}}(k)$ 中的像，它是有限循环群。我们希望用上述等变性推出 $\Kl_{\ckG}(\varphi)$ 的单值表示可取值于 $\ckG^{S^1_\varphi(k)}$。然而，由于该表示只在内部自同构意义下有定义，这需要额外的论证。

回忆我们已选定 $0$ 处形式去心圆盘 $\operatorname{Spec}K_0$ 上方的几何一般点 $\eta$。我们已把 $\Kl_{\ckG}(\varphi,1)$ 定义为从 Satake 范畴 $\mathcal S\cong\Rep(\ckG)$ 到 $\Loc(\PP^1\setminus\{0,\infty\})$ 的函子，故限制到 $\eta$ 给出张量函子
\begin{equation}
\omega_\varphi:\mathcal S\xrightarrow{\ \Kl_{\ckG}(\varphi)\ }\Loc(\PP^1\setminus\{0,\infty\})\xrightarrow{j_\eta^*}\Vec.
\end{equation}
由于这是 $\mathcal S$ 的纤维函子，可定义 ${}_\varphi\ckG=\operatorname{Aut}^\otimes(\omega_\varphi)$。当然 ${}_\varphi\ckG$ 同构于 $\ckG$，但并非典范同构。群 ${}_\varphi\ckG$ 的好处是：由 Tannaka 形式化，得到典范同态 $\operatorname{Aut}^\otimes(j_\eta^*)\to\operatorname{Aut}^\otimes(\omega_\varphi)$，其中 $\operatorname{Aut}^\otimes(j_\eta^*)$ 是基本群 $\pi_1(\PP^1\setminus\{0,\infty\},\eta)$ 的 pro-代数包络。因此得到同态
\begin{equation}
\phi:\pi_1(\PP^1\setminus\{0,\infty\},\eta)\to{}_\varphi\ckG(\mathbb{Q}_\ell).
\end{equation}
由引理 6.1 得到同态
\begin{equation}
S^1_\varphi(k)\to\operatorname{Aut}^\otimes(\mathcal S,\Hk_\varphi)\xrightarrow{(j_\eta^*)_*}\operatorname{Aut}^\otimes(\mathcal S,\omega_\varphi)=\operatorname{Aut}({}_\varphi\ckG).
\end{equation}
换言之，$S^1_\varphi$ 作用在 ${}_\varphi\ckG$ 上。我们将在引理 9.4 中证明 ${}_\varphi\ckG$ 还带自然的钉住 $\ddagger$，且上述作用保持该钉住。钉住 $\ddagger$ 与 $\dagger$ 定义典范同构 $\mathrm{can}:{}_\varphi\ckG\xrightarrow{\sim}\ckG$。利用此等同，得到：

\begin{corollary}
单值表示 $\phi$ 扩充为正合列之间的同态：
\begin{equation}
\begin{tikzcd}[column sep=large, row sep=large]
\pi_1(\PP^1\setminus\{0,\infty\},\eta) \arrow[r, "{\#S^{\mathrm{rot}}_\varphi(k)\text{-}\mathrm{th}}"] \arrow[d, "\phi"'] & \pi_1(\PP^1\setminus\{0,\infty\},\eta) \arrow[r] \arrow[d, "\tilde\phi"] & S^{\mathrm{rot}}_\varphi(k) \arrow[d, equal] \\
\ckG^{S^1_\varphi(k)} \arrow[r, hook] & \ckG^{S^1_\varphi(k)}\rtimes S^{\mathrm{rot}}_\varphi(k) \arrow[r, two heads] & S^{\mathrm{rot}}_\varphi(k)
\end{tikzcd}
\end{equation}
\end{corollary}

\begin{proof}
由 (6.4)，$S^1_\varphi(k)$ 在 ${}_\varphi\ckG$ 上的作用经 $S^1_\varphi(k)\to\operatorname{Aut}^\otimes(\mathcal S,\Hk_\varphi)$ 分解。由引理 9.2，单值表示 $\phi$ 因而经 $\pi_1(\PP^1\setminus\{0,\infty\},\eta)\to{}_\varphi\ckG^{S^1_\varphi}\subset{}_\varphi\ckG$ 分解，即
\[
\Hk_\varphi:\mathcal S\cong\Rep({}_\varphi\ckG)\xrightarrow{\mathrm{Res}}\Rep({}_\varphi\ckG^{S^1_\varphi(k)})\xrightarrow{\kappa}\Loc(\PP^1\setminus\{0,\infty\})
\]
对某个张量函子 $\kappa$ 成立。

我们用钉住同构“$\mathrm{can}$”等同 $\ckG\cong{}_\varphi\ckG$，故可把函子 $\kappa$ 视为 $\Rep(\ckG^{S^1_\varphi(k)})\to\Loc(\PP^1\setminus\{0,\infty\})$。于是 $S^{\mathrm{rot}}_\varphi(k)=S_\varphi(k)/S^1_\varphi(k)$ 经 $S_\varphi(k)\to\operatorname{Out}(G)\xrightarrow{\iota}\operatorname{Aut}^\dagger(\ckG)$ 作用在 $\ckG^{S^1_\varphi(k)}$ 上，从而作用在 $\kappa$ 的源与靶上。引理 6.1 蕴涵 $\kappa$ 带自然的 $S^{\mathrm{rot}}_\varphi(k)$-等变结构。对两个张量范畴取 $S^{\mathrm{rot}}_\varphi(k)$-不变量，得到函子 $\tilde\kappa$：
\[
\begin{tikzcd}[column sep=large, row sep=large]
\Rep(\ckG^{S^1_\varphi(k)}\rtimes S^{\mathrm{rot}}_\varphi(k)) \arrow[r, "\tilde\kappa"] \arrow[d, "\wr"'] & \Loc(\PP^1\setminus\{0,\infty\}/S^{\mathrm{rot}}_\varphi(k)) \arrow[d, "\wr"] \\
\Rep(\ckG^{S^1_\varphi(k)})^{S^{\mathrm{rot}}_\varphi(k)} \arrow[r, "(\kappa)^{S^{\mathrm{rot}}_\varphi}"] & \Loc(\PP^1\setminus\{0,\infty\})^{S^{\mathrm{rot}}_\varphi(k)}
\end{tikzcd}
\]
由于商映射 $\PP^1\setminus\{0,\infty\}\to\PP^1\setminus\{0,\infty\}/S^{\mathrm{rot}}_\varphi(k)$ 可等同于 $\PP^1\setminus\{0,\infty\}=\mathbb{G}_m$ 的 $\#S^{\mathrm{rot}}_\varphi(k)$ 次幂映射，即得图 (6.6)。
\end{proof}

\begin{remark}
把 $k$ 换成扩域 $k'$、把 $\PP^1\setminus\{0,\infty\}$ 换成 $\PP^1\setminus\{0,\infty\}\otimes_kk'$ 后，推论 6.5 仍成立。特别地，为获得几何单值性的信息，取 $k'=\bar k$。
\end{remark}

当 $\operatorname{Out}(\ckG)$ 非平凡时，有

\medskip
\begin{table}[htbp]
\centering
\caption{外自同构与 $\varphi$ 的稳定子}
\label{tab:2}
\small
\begin{tabular}{ccccc}
\toprule
$\ckG$ & $\operatorname{Aut}^\dagger(\ckG)\xrightarrow{\sim}\operatorname{Out}(\ckG)$ & $\ckG^{\operatorname{Aut}^\dagger(\ckG)}$ & $S^1_\varphi(k)$ & $S^{\mathrm{rot}}_\varphi(k)$\\
\midrule
$A_{2n-1}$（$n\ge 2$） & $\mathbb{Z}/2$ & $C_n$ & $\operatorname{Out}(\ckG)$ & $1$\\
$A_{2n}$ & $\mathbb{Z}/2$ & $B_n$ & $1$ & $\mathbb{Z}/2$\\
$D_4$ & $S_3$ & $G_2$ & $\operatorname{Out}(\ckG)$ & $1$\\
$D_n$（$n\ge 5$） & $\mathbb{Z}/2$ & $B_{n-1}$ & $\operatorname{Out}(\ckG)$ & $1$\\
$E_6$ & $\mathbb{Z}/2$ & $F_4$ & $\operatorname{Out}(\ckG)$ & $1$\\
\bottomrule
\end{tabular}
\end{table}
\medskip

\subsection{整体单值性的 Zariski 闭包}

设 $\ckG^{\mathrm{geo}}\subset\ckG$ 是几何单值表示 $\phi^{\mathrm{geo}}:\pi_1(\PP^1\setminus\{0,\infty\}\otimes_k\bar k,\eta)\to{}_\varphi\ckG\cong\ckG$ 的像的 Zariski 闭包。先证明 $\ckG^{\mathrm{geo}}$ 不太小。

\begin{proposition}[B.~Gross]
若 $\operatorname{char}(k)>2$ 且 $G$ 的秩至少为 $2$，则 $\ckG^{\mathrm{geo}}$ 不包含于 $\ckG$ 的任何主 $\mathrm{PGL}_2$ 中。
\end{proposition}

\begin{proof}
反设 $\ckG^{\mathrm{geo}}\subset\mathrm{PGL}_2\subset\ckG$，其中 $\mathrm{PGL}_2$ 含 $\ckG$ 的一个主幂幺元（$I^t_0$ 的像）。野惯性群 $I^+_\infty$ 的像必须非平凡，因为 $\Kl^{\theta^\vee}_{\ckG}(\varphi)$ 在 $\infty$ 处有非零 Swan 导子。

由于 $p=\operatorname{char}(k)>2$，$\phi(I^+_\infty)$ 落在极大环面 $\mathbb{G}_m\subset\mathrm{PGL}_2$ 中，且含 $\mu_p\subset\mathbb{G}_m$。由于 $\phi(I_\infty)$ 正规化 $\phi(I^+_\infty)$，它必包含在环面 $\mathbb{G}_m$ 的正规化子 $N(\mathbb{G}_m)\subset\mathrm{PGL}_2$ 中。

对 $\mathrm{PGL}_2$ 的任一不可约表示 $S^{2\ell}=\operatorname{Sym}^{2\ell}(\mathrm{Std})$（$\mathrm{Std}$ 是 $\mathrm{SL}_2$ 的二维表示），每对权空间 $S^{2\ell}(n)\oplus S^{2\ell}(-n)$（$0\le n\le\ell$）在 $N(\mathbb{G}_m)$ 下稳定，从而在 $I_\infty$ 下稳定。若权 $n$ 不整除 $p$，则 $S^{2\ell}(n)\oplus S^{2\ell}(-n)$ 的 Swan 导子至少为 $1$。因此
\[
\Swan_{I_\infty}(S^{2\ell})\ge\ell-[\ell/p].
\]
考虑 $I_\infty$ 在 $\ckG$ 的拟极小表示 $V_{\theta^\vee}$ 上的作用。由引理 10.1，$V_{\theta^\vee}$ 分解为主 $\mathrm{PGL}_2$ 的 $r_s(\ckG)$ 个不可约表示：$V_{\theta^\vee}=S^{2\ell_1}\oplus\cdots\oplus S^{2\ell_{r_s(\ckG)}}$。因此
\[
\Swan_{I_\infty}(V_{\theta^\vee})\ge\sum_{i=1}^{r_s(\ckG)}(\ell_i-[\ell_i/p])\ge\Big(1-\frac1p\Big)\frac{h\,r_s(\ckG)}{2}.
\]
这里用了 $\sum_i\ell_i=hr_s(\ckG)/2$（见引理 10.1 的证明）。故只要 $h>3$ 或 $p>3$，即得矛盾。即使在 $h=3$、$p=3$ 时，此时 $\ckG=\mathrm{PGL}_3$、$V_{\theta^\vee}=S^4\oplus S^2$，仍有 $\Swan(V_{\theta^\vee})\ge 2+1=3>r_s(\ckG)$，矛盾！
\end{proof}

\begin{proof}[定理 3 的证明]
由于对每个 $V\in\Rep(\ckG)$，$\Kl^V_{\ckG}(\varphi)$ 是权 $0$ 的纯层，故几何半单。由 Deligne [11, Corollaire 1.3.9]，$\ckG^{\mathrm{geo}}$ 的中性分支 $\ckG^\circ_{\mathrm{geo}}$ 是半单群。

\emph{第一步.} 设 $G$ 型不为 $A_1$、$A_{2n}$ 或 $B_3$。先确定 $\ckG^\circ_{\mathrm{geo}}$，等价地确定其 Lie 代数 $\ckg_{\mathrm{geo}}$。

一方面，由定理 1(2)，$\ckG^\circ_{\mathrm{geo}}$ 含主幂幺元，故含一个主 $\mathrm{PGL}_2$。命题 6.8 说 $\ckG^{\mathrm{geo}}$ 不能等于 $\mathrm{PGL}_2$。由于主 $\mathrm{PGL}_2$ 在 $\ckG$ 中是自正规化的，故 $\ckG^\circ_{\mathrm{geo}}$ 不能等于 $\mathrm{PGL}_2$，即 $\mathfrak{sl}_2\subsetneq\ckg_{\mathrm{geo}}$。

另一方面，由于 $G$ 型不为 $A_{2n}$，$S^1_\varphi(k)=\operatorname{Aut}^\dagger(\ckG)\xrightarrow{\sim}\operatorname{Out}(\ckG)$。在此情形，推论 6.5 蕴涵 $\ckG^{\mathrm{geo}}\subset\ckG^{\operatorname{Aut}^\dagger(\ckG)}$。

Dynkin 分类了 $\ckg$ 中含主 $\mathfrak{sl}_2$ 的所有 Lie 子代数。若 $\ckg\neq\mathfrak{so}_7$，则或者 $\ckg_{\mathrm{geo}}=\mathfrak{sl}_2$，或者 $\ckg_{\mathrm{geo}}$ 是 $\ckG$ 的某个钉住自同构的不动点代数。在我们的情形，若 $G$ 型不为 $A_{2n}$，已可断定 $\ckg_{\mathrm{geo}}=\ckg^{\operatorname{Aut}^\dagger(\ckG)}$。

\emph{第二步.} 设 $\ckG$ 型为 $B_3$ 且 $\operatorname{char}(k)>3$，断言 $\ckG^\circ_{\mathrm{geo}}=G_2$。为此只需证明 $\Kl^{\wedge^3V_7}_{\mathrm{SO}_7}(\varphi)$ 在 $\PP^1\setminus\{0,\infty\}\otimes_k\bar k$ 上有整体截面，其中 $V_7$ 是 $\mathrm{SO}_7$ 的 $7$ 维标准表示。设不然，则把 (5.4)–(5.6) 的长正合列中的 $\ckg$ 换成 $\wedge^3V_7$，将推出
\begin{equation}
\dim(\wedge^3V_7)^{I_0}+\dim(\wedge^3V_7)^{I_\infty}\le\dim H^1_c(\PP^1\setminus\{0,\infty\},\Kl^{\wedge^3V_7}_{\mathrm{SO}_7}(\varphi))=\Swan_{I_\infty}(\wedge^3V_7).
\end{equation}
对 $V_7$ 取标准基 $\{e_{-3},\dots,e_0,\dots,e_3\}$（二次型为 $e_0^2+e_1e_{-1}+e_2e_{-2}+e_3e_{-3}$）。$\wedge^3V_7$ 在主 $\mathfrak{sl}_2$ 下的 $0$-权空间由 $\{e_1\wedge e_{-2}\wedge e_{-3},\ e_{-1}\wedge e_2\wedge e_3,\ e_0\wedge e_i\wedge e_{-i}\ (i=1,2,3)\}$ 张成，故维数为 $5$。由定理 1(3)，$I_0$ 作为主幂幺元作用在 $V_7$ 上，故 $\dim(\wedge^3V_7)^{I_0}=5$。

由于 $\operatorname{char}(k)>3$ 不整除 $\ckG$ 的 Coxeter 数，推论 2.13 可用。由于 $I^+_\infty$ 在 $V_7$ 上作用的 breaks 为 $1/h=1/6$，$I^+_\infty$ 在 $\wedge^3V_7$ 的各非零权空间上的 breaks 至多为 $1/6$。因此 $\Swan_{I_\infty}(\wedge^3V_7)\le[\,\tfrac16\dim\wedge^3V_7\,]=[35/6]=5$。此外，由推论 2.13 中 $\phi(I^+_\infty)$ 的描述，$I^+_\infty$ 平凡地作用在 $\operatorname{Span}\{e_1\wedge e_2\wedge e_3,\ e_{-1}\wedge e_{-2}\wedge e_{-3}\}$ 上，而 Coxeter 置换 $e_1\to e_2\to e_3\to e_{-1}\to e_{-2}\to e_{-3}\to e_1$ 置换 $e_1\wedge e_2\wedge e_3$ 与 $e_{-1}\wedge e_{-2}\wedge e_{-3}$。故 $\dim(\wedge^3V_7)^{I_\infty}\ge 1$。于是
\[
\dim(\wedge^3V_7)^{I_0}+\dim(\wedge^3V_7)^{I_\infty}\ge 5+1>5\ge\Swan_{I_\infty}(\wedge^3V_7),
\]
与 (6.9) 矛盾。这证明当 $\ckG=\mathrm{SO}_7$ 时 $\ckG^{\mathrm{geo}}=G_2$。

\emph{第三步.} 当 $\ckG$ 型为 $A_{2n}$ 时，$\Kl_{\ckG}(\varphi)$ 来自经典 Kloosterman 层 $\mathrm{Kl}_n(\varphi)$，其整体单值性由 Katz 在 [25] 中处理，见 (0.2)。

\emph{第四步.} 剩下证明 $\ckG^{\mathrm{geo}}$ 连通。$A_n$ 的情形由 Katz 在上引文献中处理。$\ckG$ 型为 $E_6$ 时，容易验证 $\ckG^{\operatorname{Aut}^\dagger(\ckG)}$ 已是连通的。

一般地有 $\ckG^{\mathrm{geo}}\subset N_{\ckG}(\ckG^\circ_{\mathrm{geo}})=\ckG^\circ_{\mathrm{geo}}Z_{\ckG}$。考虑满同态
\[
\phi^{\mathrm{geo}}:\pi_1^{\mathrm{geo}}(\PP^1\setminus\{0,\infty\},\eta)\xrightarrow{\phi^{\mathrm{geo}}}\ckG^{\mathrm{geo}}\to\pi_0(\ckG^{\mathrm{geo}})=Z_{\ckG}/Z_{\ckG}\cap\ckG^\circ_{\mathrm{geo}}.
\]
假设 $\ckG$ 型不为 $A_n$ 或 $E_6$，则 $\pi_0(\ckG)$ 是 $2$-群。由于 $\operatorname{char}(k)>2$，$\phi^{\mathrm{geo}}$ 经驯商分解。另一方面，$I_0\cong\pi_1^{\mathrm{geo}}(\PP^1\setminus\{0,\infty\},\eta)$ 中的驯生成元必须映到 $\ckG^{\mathrm{geo}}$ 中的幂幺元，从而落在 $\ckG^\circ_{\mathrm{geo}}$ 内，因此映射 $\phi^{\mathrm{geo}}$ 平凡，即 $\ckG^{\mathrm{geo}}$ 连通。
\end{proof}

% ============================================================
%  ch11 —— §7（原文 p.44–p.46）
% ============================================================
\routenote{p.44--p.46}

\section{Kloosterman 层的函子性——猜想}

本节提出若干猜想，供进一步研究 Kloosterman 层之用。特别地，依据 Katz 与 Gabber 在 $GL_n$ 情形已知的刚性性质，表 1 中具有相同几何单值性的 Kloosterman 层在匹配加性特征 $\varphi$ 之后应当同构。我们还将对某些拟分裂群 $G$ 的 $\Kl^{\LG}(\varphi)$ 的局部与整体单值性给出猜想性描述。

\subsection{刚性}

\begin{conjecture}[Kloosterman 层的物理刚性]
设 $L$ 是 $\PP^1\setminus\{0,\infty\}\otimes_k\bar k$ 上的 $\ckG$-局部系统，它作为 $I_0$-与 $I_\infty$-表示与 $\Kl_{\ckG}(\varphi,\chi)$ 有相同的同构型。则 $L\cong\Kl_{\ckG}(\varphi,\chi)$（在 $\PP^1\setminus\{0,\infty\}\otimes_k\bar k$ 上）。
\end{conjecture}

我们期待更强的结论。

\begin{conjecture}
设 $L$ 是 $\PP^1\setminus\{0,\infty\}$ 上满足下述条件的 $\ckG$-局部系统：
\begin{itemize}[leftmargin=1.6em, itemsep=.2em]
\item $L$ 在 $\{0\}$ 处驯顺，$I^t_0$ 的一个拓扑生成元的像的半单部分共轭于 $\ckT[q-1]$ 中对应某个乘性特征 $\chi:T(k)\to\overline{\mathbb{Q}}_\ell^\times$ 的元素；
\item $\Swan_\infty(L^{\mathrm{Ad}})=r$，且 $(L^{\mathrm{Ad}})^{I_\infty}=0$。
\end{itemize}
则存在一般线性函数 $\varphi:I(1)/I(2)\to\mathbb{G}_a$ 使 $L\cong\Kl_{\ckG}(\varphi,\chi)$，相差一个非分歧扭转（由 $\operatorname{Gal}(\bar k/k)\to Z_{\ckG}$ 给出）。
\end{conjecture}

受上述猜想以及定理 1(2)、推论 2.13 与定理 3 中给出的 Kloosterman 层局部与整体单值性的计算启发，我们猜想 Kloosterman 层之间应有函子性关系。

\begin{conjecture}
设 $G,G'$ 是 $k$ 上分裂几乎单群，其对偶群 $\ckG\supset\ckG'$ 出现在表 1 的同一行。则对 $G$ 的每个一般线性函数 $\varphi$，存在 $G'$ 的一般线性函数 $\varphi'$，使得在 $\PP^1\setminus\{0,\infty\}\otimes_k\bar k$ 上，Kloosterman 层 $\Kl_{\ckG}(\varphi)$ 是 Kloosterman 层 $\Kl_{\ckG'}(\varphi')$ 的推出。
\end{conjecture}

\subsection{拟分裂群}

设 $G$ 是 $k$ 上分裂、几乎单且单连通的群。设 $\mathbb G$ 是 $\PP^1$ 上的拟分裂群概形，其在 $\PP^1\setminus\{0,\infty\}$ 上的限制由 §1.1 中的扭 $\sigma:\mu_N\hookrightarrow\operatorname{Aut}^\dagger(G)$ 给出。回忆 $N$ 假设与 $\operatorname{char}(k)$ 互素，且除 $G$ 型为 $D_4$ 时 $N=3$ 外，$N=2$。我们滥用记号，仍以 $\sigma$ 记 $\mu_N$ 的生成元在 $\operatorname{Aut}^\dagger(G)$ 中的像。利用引理 9.3(2) 中的同构把 $\operatorname{Aut}^\dagger(G)$ 与 $\operatorname{Aut}^\dagger(\ckG)$ 等同，记 $\langle\sigma\rangle\subset\operatorname{Aut}^\dagger(G)\cong\operatorname{Aut}^\dagger(\ckG)$ 为由 $\sigma$ 生成的子群。相应的 $L$-群可取为 $\LG=\ckG\rtimes\langle\sigma\rangle$。定理 1 中构造的 Kloosterman 层 $\Kl^{\LG}(\varphi)$ 给出单值表示
\[
\phi:\pi_1(\PP^1\setminus\{0,\infty\},\eta)\to\LG=\ckG\rtimes\langle\sigma\rangle.
\]
伴随表示 $\mathrm{Ad}:\ckG\to GL(\ckg)$ 可扩充为同态
\[
\mathrm{Ad}_\sigma:\LG\to GL(\ckg),
\]
使 Kloosterman 层 $\Kl^{\LG}(\varphi)$ 诱导局部系统 $\Kl^{\mathrm{Ad}}_{\LG}(\varphi)$。

对 $\Kl^{\LG}(\varphi)$ 的局部单值性，我们有如下预测。$\{0\}$ 处的驯单值应由元素 $(u,\sigma)\in\LG=\ckG\rtimes\langle\sigma\rangle$ 生成，其中 $u\in\ckG^\sigma$ 是主幂幺元。Swan 导子 $\Swan_\infty(\Kl^{\LG}(\varphi))=-\chi_c(\PP^1\setminus\{0,\infty\},\Kl^{\mathrm{Ad}}_{\LG}(\varphi))$ 应等于 $\ckG^\sigma$ 的中性分支的秩。此外，我们还期待推论 2.13 的类比对 $\Kl^{\LG}(\varphi)$ 成立，只需把 Coxeter 元换成 $\sigma$-扭 Coxeter 元 $\mathrm{Cox}_\sigma\in W\times\{\sigma\}\subset W\rtimes\langle\sigma\rangle$（见 [35] 与 [34, §5]），把 Coxeter 数换成 $\sigma$-扭 Coxeter 数 $h_\sigma$（即 $\mathrm{Cox}_\sigma$ 的阶）。

至于整体单值性，我们期待当 $\operatorname{char}(k)$ 不太小时，$\Kl^{\LG}(\varphi)$ 的整体几何单值表示 $\phi^{\mathrm{geo}}$ 的像 Zariski 稠密。

由表 1 观察到：对型不为 $A_{2n}$ 的单链分裂群 $G$，$\Kl_{\ckG}(\varphi)$ 的几何单值性的 Zariski 闭包 $\ckG^{\mathrm{geo}}$ 小于 $\ckG$。现取由阶为 $N$ 的非平凡 $\sigma\in\operatorname{Aut}^\dagger(G)$ 构造的 $G$ 的拟分裂形式 $\mathbb G$，则按我们的期待，$\widetilde{\PP^1\setminus\{0,\infty\}}$（$\PP^1\setminus\{0,\infty\}$ 的 $N$ 次 Kummer 覆盖，仍同构于 $\PP^1\setminus\{0,\infty\}$）上的 $\ckG$-局部系统 $[N]^*\Kl^{\LG}(\varphi)$ 的几何单值性在 $\ckG$ 中应 Zariski 稠密。这补偿了单链 $G$ 时 $\ckG^{\mathrm{geo}}$ 偏小的情况。

当 $G$ 型为 $A_{2n}$ 时，由推论 6.5，$\phi$ 扩充为
\[
\pi_1(\PP^1\setminus\{0,\infty\},\eta)\to\ckG\rtimes\mu_2,
\]
从而给出 $\PP^1\setminus\{0,\infty\}/\mu_2$ 上的 $\ckG\rtimes\mu_2$-局部系统 $\Kl_{\ckG}(\varphi)$。另一方面，设 $\mathbb G$ 是由非平凡 $\sigma\in\operatorname{Aut}^\dagger(G)$ 给出的拟分裂酉群，我们在 $\PP^1\setminus\{0,\infty\}$ 上有 Kloosterman 层 $\Kl^{\LG}(\varphi')$。在把 $\LG$ 与 $\ckG\rtimes\mu_2$ 等同、并把 $\varphi$ 与 $\varphi'$ 匹配之后，我们期待在 $\PP^1\setminus\{0,\infty\}\otimes_k\bar k\cong(\PP^1\setminus\{0,\infty\}/\mu_2)\otimes_k\bar k$ 上有 $\Kl^{\LG}(\varphi)\cong\Kl_{\ckG}(\varphi')$。

% ============================================================
%  ch12 —— §8（原文 p.46–p.48）
% ============================================================
\routenote{p.46--p.48}

\section{附录：$\PP^1$ 上 $G$-丛模空间几何的命题 1.1 的证明}

本附录给出命题 1.1 的证明。

\begin{proof}
前两条断言的证明用与 [23, Proposition 3] 相同的论证。先假设基域 $k$ 代数闭。

\emph{A.} 设 $G=T$ 是环面。此时 (2) 已在 [24, Lemma 16] 中证明。下面证明 (1)。由于 $T$ 交换，对任何 $T$-丛 $P$ 有 $H^1(\PP^1,P\times^T\mathrm{Lie}(T))=H^1(\PP^1,\mathrm{Lie}(T))$。我们断言该群消没。

按构造，$\mathrm{Lie}(T)\subset\mathrm{Lie}([N^*](\mathbb{G}_m\times\PP^1))$ 是直和项。又 $\mathrm{Lie}([N]^*(\mathbb{G}_m\times\PP^1))=[N]^*\mathcal O_{\PP^1}$，而 $H^1(\PP^1,[N]^*(\mathcal O_{\PP^1}))=H^1(\PP^1,\mathcal O_{\PP^1})=0$。故对任何 $P$ 有 $H^1(\PP^1,P\times^T\mathrm{Lie}(T))=0$，于是 $T$-丛没有非平凡变形。由于 $k$ 代数闭，$H^1(k(\PP^1),T)=0$，即所有 $T$-丛都一般平凡。因此映射 $\oplus_{x\in\PP^1}T(K_x)/T(\mathcal O_x)\to\Bun_T(k)$ 是满的。按构造，同构 $\pi_0(\Bun_T)\cong\pi_1(T)^{\pi_1(\mathbb{G}_m)}$ 由 Kottwitz 同态 $\pi_0(\mathrm{Gr}_{T,x})\to X_*(T)^{\operatorname{Gal}(K^{\mathrm{sep}}_x/K_x)}$ 诱导。故对任何 $x$，映射 $T(K_x)\to\Bun_T(k)$ 是满的。这证明了环面情形的 (1)。最后，(1) 与 (2) 蕴涵环面情形的 (4)。

\emph{B.} 对半单单连通群 $G$，陈述 (1) 与 (2) 已在 [24, Theorem 4] 中证明（同样假设 $k$ 代数闭）。

\emph{C.} 设 $G|_{\PP^1\setminus\{0,\infty\}}$ 的导群 $G^{\mathrm{der}}|_{\PP^1\setminus\{0,\infty\}}$ 单连通，并记 $D|_{\PP^1\setminus\{0,\infty\}}:=G/G^{\mathrm{der}}|_{\PP^1\setminus\{0,\infty\}}$。Haines 与 Rapoport 证明 $\PP^1$ 上有群概形的诱导正合列
\[
1\to G^{\mathrm{der}}\to G\to D\to 1.
\]
由于任何 $D$-丛在 $\PP^1\setminus\{x\}$ 上平凡，且对 $G^{\mathrm{der}}$ 的所有形式也有同样的结论，这蕴涵 $G$ 的 (1)。此外，由（上引文献）知 $G(k((s)))\to D(k((s)))$ 是满的。于是 $\pi_0(\Bun_G)\cong\pi_0(\Bun_D)$，这证明了 $G$ 的 (2)。

\emph{D.} 对一般群 $G$，可取 $z$-扩张
\[
1\to Z\to G'\to G\to 1,
\]
使 $Z$ 是诱导环面且 $G'^{\mathrm{der}}$ 单连通。由 Tsen 定理 $H^2(\PP^1,Z)=0$，故任何 $G$-丛都可提升为 $G'$-丛。由于已知 $G'$ 的 (1)，这蕴涵 $G$ 的 (1)。为证 (2)，利用在我们的假设下 $(X_*(Z))^{\pi_1(\PP^1\setminus\{0,\infty\})}$ 无挠，故如 [23] 那样有正合列
\[
0\to X_*(\check Z^{\pi_1(\PP^1\setminus\{0,\infty\})})\to X_*(Z(\ckG')^{\pi_1(\PP^1\setminus\{0,\infty\})})\to X_*(Z(\ckG)^{\pi_1(\PP^1\setminus\{0,\infty\})})\to 0.
\]

为证第三条断言，我们遵循 Ramanathan [33] 或 Faltings [14, Lemma 4] 的论证。

先固定支配权的记号。首先 $X_*(T)^{\pi_1(\mathbb{G}_m)}\otimes\mathbb{Q}=X_*(S)\otimes\mathbb{Q}$。此外，相对根 $\Phi(G,S)$ 张成 $S$ 的在 $G$ 的中心的限制下平凡的特征构成的子空间。记 $X_*(T)^{\pi_1(\mathbb{G}_m),+}\subset\pi_0(\Bun_T)$ 为使对任何正相对根 $a\in\Phi(G,S)^+$ 有 $a(\gamma)\ge 0$ 的那些 $\gamma\in X_*(T)^{\pi_1(\mathbb{G}_m)}$ 构成的子集。

首先证明
\[
G(k((s)))=\bigsqcup_{\lambda\in X_*(T)^{\pi_1(\mathbb{G}_m),+}}G(k[t])\lambda(s)G(k[[s]]).
\]
由 (1) 知每个 $G$-丛在任一点之外平凡。特别地，任何 $G$-丛可约化到 Borel 子群 $B$。设 $E$ 是 $G$-丛，并选其到 $B$ 的约化 $E_B$。

对任何特征 $\alpha:T\to\mathbb{G}_m$，记 $E_B(\alpha)$ 为 $\PP^1$ 上相伴的线丛。由于 $X_*(T)^{\pi_1(\mathbb{G}_m)}\otimes\mathbb{Q}\cong X_*(S)\otimes\mathbb{Q}$，次数 $\deg(E_B(\alpha))\in\mathbb{Q}$ 对 $\alpha\in X_*(S)$ 也有定义。我们断言：若对所有正单根 $a_i\in\Phi(G,S)^+$ 有 $\deg(E_B(a_i))\ge 0$，则丛 $E_B$ 可约化到 $T$。

为证此，记 $U$ 为 $B$ 的幂幺根，$E_T:=E_B/U$ 为诱导的 $T$-丛。为证 $E_B$ 由 $E_T$ 诱导，只需证明 $H^1(\PP^1,E_T\times^TU)=0$。群 $U$ 有滤链，其子商由根子群给出。考虑正相对根 $a\in\Phi(G,S)$。根子群 $U_a$ 是 $[N]^*(\oplus U_{\alpha'})$ 的直和项，其中和跑遍限制到 $S$ 上等于 $a$ 的那些根 $\alpha'\in\Phi(G,T)$。故 $U_a$ 是满足 $H^1(\PP^1,V)=0$ 的向量丛 $V$ 的直和项。

类似地，$E_T\times^TU_a$ 是 $E_T\times^T[N]^*(\oplus U_{\alpha'})$ 的直和项。由于 $[N]^*E_T$ 是 $\pi_1(\mathbb{G}_m)$-不变的，这蕴涵当 $\deg(E_B(a))\ge 0$ 时该丛的 $H^1$ 为 $0$。于是也得 $H^1(\PP^1,E_T\times^TU)=0$。

若约化 $E_B$ 对某个单根 $a_i$ 不满足 $\deg(E_B(a_i))\ge 0$，我们想修改 $E_B$。考虑由 $B$ 与 $U_{-a_i}$ 生成的抛物子群 $P_{a_i}\subset G$。根子群 $U_{\pm a_i}$ 定义了 $P$ 的一个子群 $L$，使 $L$ 的单连通覆盖或同构于 $[n]^*SL_2$，或同构于 $[n]^*SU_3$，其中 $n$ 整除 $N$，见 [7, §4.1.4]。半单商 $P^{ss}/Z(P^{ss})$ 同构于 $L^{\mathrm{ad}}$，且 $P/B\cong L/(L\cap B)$。我们断言结论对 $L^{\mathrm{ad}}$-丛成立，因为这些丛可用向量丛描述。若 $L^{\mathrm{ad}}=[n]^*\mathrm{PGL}_2$，则结论由秩 $2$ 向量丛的相应结论推出。酉群的情形类似，将在下面引理 8.1 中说明。

于是可找到新的到 $B$ 的约化 $E'_B$，使 $E_{P_{a_i}}$ 不变，但 $E'_B(a_i)\ge 0$。这意味着：对基本权 $\epsilon_k$（即由 $B$ 与所有 $j\neq k$ 的 $U_{a_j}$ 生成的极大抛物子群 $P_k$ 的伴随表示的行列式的倍数），当 $k\neq i$ 时 $\deg E_B(\epsilon_k)$ 不变，而 $\deg E'_B(\epsilon_i)$ 大于 $\deg E_B(\epsilon_i)$。

然而 $E'_B(\epsilon_i)$ 是 $\wedge^{\dim(\mathrm{Lie}(P_i))}E(\mathrm{Lie}(G))$ 的子丛，故所有这些线丛的次数有界，于是该过程必然终止于满足对所有单根 $a_i$ 有 $\deg(E_B(a_i))\ge 0$ 的 $B$-约化。这证明了 (3)。

下面推出 Birkhoff 分解 (4)。对常群的情形，这通常由分解 $G=B^-W_0B$ 与 (3) 推出。对一般 $\mathbb G$，其类比由 [7, Theorem 4.6.33] 提供：记 $\mathbb G_0:=\mathbb G^0_{\mathrm{red}}$ 为 $\mathbb G$ 在 $0$ 处纤维的约化商，$B_0\subset\mathbb G_0$ 为 $B$ 的像。记 $U_0$ 为 $B_0$ 的幂幺根，使 $\mathbb G_0=U_0W_0B_0$。所引结果表明 $B_0$ 在 $\mathbb G$ 中的原像正是 $I(0)$。按 $\mathbb G$ 的构造，$U_0(k)$ 的元素可提升到 $U(\PP^1)$。于是
\[
\mathbb G(k[[s]])=U_0W_0I(0).
\]
类似地，按 $\mathbb G$ 的构造，该分解蕴涵 $\mathbb G(\mathbb{A}^1)=I^-(0)W_0U_0$。对支配的 $t\in T(k((s)))$ 与 $b\in N_0$ 有 $t^{-1}bt\in I(0)$。于是利用 (3) 得
\[
\mathbb G(k((s)))=I^-(0)T(k((s)))\mathbb G(k[[s]]).
\]
现以同样方式分解 $\mathbb G(k[[s]])$。对任何 $t$ 可取 $w\in W_0$ 使 $wtw^{-1}$ 支配。再取 $B'_0\subset\mathbb G_0$ 为 $wB_0w^{-1}$，并写 $\mathbb G_0=B'_0W_0B_0$。则可用与前相同的论证推出 (4)（仍假设 $k$ 代数闭）。

下面推出 $k$ 为有限域的情形。先假设 $\mathbb G$ 在全分歧覆盖 $[N]:\PP^1\to\PP^1$ 上分裂。此时嵌入 $\widetilde W\to\mathrm{Gr}_{\mathbb G,x}$ 定义在 $k$ 上，故 $\Bun_{\mathbb G}$ 的所有几何点都定义在 $k$ 上。此外我们断言任何 $\mathbb G(0,0)$-丛的自同构群连通。一旦证明这一点，(1)、(3)、(4) 便由 Lang 定理在 $k$ 上成立。

有 $H^0(\PP^1,T)\hookrightarrow H^0(\PP^1,[N]^*(\mathbb{G}_m))=\mathbb{G}_m$。因此 $H^0(\PP^1,T)\cong T^{0,\mathrm{red}}$，它连通。又对任何 $w\in\widetilde W$，相应的 $k$ 上 $\mathbb G(0,0)$-丛的自同构群是 $I^-(0)\cap wI(0)w^{-1}$。该群以 $H^0(\PP^1,T)$ 为商，其核是仿射根对应的根子群之积，也是连通的。

一般情形由此经 Galois 下降得到，利用 [23, Remark 9]：Iwahori--Weyl 群可计算为 $k$ 的可分闭包上的 Iwahori--Weyl 群的 Galois 不变量。
\end{proof}

上述证明中用到了下述特殊情形。记 $SU_3$ 为覆盖 $[2]:\PP^1\to\PP^1$ 的拟分裂酉群。它可描述为 Hermite 型 $h(x_1,x_2,x_3)=x_1x_3^\sigma+x_2x_2^\sigma+x_3x_1^\sigma$ 的特殊酉群。记 $PSU_3$ 为相应的伴随群。

\begin{lemma}
任何 $PSU_3$-丛 $P$ 都有到 $B$ 的约化 $P_B$，使对正根 $\alpha$ 有 $\deg(P_B(\alpha))\ge 0$。
\end{lemma}

\begin{proof}
定义 $GU_3$ 为由 $SU_3$ 把其中心扩充为 $[2]^*\mathbb{G}_m$ 所得的群，于是有正合列 $1\to[2]^*\mathbb{G}_m\to GU_3\to PSU_3\to1$。同样，每个 $PSU_3$-丛都由某个 $GU_3$-丛诱导。这样的丛可视为覆盖 $\PP^1\xrightarrow{[2]}\PP^1$ 上的秩 $3$ 向量丛 $E$ 带一个取值于形如 $[2]^*L$ 的线丛的 Hermite 型。此时，给出到 $B$ 的约化只需给出各向同性的线子丛 $E_1\to E$，因为它定义旗 $E_1\subset E_1^\perp\subset E$。若 $E$ 非半稳定，则 $E$ 的典范子丛给出正次数的各向同性子丛，故此情形约化存在。若 $E$ 半稳定，则 Hermite 型给出整体同构 $E\xrightarrow{\sim}\sigma^*E^\vee\otimes[2]^*L$。但这样的同构必为常值，故可找到次数为 $\frac{\deg(E)}{3}$ 的各向同性子丛。
\end{proof}

% ============================================================
%  ch13 —— §9（原文 p.48–p.52）
% ============================================================
\routenote{p.48--p.52}

\section{附录：几何 Satake 等价与 Langlands 对偶群的钉住}

先需要张量函子的若干一般性质。设 $\mathcal C$ 是带纤维函子 $\omega:\mathcal C\to\Vec_F$（$F$ 是域）的刚性张量范畴，$H=\operatorname{Aut}^\otimes(\omega)$ 为由 $(\mathcal C,\omega)$ 决定的 $F$ 上代数群。

对任何 $F$-代数 $R$，令 $\omega_R:\mathcal C\xrightarrow{\omega}\Vec_F\xrightarrow{\otimes_FR}\mathrm{Mod}_R$。记 $\operatorname{Aut}^\otimes(\mathcal C,\omega_R)$ 为对 $(\sigma,\alpha)$ 的同构类构成的集合，其中 $\sigma:\mathcal C\xrightarrow{\sim}\mathcal C$ 是张量自等价，$\alpha:\omega_R\circ\sigma\Rightarrow\omega_R$ 是函子的自然同构。则 $\operatorname{Aut}^\otimes(\mathcal C,\omega_R)$ 带自然的群结构。记 $\operatorname{Aut}^\otimes(\mathcal C,\omega)$ 为函子 $R\mapsto\operatorname{Aut}^\otimes(\mathcal C,\omega_R)$，它定义 fppf 群层。另一方面，记 $\operatorname{Aut}(H)$ 为 $F$ 上 pro-代数群 $H$ 的自同构 fppf 层。

\begin{lemma}
存在 fppf 群层的自然同构 $\operatorname{Aut}^\otimes(\mathcal C,\omega)\xrightarrow{\sim}\operatorname{Aut}(H)$。特别地，有群的同构 $\operatorname{Aut}^\otimes(\mathcal C,\omega)\xrightarrow{\sim}\operatorname{Aut}(H)$，它诱导群的同构 $[\operatorname{Aut}^\otimes(\mathcal C)]\xrightarrow{\sim}\operatorname{Out}(H)$，其中 $[\operatorname{Aut}^\otimes(\mathcal C)]$ 是 $\mathcal C$ 的张量自等价的同构类集合。
\end{lemma}

在 $F$-点的层面，对 $(\sigma,\alpha)\in\operatorname{Aut}^\otimes(\mathcal C,\omega)$，它给出 $H=\operatorname{Aut}^\otimes(\omega)$ 的下述自同构：把 $h:\omega\Rightarrow\omega$ 送到自然变换
\[
\omega\xrightarrow{\alpha}\omega\circ\sigma\xrightarrow{h\circ\mathrm{id}_\sigma=\mathrm{id}}\omega\circ\sigma\xrightarrow{\alpha^{-1}}\omega.
\]

更一般地，设给定张量函子 $\Phi:\mathcal C\to\mathcal C'$，其中 $\mathcal C'$ 是另一刚性张量范畴，则可类似地定义群层 $\operatorname{Aut}^\otimes(\mathcal C,\Phi)$。

设 $\omega':\mathcal C'\to\Vec$ 是纤维函子，$\omega=\omega'\circ\Phi$。则把 $(\sigma,\alpha)\in\operatorname{Aut}^\otimes(\mathcal C,\Phi)$ 送到 $(\sigma,\mathrm{id}_{\omega'}\circ\alpha)\in\operatorname{Aut}^\otimes(\mathcal C,\omega)$，给出自然同态 $\omega'_*:\operatorname{Aut}^\otimes(\mathcal C,\Phi)\to\operatorname{Aut}^\otimes(\mathcal C,\omega)=\operatorname{Aut}(H)$。换言之，$\operatorname{Aut}^\otimes(\mathcal C,\Phi)$ 作用在 pro-代数群 $H$ 上。另一方面，有 pro-代数群的自然同态 $\Phi^*:H'=\operatorname{Aut}^\otimes(\omega')\to H=\operatorname{Aut}^\otimes(\omega)$。

\begin{lemma}
同态 $\Phi^*:H'\to H$ 经 $H'\to H^{\operatorname{Aut}^\otimes(\mathcal C,\Phi)}\subset H$ 分解。
\end{lemma}

现考虑 §2.3 中标准化的半单 Satake 范畴 $\mathcal S$。沿袭 [30] 与 [18]，用整体截面函子 $h$ 定义 $\ckG=\operatorname{Aut}^\otimes(h)$，并得到几何 Satake 等价 $\mathcal S\cong\Rep(\ckG)$。

\begin{lemma}
\begin{enumerate}[label=(\arabic*), leftmargin=1.9em, itemsep=.2em]
\item 存在自然同态 $\operatorname{Aut}(G)\to\operatorname{Aut}^\otimes(\mathcal S,h)\cong\operatorname{Aut}(\ckG)$，它经 $\tilde\iota:\operatorname{Out}(G)\to\operatorname{Aut}(\ckG)$ 分解；
\item 存在 $\ckG$ 的自然钉住 $\dagger=(\check B,\check T,\{x_{\alpha_i^\vee}\})$，它被经 $\tilde\iota$ 的 $\operatorname{Out}(G)$-作用保持。记 $\operatorname{Aut}^\dagger(\ckG)$ 为 $\ckG$ 的保持该钉住的自同构群。则 $\tilde\iota$ 诱导同构 $\iota:\operatorname{Out}(G)\xrightarrow{\sim}\operatorname{Aut}^\dagger(\ckG)$。
\end{enumerate}
\end{lemma}

\begin{proof}
(1) $\operatorname{Aut}(G)$ 在 $(\mathcal S,h)$ 上的作用由其在对 $\mathrm{Gr}_G$ 的作用诱导。由于 $\mathcal S$ 中的对象带 $G$ 在 $\mathrm{Gr}_G$ 上共轭作用的 $G$-等变结构，$G$ 的内自同构在 $(\mathcal S,h)$ 上平凡作用，即 $\operatorname{Aut}(G)$ 在 $(\mathcal S,h)$ 上的作用经 $\operatorname{Out}(G)$ 分解。

(2) 先需给出一个被 $\operatorname{Aut}(G)$-作用保持的 $\ckG$ 的钉住。为此需给出极大环面 $\check T\subset\ckG$、余特征 $2\rho\in X_*(T)=X_*(\check T)$（钉住中正余根和之半）与主幂幺元 $e\in\ckg$（单根向量之和）。等价地，需要
\begin{itemize}[leftmargin=1.6em, itemsep=.2em]
\item[(P1)] 把纤维函子 $h$ 分解为张量函子
\[
h=\bigoplus_{\mu\in X_*(T)}h_\mu:\mathcal S\to\Vec^{X_*(T)}\xrightarrow{\text{遗忘}}\Vec;
\]
\item[(P2)] 找张量导子 $e:h\to h$（即 $e(K_1*K_2)=e(K_1)\otimes\mathrm{id}_{K_2}+\mathrm{id}_{K_1}\otimes e(K_2)$），它把 $h_\mu$ 送到 $\oplus_ih_{\mu+\alpha_i}$，且对每个单根 $\alpha_i$，分量 $h_\mu\to h_{\mu+\alpha_i}$ 作为函子非零。
\end{itemize}

分解 $h_\mu$ 由 Mirković--Vilonen 的“权函子”给出，见 [30, Theorem 3.5, 3.6]。他们还证明 $h=\bigoplus_ih_i$（$h_i$ 是满足 $\langle2\rho,\mu\rangle=i$ 的 $h_\mu$ 之和）与 $h=H^*(\mathrm{Gr}_G,-)$ 的上同调次数分次一致。张量导子 $e$ 由与 $c_1(\mathcal L_{\det})\in H^2(\mathrm{Gr}_G,\mathbb{Q}_\ell)$ 的 cup 积给出，其中 $\mathcal L_{\det}$ 是 $\mathrm{Gr}_G$ 上的行列式线丛。

$\operatorname{Aut}^\dagger(G)\subset\operatorname{Aut}(G)$ 在 $\mathcal S$ 上的作用以与置换 $\mu\in X_*(T)$ 相同的方式置换权函子 $h_\mu$，保持 $\langle2\rho,\mu\rangle$（从而保持 $2\rho\in X_*(T)=X_*(\check T)$），并与 $e=c_1(\mathcal L_{\det})$ 交换。因此 $\operatorname{Aut}^\dagger(G)\xrightarrow{\sim}\operatorname{Out}(G)$ 在 $\mathcal S$ 上的作用保持上述钉住。

$\sigma\in\operatorname{Aut}^\dagger(G)$ 诱导 $\ckG$ 的 Dynkin 图的对偶自同构 $\check\sigma$。由于 $\sigma^*\mathrm{IC}_\mu\cong\mathrm{IC}_{\sigma^{-1}(\mu)}$，$\mathcal S\cong\Rep(\ckG)$ 的自等价 $\sigma^*$ 同构于 $\Rep(\ckG)$ 的由 $\ckG$ 的 Dynkin 图上对偶自同构 $\check\sigma^{-1}$ 给出的钉住自同构所诱导的自等价。这证明
\[
\operatorname{Out}(G)\to[\operatorname{Aut}^\otimes(\mathcal S)]\cong[\operatorname{Aut}^\otimes(\Rep(\ckG))]\cong\operatorname{Out}(\ckG)
\]
是同构（最后一个同构来自引理 9.1）。因此 $\iota:\operatorname{Out}(G)\to\operatorname{Aut}^\dagger(\ckG)$ 也是同构。
\end{proof}

为了 §6 的目的，还需要 (2.11) 中定义的另一纤维函子 $\omega_\varphi$。回忆 $S_\varphi$ 是 $\varphi$ 在 $T\rtimes\operatorname{Aut}^\dagger(G)\times\mathbb{G}_m^{\mathrm{rot}}$ 作用下的稳定子，$S^1_\varphi(k)=\ker(S_\varphi(k)\to\mathbb{G}_m^{\mathrm{rot}}(k))$。由 (6.4) 回忆 $S^1_\varphi(k)$ 作用在 ${}_\varphi\ckG$ 上。

\begin{lemma}
存在 ${}_\varphi\ckG$ 的自然钉住 $\ddagger=({}_\varphi\check B,{}_\varphi\check T,\{{}_\varphi x_{\alpha_i^\vee}\})$，它被 $S^1_\varphi(k)$-作用保持。
\end{lemma}

\begin{proof}
利用 (4.4)，可把纤维函子 $\omega_\varphi$ 改写为
\[
\mathcal S\xrightarrow{\Psi}\mathrm{Perv}_{I_0}(\mathrm{Fl}_G)\xrightarrow{V_\varphi}\Vec,
\]
其中 $\mathrm{Fl}_G=G((t))/I_0$ 是 $\{0\}$ 处的仿射旗簇，$\Psi:\mathcal S\to\mathrm{Perv}(I_0\backslash G((t))/I_0)=\mathrm{Perv}_{I_0}(\mathrm{Fl})$ 是 Gaitsgory 的邻近闭链函子 [16]，$V_\varphi(K):=R\Gamma_c(\mathrm{Fl}_G,K\otimes\mathrm{pr}_1^*A_\varphi)$ 如 (4.4)。为给出 ${}_\varphi\ckG$ 的钉住，需找到引理 9.3(2) 证明中 (P1) 与 (P2) 的类比。

据 Arkhipov--Bezrukavnikov [1, Theorem 4]，每个对象 $\Psi(K)$ 带 $X_*(T)$-滤链，其相伴分次片为 Wakimoto 层。更精确地，他们构造了函子
\begin{equation}
\bigoplus_{\mu\in X_*(T)}W^\mu\circ\Psi(-):\mathcal S\to\Vec^{X_*(T)}.
\end{equation}
这里对带 Wakimoto 滤链 $\Psi(K)_{\le\mu}$ 的 $\Psi(K)\in\mathrm{Perv}_{I_0}(\mathrm{Fl}_G)$，记 $\mathrm{gr}_\mu\Psi(K)=\Psi(K)_{\le\mu}/\Psi(K)_{<\mu}=J_\mu\otimes W^\mu(K)$，其中 $J_\mu$ 是 Wakimoto 层 [1, §3.2]，$W^\mu(K)$ 是带 Frobenius 作用的向量空间。在 [1, Theorem 6] 中证明了 (9.5) 是张量的。

对每个 $\mu$，设 $j_\mu:\mathrm{Fl}_\mu=I_0t^\mu I_0/I_0\hookrightarrow\mathrm{Fl}_G$ 为包含。按如下方式固定 $J_\mu$ 的 Frobenius 结构：对正则支配的 $\mu$，令 $J_\mu=j_{\mu,*}\mathbb{Q}_\ell[\langle2\rho,\mu\rangle](\langle2\rho,\mu\rangle)$；对正则反支配的 $\mu$，令 $J_\mu=j_{\mu,!}\mathbb{Q}_\ell[\langle2\rho,\mu\rangle]$；对一般的 $\mu=\mu_1+\mu_2$（$\mu_1$ 正则支配、$\mu_2$ 正则反支配），令 $J_\mu=J_{\mu_1}*^{I_0}J_{\mu_2}$（其中 $*^{I_0}$ 是 $\mathrm{Fl}_G$ 上的卷积）。由 [1, Lemma 8, Corollary 1] 知 $J_\mu$ 良定义。此归一化保证在 $J_\mu$ 的合成列中 $\delta=\mathrm{IC}_1$ 恰出现一次，重数空间为 $\mathbb{Q}_\ell$ 且作为平凡 Frobenius 模（见 [1, Lemma 3(a)]，作混合情形的显然调整）。由于除 $e_w=1$ 外所有 $\mathrm{IC}_{e_w}$ 都被 $V_\varphi$ 消去，得
\begin{equation}
V_\varphi(J_\mu)=\mathbb{Q}_\ell\ \text{作为平凡 Frobenius 模，对一切}\ \mu\in X_*(T).
\end{equation}

\medskip\noindent\textbf{断言.} 对 $K\in\mathcal S$，$W^\mu(K)$ 是权 $\langle2\rho,\mu\rangle$ 的纯模。事实上，有函子的自然同构 $h_\mu\cong W^\mu$（$h_\mu$ 是 [30, Theorem 3.5, 3.6] 中的权函子，已在引理 9.3 的证明中用到）。

\begin{proof}[断言的证明]
先回忆上引文献中权函子的定义。对每个 $\mu\in X_*(T)$，设 $S_\mu\subset\mathrm{Gr}_G$ 是含 $t^\mu$ 的 $U((t))$-轨道。上引文献中定义的权函子为 $h_\mu(K)=H^*_c(S_\mu,K)$，它集中在次数 $\langle2\rho,\mu\rangle$。

设 $\pi:\mathrm{Fl}_G\to\mathrm{Gr}_G$ 为投影，则 $\pi^{-1}(S_\mu)=\bigsqcup_{w\in W}\widetilde S_{\mu w}$，其中 $\widetilde S_{\mu w}\subset\mathrm{Fl}_G$ 是含 $t^\mu w$ 的 $U((t))$-轨道。对 $K\in\mathcal S$ 有自然同构
\[
h_\mu(K)=H^*_c(S_\mu,\pi_!\Psi(K))\cong H^*_c(\pi^{-1}(S_\mu),\Psi(K))\xrightarrow{\sim}H^*_c(\widetilde S_\mu,\Psi(K))\xrightarrow{\sim}W^\mu(K).
\]
这里第一个等式来自 $\pi_!\Psi(K)=K$，最后两个同构来自 [1, Theorem 4(2)]（须对 Frobenius 结构多加小心）。

由于 $h_i(K)=H^i(\mathrm{Gr}_G,K)$ 是权 $i$ 的纯模，而 $h_\mu(K)$ 是 $h_{\langle2\rho,\mu\rangle}$ 的直和项，故 $h_\mu(K)$ 是权 $\langle2\rho,\mu\rangle$ 的纯模。因此 $W^\mu(K)$ 也是权 $\langle2\rho,\mu\rangle$ 的纯模。
\end{proof}

下面构造自然同构 $\bigoplus_\mu W^\mu\cong\omega_\varphi$。对每个 $K\in\mathcal S$，$\Psi(K)$ 上的 Wakimoto 滤链给出计算 $V_\varphi(\Psi(K))$ 的谱序列，其 $E_1$ 页为 $V_\varphi(\mathrm{gr}_\mu\Psi(K))$。由 (9.6)，
\begin{equation}
V_\varphi(\mathrm{gr}_\mu\Psi(K))=V_\varphi(J_\mu)\otimes W^\mu(K)=W^\mu(K)
\end{equation}
集中在次数 $0$，故谱序列在 $E_1$ 处退化。谱序列的极限给出 Frobenius 模 $\omega_\varphi(K)=V_\varphi(\Psi(K))$ 上的 Wakimoto 滤链，记作 $w_{\le\mu}$。由断言与 (9.7)，该滤链细化权滤链 $w_{\le i}$：
\[
\mathrm{gr}^w_i\omega_\varphi(K)=\bigoplus_{\langle2\rho,\mu\rangle=i}\mathrm{gr}^w_\mu\omega_\varphi(K).
\]
由于 $\omega_\varphi(K)$ 是 $\mathbb{Q}_\ell$-向量空间中的 Frobenius 模，权滤链典范分裂。因此 Wakimoto 滤链 $w_{\le\mu}\omega_\varphi(K)$ 也典范分裂。这给出典范同构 $\bigoplus_\mu W^\mu\cong\omega_\varphi$。

最后，主幂幺元由单值作用（$0$ 处驯惯性群的拓扑生成元）在邻近闭链上的对数给出（[16, Theorem 2]）：$M_K:\Psi(K)\to\Psi(K)(-1)$。

$S^1_\varphi(k)$ 的作用与邻近闭链函子交换，故与单值 $M_K$ 交换。它置换 Wakimoto 层，从而经 $S^1_\varphi(k)\to\operatorname{Aut}^\dagger(G)$ 在 $X_*(T)$ 上的作用置换函子 $W^\mu$，并保持 $\langle2\rho,\mu\rangle$。因此 $S^1_\varphi(k)$-作用保持钉住 $\ddagger$。
\end{proof}

% ============================================================
%  ch14 —— §10、§11（原文 p.52–p.55）
% ============================================================
\routenote{p.52--p.55}

\section{附录：拟极小的组合论}

本节假设 $G$ 几乎是单的且秩至少为 $2$。

设 $\theta$ 是最高根（它是长根），$\theta^\vee$ 是对应的余根（它是短余根）。根集 $\Phi$ 被划分为
\[
\Phi^\theta_n=\{\alpha\in\Phi\mid\langle\alpha,\theta^\vee\rangle=n\},\qquad n=0,\pm1,\pm2,
\]
其中 $\Phi^\theta_{\pm2}=\{\pm\theta\}$。$\Phi^\theta_1$ 中的根成对配为 $(\alpha,\beta)$，满足 $\alpha+\beta=\theta$。

设 $V_{\theta^\vee}$ 是 $\ckG$ 的以 $\theta^\vee$ 为最高权的不可约表示。$V_{\theta^\vee}$ 的非零权正是 $\ckG$ 的短根，每个重数为 $1$。设 $\{e=\sum_ix_i,\ 2\rho,\ f=\sum_iy_i\}\in\ckg$ 是主 $\mathfrak{sl}_2$-三元组，其中 $x_i\in\ckg_{\alpha_i^\vee}$、$y_i\in\ckg_{-\alpha_i^\vee}$ 非零，$2\rho\in\ck t$ 是 $\ckG$ 的正余根之和。

\begin{lemma}
下列数相等：
\begin{enumerate}[label=(\arabic*), leftmargin=1.9em, itemsep=.2em]
\item $\dim V_{\theta^\vee}(0)$，其中 $V_{\theta^\vee}(0)$ 是 $2\rho$-作用下的零权空间，也是 $\check T$-作用下的零权空间；
\item $\dim V^e_{\theta^\vee}$，其中 $V^e_{\theta^\vee}=\ker(\operatorname{ad}(e)|_{V_{\theta^\vee}})$；
\item $\ckG$ 的短单根数（短秩 $r_s(\ckG)$）；
\item $\#(W\theta^\vee)/h=\#\{\ckG\ \text{的短根}\}/h$（这里 $h$ 是 $W$ 的 Coxeter 数）。
\end{enumerate}
\end{lemma}

\begin{proof}
在主 $\mathfrak{sl}_2$-作用下，$V_{\theta^\vee}$ 可分解为 $\mathfrak{sl}_2$ 的不可约表示之和：
\[
V_{\theta^\vee}=\bigoplus_{i=1}^{r_s(\ckG)}\operatorname{Sym}^{2\ell_i}(\mathrm{Std}),
\]
其中 $\mathrm{Std}$ 是 $\mathfrak{sl}_2$ 的二维表示。由于 $2\rho$-作用在 $V_{\theta^\vee}$ 上的权都是偶数，$V_{\theta^\vee}$ 中只出现 $\mathrm{Std}$ 的偶对称幂。每个 $\operatorname{Sym}^{2\ell_i}(\mathrm{Std})$ 对 $V_{\theta^\vee}(0)$ 与 $V^e_{\theta^\vee}$ 各贡献 $1$ 维，故
\[
r_s(\ckG)=\dim V_{\theta^\vee}(0)=\dim V^e_{\theta^\vee}.
\]
这证明了 (1) 与 (2) 中的数相等。

设 $V_{\theta^\vee}(n)$ 是 $\rho$-作用的权 $n$-特征空间，则
\[
V_{\theta^\vee}(n)=\bigoplus_{\alpha^\vee\ \text{短},\ \langle\rho,\alpha^\vee\rangle=n}V_{\theta^\vee}(\alpha^\vee).
\]
特别地，$\dim V_{\theta^\vee}(1)$ 是 $\ckG$ 的短单根数。映射 $e:V_{\theta^\vee}(0)\to V_{\theta^\vee}(1)$ 显然满。它也是单的：若 $v\in V_{\theta^\vee}(0)$ 且 $ev=0$，则对一切单根 $\alpha_i^\vee$ 有 $\ckg_{\alpha_i^\vee}v=0$，即 $v$ 是最高权向量，矛盾。这证明了
\begin{equation}
\dim V_{\theta^\vee}(0)=\dim V_{\theta^\vee}(1)=r_s(\ckG).
\end{equation}

还需证明 (4) 与其余各数相同。论证与 [28, §6.7] 类似，那里 Kostant 考虑的是伴随表示而非 $V_{\theta^\vee}$。这里只给梗概。设 $\zeta$ 是本原 $h$ 次单位根，$P=\rho(\zeta)\in\ckG$。设 $\gamma^\vee$ 是 $\ckG$ 的最高根，$z=e+x_{-\gamma^\vee}$。则 $z$ 是 $\ckg$ 中的正则半单元，且 $\operatorname{Ad}(P)z=\zeta z$。设 $\ckG_z$ 是 $z$ 的中心化子（它是极大环面），则 $P\in N_{\ckG}(\ckG_z)$，其在 Weyl 群中的像是 Coxeter 元。取 $\operatorname{Sym}^{2\ell_i}(\mathrm{Std})\subset V_{\theta^\vee}$ 中最高权线的一组基 $u_i$，则存在唯一的 $v_i\in V_{\theta^\vee}(\ell_i-h)$ 使 $u_i+v_i\in V^z_{\theta^\vee}$（即 $z$-作用在 $V_{\theta^\vee}$ 上的核）。于是 $\{u_i+v_i\}_{1\le i\le r_s(\ckG)}$ 构成 $V^z_{\theta^\vee}$ 的基，在 $\operatorname{Ad}(P)$ 作用下的特征值为 $\zeta^{\ell_i}$。

$\ckG$ 的表示 $V_{\theta^\vee}$ 显然自对偶。按 $\ckg^z$ 的零权与非零权，可把 $V_{\theta^\vee}$ 写成 $N_{\ckG}(\ckG_z)$-模的直和 $V_{\theta^\vee}=V^z_{\theta^\vee}\oplus V'$。$\ckG$-模的任何自对偶 $V_{\theta^\vee}\cong V^*_{\theta^\vee}$ 必然限制为 $N_{\ckG}(\ckG_z)$-模的自对偶 $V^z_{\theta^\vee}\cong(V^z_{\theta^\vee})^*$。因此 $\operatorname{Ad}(P)$ 在 $V^z_{\theta^\vee}$ 上的特征值在求逆下不变：即多重集 $\{\zeta^{\ell_i}\}_{1\le i\le r_s(\ckG)}$ 在求逆下不变。这蕴含 $\sum_i\ell_i=r_s(\ckG)h/2$。故 $\dim V_{\theta^\vee}=\sum_i(2\ell_i+1)=(h+1)r_s(\ckG)$。由于 $V_{\theta^\vee}$ 的非零权空间以 $\ckG$ 的短根为指标，从而以轨道 $W\theta^\vee$ 为指标，得 $\#(W\theta^\vee)=\dim V_{\theta^\vee}-\dim V_{\theta^\vee}(0)=hr_s(\ckG)$。这证明 (4) 与其余各数一致。
\end{proof}

\section{附录：$G_2$ 的伴随 Schubert 簇}

本节假设 $\operatorname{char}(k)>3$。回忆（见 [8] 中的表）$G_2$ 有四个非正则幂幺轨道：
\begin{itemize}[leftmargin=1.6em, itemsep=.2em]
\item 次正则轨道，含 $P_{-\theta}$ 的幂幺根中的一般元；其维数为 $10$；
\item 含 $U^\times_{-\gamma}$ 的轨道，维数为 $8$；
\item 含 $U^\times_{-\theta}$ 的轨道，维数为 $6$；
\item 恒等轨道，维数为 $0$。
\end{itemize}
对 $G_2$，$\mathrm{Gr}^{\mathrm{triv}}_{\le\gamma^\vee}$ 的 $G$-轨道与这些幂幺轨道密切相关。更精确地：

\begin{lemma}
\begin{enumerate}[label=(\arabic*), leftmargin=1.9em, itemsep=.25em]
\item $\mathrm{Gr}^{\mathrm{triv}}_{\le\gamma^\vee}$ 上有四个 $\operatorname{Ad}(G)$-轨道
\[
\mathrm{Gr}^{\mathrm{triv}}_{\le\gamma^\vee}=\mathrm{Gr}^{\mathrm{subr}}_G\sqcup\operatorname{Ad}(G)U^\times_{-\gamma,-1}\sqcup\mathrm{Gr}^{\mathrm{triv}}_{\theta^\vee}\sqcup\{\star\},
\]
其维数分别为 $10,8,6,0$；在赋值映射 $\mathrm{ev}_{\tau=1}$ 下它们映到上述四个非正则幂幺轨道。
\item 态射
\[
\begin{aligned}
\nu:\ \frac{G}{P_{-\theta}}\times\Bigl(\prod_{\langle\beta,-\theta^\vee\rangle\ge1}U_{\beta,-1}\Bigr)&\longrightarrow\mathrm{Gr}^{\mathrm{triv}}_{\le\gamma^\vee},\\
\Bigl(g,\prod u_\beta(c_\beta\tau^{-1})\Bigr)&\mapsto\operatorname{Ad}(g)\Bigl(\prod u_\beta(c_\beta\tau^{-1})\Bigr)
\end{aligned}
\]
是消解。它在各 $G$-轨道上的纤维为：
\begin{itemize}[leftmargin=1.4em, itemsep=.1em]
\item $\nu$ 在 $\mathrm{Gr}^{\mathrm{subr}}$ 上是同构；
\item $\nu^{-1}(u_{-\gamma}(\tau^{-1}))\cong\PP^1$；
\item $\nu^{-1}(u_{-\theta}(\tau^{-1}))$ 是 $\PP^1$ 上的射影锥（它含一个点，其补是 $\PP^1$ 上的线丛）；
\item $\nu^{-1}(\star)\cong G/P_{-\theta}$。
\end{itemize}
\item 按与 $\nu$ 相同方式定义的态射
\[
\nu':\ \frac{G}{P_{-\gamma}}\times\Bigl(\prod_{\langle\beta,-\gamma^\vee\rangle\ge2}U_{\beta,-1}\Bigr)\longrightarrow\mathrm{Gr}^{\mathrm{triv}}_{\le\gamma^\vee}
\]
是 $\operatorname{Ad}(G)U^\times_{-\gamma,-1}$ 的闭包的消解。它在各 $G$-轨道上的纤维为：
\begin{itemize}[leftmargin=1.4em, itemsep=.1em]
\item $\nu'$ 在 $\operatorname{Ad}(G)U^\times_{-\gamma,-1}$ 上是同构；
\item $\nu'^{-1}(u_{-\theta}(\tau^{-1}))\cong\PP^1$；
\item $\nu'^{-1}(\star)\cong G/P_{-\gamma}$。
\end{itemize}
\end{enumerate}
\end{lemma}

\noindent\textbf{对 $G=G_2$ 证明定理 4(2)。} 经与其他型情形相同的化归步骤，只需证
\begin{equation}
\chi_c\bigl(\mathrm{Gr}^{\mathrm{triv},a=0}_{\gamma^\vee},\mathrm{ev}^*J\bigr)=-r_s(G)=-1.
\end{equation}

\medskip\noindent\textbf{第一步.} $\chi_c\bigl((\operatorname{Ad}(G)U^\times_{-\gamma,-1})^{a=0},\mathrm{ev}^*J\bigr)=-1$。

对任何短根 $\beta$，令 $V_\beta=\prod_{\langle\beta',\beta^\vee\rangle\ge2}U_{\beta',-1}$。则 $\nu'$ 的源有 Bruhat 分解
\[
\frac{G}{P_{-\gamma}}\times V_{-\gamma}=\bigsqcup_{\beta\ \text{短根}}\operatorname{Ad}(U)V_\beta.
\]
下述断言可仿 §5.6 中的断言证明。

\medskip\noindent\textbf{断言.} 对短根 $\beta$，
\[
\chi_c\bigl(\operatorname{Ad}(U)V^{>-\theta}_\beta,\mathrm{ev}^*J\bigr)=
\begin{cases}
0, & \Phi^\beta_{\ge2}\ \text{含单根},\\
1, & \text{否则}.
\end{cases}
\]

观察根系 $G_2$，有 $3$ 个短根 $\beta$ 使 $\Phi^\beta_{\ge2}$ 不含单根。因此 $\chi_c\bigl((\frac{G}{P_{-\gamma}}\times V_{-\gamma})^{a=0},\nu'^*\mathrm{ev}^*J\bigr)=3$。另一方面，由引理 11.1(3) 有
\[
\begin{aligned}
3&=\chi_c\Bigl(\bigl(\tfrac{G}{P_{-\gamma}}\times V_{-\gamma}\bigr)^{a=0},\nu'^*\mathrm{ev}^*J\Bigr)\\
&=\chi_c(G/P_{-\gamma})+\chi_c(\PP^1)\chi_c\bigl(\mathrm{Gr}^{\mathrm{triv},a=0}_{\theta^\vee},\mathrm{ev}^*J\bigr)\\
&\quad+\chi_c\bigl((\operatorname{Ad}(G)U^\times_{-\gamma,-1})^{a=0},\mathrm{ev}^*J\bigr).
\end{aligned}
\]
代入 $\chi_c(G/P_{-\gamma})=\#W/W_{\gamma^\vee}=6$、$\chi_c(\PP^1)=2$ 以及由定理 4(1) 的证明得到的 $\chi_c(\mathrm{Gr}^{\mathrm{triv},a=0}_{\theta^\vee},\mathrm{ev}^*J)=-1$，即得 $\chi_c\bigl((\operatorname{Ad}(G)U^\times_{-\gamma,-1})^{a=0},\mathrm{ev}^*J\bigr)=-1$。

\medskip\noindent\textbf{第二步.} $\chi_c\bigl(\mathrm{Gr}^{a=0}_{\mathrm{subr}},\mathrm{ev}^*J\bigr)=0$。

对任何长根 $\alpha$，令 $V_\alpha=\prod_{\langle\alpha',\alpha^\vee\rangle\ge1}U_{\alpha',-1}$。则 $\nu$ 的源有 Bruhat 分解
\[
\frac{G}{P_{-\theta}}\times V_{-\theta}=\bigsqcup_{\alpha\ \text{长根}}\operatorname{Ad}(U)V_\alpha.
\]
下述断言可仿 §5.6 中的断言证明。

\medskip\noindent\textbf{断言.} 对长根 $\alpha$，
\[
\chi_c\bigl(\operatorname{Ad}(U)V^{>-\theta}_\alpha,\mathrm{ev}^*J\bigr)=
\begin{cases}
0, & \Phi^\alpha_{\ge1}\ \text{含单根},\\
1, & \text{否则}.
\end{cases}
\]

观察根系 $G_2$，$\alpha=-\theta$ 是唯一使 $\Phi^\alpha_{\ge1}$ 不含单根的长根。因此 $\chi_c\bigl((\frac{G}{P_{-\theta}}\times V_{-\theta})^{a=0},\nu^*\mathrm{ev}^*J\bigr)=1$。另一方面，由引理 11.1(2) 有
\begin{align*}
1&=\chi_c\Bigl(\bigl(\tfrac{G}{P_{-\theta}}\times V_{-\theta}\bigr)^{a=0},\nu^*\mathrm{ev}^*J\Bigr)\\
&=\chi_c(G/P_{-\theta})+\chi_c(\PP^1)\chi_c\bigl((\operatorname{Ad}(G)U^\times_{-\gamma,-1})^{a=0},\mathrm{ev}^*J\bigr)\\
&\qquad+\chi_c\bigl(\nu^{-1}(u_{-\theta}(\tau^{-1}))\bigr)\chi_c\bigl(\mathrm{Gr}^{\mathrm{triv},a=0}_{\theta^\vee},\mathrm{ev}^*J\bigr)
+\chi_c\bigl(\mathrm{Gr}^{a=0}_{\mathrm{subr}},\mathrm{ev}^*J\bigr).
\end{align*}
代入 $\chi_c(G/P_{-\theta})=\#W/W_\theta=6$、$\chi_c(\PP^1)=2$、由引理 11.1(2) 得到的 $\chi_c(\nu^{-1}(u_{-\theta}(\tau^{-1})))=3$、第一步的 $\chi_c\bigl((\operatorname{Ad}(G)U^\times_{-\gamma,-1})^{a=0},\mathrm{ev}^*J\bigr)=-1$ 以及定理 4(1) 的证明中的 $\chi_c(\mathrm{Gr}^{\mathrm{triv},a=0}_{\theta^\vee},\mathrm{ev}^*J)=-1$，得 $\chi_c\bigl(\mathrm{Gr}^{a=0}_{\mathrm{subr}},\mathrm{ev}^*J\bigr)=0$。\zhnote{原文此处写作 ``we conclude that $\chi_c(\mathrm{Gr}^{a=0}_{\mathrm{subr}},\mathrm{ev}^*J)=-1$''，与其小标题（$=0$）及上式代入的结果（$6-2-3+x=1$ 给出 $x=0$）均不符；且只有取 $0$ 才能在末尾与第一步的 $-1$ 相加得到 (11.2)。故按 $0$ 译。[来源：原文 p.55 第 2 步末尾与 (11.2)]}

合并第一步与第二步，由 $\mathrm{Gr}^{\mathrm{triv},a=0}_{\gamma^\vee}=\mathrm{Gr}^{a=0}_{\mathrm{subr}}\sqcup(\operatorname{Ad}(G)U^\times_{-\gamma,-1})^{a=0}$ 即得 (11.2)。

这证明了 $G_2$ 情形的定理 4(2)。

% ============================================================
%  ch15 —— 参考文献（原文 p.55–p.56）
% ============================================================
\routenote{p.55--p.56}

\section*{参考文献}
\addcontentsline{toc}{section}{参考文献}

\begin{thebibliography}{99}\small

\bibitem{arkhipov} S. Arkhipov, R. Bezrukavnikov. \emph{Perverse sheaves on affine flags and Langlands dual group}（附 R. Bezrukavnikov 与 I. Mirković 的附录）. Israel J. Math. 170 (2009), 135--183.

\bibitem{bb} A. Beilinson, J. Bernstein. \emph{A proof of Jantzen conjectures}. I. M. Gelfand Seminar, 1--50, Adv. Soviet Math., 16, Part 1, Amer. Math. Soc., Providence, RI, 1993.

\bibitem{bbd} A. Beilinson, J. Bernstein, and P. Deligne. \emph{Faisceaux pervers}. Astérisque, 100, 1983.

\bibitem{bbm} R. Bezrukavnikov, A. Braverman, I. Mirković. \emph{Some results about the geometric Whittaker model}. Adv. Math. 186, (2004), no. 1, 143--152.

\bibitem{bez} R. Bezrukavnikov. \emph{On tensor categories attached to cells in affine Weyl groups}. Representation theory of algebraic groups and quantum groups, Adv. Stud. Pure Math., 40, Math. Soc. Japan, Tokyo, 2004, 69--90.

\bibitem{bourbaki} N. Bourbaki. \emph{Éléments de mathématique. Fasc. X\mbox{}XXIV. Groupes et algèbres de Lie. Chapitre IV--VI}. Actualités Scientifiques et Industrielles, No. 1337, Hermann, Paris 1968, 288 pp.

\bibitem{bt} F. Bruhat, J. Tits. \emph{Groupes réductifs sur un corps local. II. Schémas en groupes. Existence d'une donnée radicielle valuée}. Publ. Math. IHES 60 (1984), 197--376.

\bibitem{carter} R.W. Carter, \emph{Finite groups of Lie type. Conjugacy classes and complex characters}. Pure and Applied Mathematics. A Wiley-Interscience Publication. John Wiley \& Sons, Inc., New York, 1985. xii+544 pp.

\bibitem{chev} C. Chevalley. \emph{Sur certains groupes simples}. Tôhoku Math. J. (2), 7, (1955), 14--66.

\bibitem{sga4} P. Deligne. \emph{Cohomologie étale}. Séminaire de Géométrie Algébrique du Bois-Marie SGA $4\frac{1}{2}$. Avec la collaboration de J. F. Boutot, A. Grothendieck, L. Illusie et J. L. Verdier. Lecture Notes in Mathematics, Vol. 569. Springer-Verlag, Berlin-New York, 1977. iv+312pp.

\bibitem{weil2} P. Deligne. \emph{La conjecture de Weil II}. Publ. Math. IHES No. 52 (1980), 137--252.

\bibitem{dm} P. Deligne, J.S. Milne. \emph{Tannakian Categories}. In Hodge cycles, motives, and Shimura varieties., Lecture Notes in Mathematics, 900. Springer-Verlag, Berlin-New York, 1982.

\bibitem{sga3} M. Demazure, A. Grothendieck. \emph{Schémas en groupes. III: Structure des schémas en groupes réductifs}. Séminaire de Géométrie Algébrique du Bois Marie 1962/64 (SGA 3). Lecture Notes in Mathematics, Vol. 153 Springer-Verlag, Berlin-New York 1962/1964 viii+529 pp.

\bibitem{faltings} G. Faltings. \emph{Algebraic loop groups and moduli spaces of bundles}. J. Eur. Math. Soc. 5, (2003), 41--68.

\bibitem{fg} E. Frenkel, B. Gross. \emph{A rigid irregular connection on the projective line} Ann. of Math. (2) 170, (2009), no. 3, 1469--1512.

\bibitem{gaits} D. Gaitsgory. \emph{Construction of central elements in the affine Hecke algebra via nearby cycles}. Invent. Math. 144 (2001), no. 2, 253--280.

\bibitem{gaits2} D. Gaitsgory \emph{On de Jong's conjecture}. Israel J. Math. 157 (2007), 155--191.

\bibitem{ginzburg} V. Ginzburg \emph{Perverse sheaves on a Loop group and Langlands' duality}. Preprint, arXiv: alg-geom/9511007.

\bibitem{gh} U. Görtz, T. Haines. \emph{The Jordan-Hölder series for nearby cycles on some Shimura varieties and affine flag varieties}. J. Reine Angew. Math. 609 (2007), 161--213.

\bibitem{gr} B. Gross, M. Reeder. \emph{Arithmetic invariants of discrete Langlands parameters}. Preprint, \url{http://www.math.harvard.edu/~gross/preprints/}

\bibitem{gross} B. Gross. \emph{Irreducible cuspidal representations with prescribed local behavior}. Preprint, \url{http://www.math.harvard.edu/~gross/preprints/}

\bibitem{letter} B. Gross. \emph{Letter to Deligne and Katz}. July 14, 2008.

\bibitem{hr} T. Haines, M. Rapoport. \emph{Appendix: On parahoric subgroups}. Adv. in Math. 219, no. 1, (2008), 188--198; appendix to [32].

\bibitem{heinloth} J. Heinloth. \emph{Uniformization for $G$-bundles}. Math. Ann. (2010)...

\bibitem{katz88} N. Katz. \emph{Gauss sums, Kloosterman sums, and monodromy groups}. Annals of Mathematics Studies, 116. Princeton University Press, Princeton, NJ, 1988. x+246 pp.

\bibitem{katz90} N. Katz. \emph{Exponential sums and differential equations}. Annals of Mathematics Studies, 124. Princeton University Press, Princeton, NJ, 1990. xii+430 pp.

\bibitem{katz96} N. Katz. \emph{Rigid local systems}. Annals of Mathematics Studies, 139. Princeton University Press, Princeton, NJ, 1996. viii+223 pp.

\bibitem{kostant} B. Kostant. \emph{The principal three-dimensional subgroup and the Betti numbers of a complex simple Lie group}. Amer. J. Math., 81 (1959), 973--1032.

\bibitem{lusztig} G. Lusztig. \emph{Singularities, character formulas, and a q-analog of weight multiplicities}. Analysis and topology on singular spaces, II, III (Luminy, 1981), 208--229. Astérisque, 101-102, Soc. Math. France, Paris, 1983.

\bibitem{mv} I. Mirković, K. Vilonen. \emph{Geometric Langlands duality and representations of algebraic groups over commutative rings}. Ann. of Math. (2) 166 (2007), no. 1, 95--143.

\bibitem{npo} B.C. Ngô and P. Polo. \emph{Résolutions de Demazure affines et formule de Casselman-Shalika géométrique}. J. Alg. Geom. 10 (2001), 515--547.

\bibitem{pr} G. Pappas, M. Rapoport. \emph{Twisted loop groups and their affine flag varieties}. Adv. Math. 219, (2008), no. 1, 118--198.

\bibitem{ramanathan} A. Ramanathan. \emph{Deformations of principal bundles on the projective line}. Invent. Math. 71, (1983), no. 1, 165--191.

\bibitem{reeder} M. Reeder. \emph{Torsion automorphisms of simple Lie algebras}. To appear in L'Enseignement Mathématique, available at \url{http://www2.bc.edu/~reederma/Torsion.pdf}

\bibitem{springer} T.A. Springer. \emph{Regular elements of finite reflection groups}. Invent. Math. 25, (1974), 159--198.

\bibitem{yu} J-K. Yu. \emph{Smooth models associated to concave functions in Bruhat-Tits theory}. Preprint, \url{http://www.math.purdue.edu/~jyu/prep/model.pdf}

\end{thebibliography}

\vspace{1em}
\noindent{\small\itshape 参考文献保留原文语种与著录格式，仅将连接符、破折号与引号按中文排版规范统一；方括号中的编号与原文一致。}


\end{document}
